% Quelques machines simples — PCT 3ème — Cours signé OnBuch+
% Programme officiel MINESEC. Compiler avec : tectonic N11-machines-simples.tex
\newcommand{\DOCMATIERE}{PCT}
\newcommand{\DOCNIVEAU}{3ème}
\newcommand{\DOCMODULE}{Module 2 : Électricité}
\newcommand{\DOCLECON}{N11}
\newcommand{\DOCTITRE}{Quelques machines simples}
% OnBuch+ — préambule « cours signés » ENRICHI (compilé avec tectonic / XeLaTeX)
% Système visuel « Pop » d'OnBuch+ : fond crème (jamais blanc), bordures encre
% épaisses, ombres « dures » décalées, titres Archivo Black en capitales.
% Version MATHÉMATIQUES : graphes (pgfplots/tikz), boîtes pédagogiques
% (définition, propriété, méthode, attention, le savais-tu, exercice résolu,
% exercice, TP, bilan), page de garde et signature OnBuch+ ancrée en bas.
%
% Macros attendues dans main.tex AVANT \input{preamble} :
%   \DOCMATIERE  \DOCNIVEAU  \DOCMODULE  \DOCLECON (numéro)  \DOCTITRE
\documentclass[11pt]{article}

\usepackage{fontspec}
\usepackage[french]{babel}
\usepackage[a4paper,margin=2cm,top=3.4cm,bottom=2.8cm,headheight=1.4cm,footskip=1.3cm]{geometry}
\usepackage{amsmath,amssymb,amsfonts}
\usepackage{mathtools} % \xmapsto, \coloneqq... souvent utilisés par les modèles
\usepackage{cancel} % simplifications barrées (bilans de forces)
% \abs{z} (module, valeur absolue) et \norm{u} (norme), utilisés par les modèles
\providecommand{\abs}{}\let\abs\relax\DeclarePairedDelimiter{\abs}{\lvert}{\rvert}
\providecommand{\norm}{}\let\norm\relax\DeclarePairedDelimiter{\norm}{\lVert}{\rVert}
\usepackage[table]{xcolor} % [table] : fournit aussi \rowcolors (alternance de lignes)
\usepackage{pagecolor}
\usepackage{tikz}
\usetikzlibrary{arrows.meta,positioning,calc,shapes.geometric,shapes.misc,shapes.arrows,decorations.pathmorphing,decorations.pathreplacing,decorations.markings,patterns,shadows,mindmap,trees,fit,backgrounds,matrix,intersections,angles,quotes}
\usepackage{pgfplots}
\usepackage[version=4]{mhchem} % équations nucléaires \ce{^{235}_{92}U}
\usepackage[european,siunitx]{circuitikz} % schémas électriques
\pgfplotsset{compat=1.18}
\usepgfplotslibrary{fillbetween,groupplots} % aires entre courbes (calcul intégral)
\usepackage{enumitem}
\usepackage{fancyhdr}
% [most] charge toutes les bibliothèques tcolorbox (listings, minted, raster,
% documentation...) et gonfle inutilement la mémoire TeX sur un document long
% (~300 boîtes) : on ne charge que ce qui est réellement utilisé.
\usepackage[skins,breakable]{tcolorbox}
\usepackage{graphicx}
\usepackage{colortbl}
\usepackage{booktabs}
\usepackage{tabularx}
\usepackage{multirow}
\usepackage{diagbox} % case d'en-tête diagonale des tableaux à double entrée
% Colonnes extensibles centrée / gauche / droite (C, L, R), souvent utilisées par les modèles
\newcolumntype{C}{>{\centering\arraybackslash}X}
\newcolumntype{L}{>{\raggedright\arraybackslash}X}
\newcolumntype{R}{>{\raggedleft\arraybackslash}X}
\usepackage{array}
\usepackage{multicol}
\usepackage{siunitx}
% parse-numbers=false : accepte aussi \SI{2,5 \times 10^{-3}}{...}
\sisetup{locale=FR,output-decimal-marker={,},group-minimum-digits=4,parse-numbers=false}
\DeclareSIUnit\ohm{\ensuremath{\symup{\Omega}}} % U+2126 absent de la police
\DeclareSIUnit\division{div} % réglages d'oscilloscope (ms/div, V/div)
\DeclareSIUnit\tr{tr} % tours (vitesse de rotation, ex. \SI{1500}{\tr\per\minute})
\DeclareSIUnit\molar{M} % concentration molaire (ex. \SI{3}{\molar}, \SI{1}{\micro\molar})
\DeclareSIUnit\an{an} % année (le modèle confond parfois avec \year, un compteur TeX) — ex. \SI{5}{\kilo\meter\per\an}
\DeclareSIUnit\FCFA{FCFA} % franc CFA (ex. \SI{500}{\FCFA\per\kilo\gram})
\DeclareSIUnit\degreeBrix{\textdegree Brix} % teneur en sucre d'un sirop/jus (réfractométrie)
\DeclareSIUnit\month{mois} % le modèle utilise parfois \month, un compteur TeX, comme unité de durée
\DeclareSIUnit\microliter{\micro\liter} % le modèle écrit parfois \microliter en un seul token
\DeclareSIUnit\million{millions} % dénombrement (ex. \SI{15}{\million\per\mL} spermatozoïdes)
\usepackage{qrcode}
% \degree : commande fréquemment utilisée par les modèles pour un angle ou une
% température en dehors de \SI{}{} (non fournie par les paquets chargés).
\providecommand{\degree}{^{\circ}}
% \qed (fin de démonstration), utilisé par les modèles sans amsthm
\providecommand{\qed}{\hfill$\square$}
% \barre : double barre des valeurs interdites dans un tableau de variations (tkz-tab)
\providecommand{\barre}{\ensuremath{\|}}
% environnement proof (amsthm non chargé) utilisé par les modèles
\ifdefined\proof\else\newenvironment{proof}[1][Démonstration]{\par\noindent\textit{#1.}\ }{\qed\par}\fi
\usepackage{lastpage}
\usepackage{hyperref}
\hypersetup{hidelinks}

% ============================================================
% Palette « Pop » OnBuch+ — mêmes hex que la classe P (plus_ui.dart)
% ============================================================
\definecolor{poporange}{HTML}{F59321}
\definecolor{popdark}{HTML}{E07A0C}
\definecolor{popcream}{HTML}{FFF6E8}
\definecolor{popink}{HTML}{1C1714}
\definecolor{poppurple}{HTML}{7A5AE0}
\definecolor{popgreen}{HTML}{2E9E6A}
\definecolor{poppink}{HTML}{E0457B}
\definecolor{popblue}{HTML}{2F7DE1}
\definecolor{popgold}{HTML}{C98A12}
\definecolor{popmuted}{HTML}{8A7F76}
\definecolor{popline}{HTML}{EADFD2}
\definecolor{popgray}{HTML}{8A7F76}
\colorlet{popgrey}{popgray}
% Alias tolérants (le modèle nomme parfois les couleurs différemment)
\colorlet{popwhite}{white}
\colorlet{popblack}{popink}
\colorlet{popred}{poppink}
\colorlet{popyellow}{popgold}
% Teintes claires (fonds de tableaux/encadrés) — le modèle nomme parfois
% les couleurs avec un suffixe L (ex. popblueL) sans les avoir définies.
\colorlet{poporangeL}{poporange!20}
\colorlet{popdarkL}{popdark!20}
\colorlet{poppurpleL}{poppurple!20}
\colorlet{popgreenL}{popgreen!20}
\colorlet{poppinkL}{poppink!20}
\colorlet{popblueL}{popblue!20}
\colorlet{popgoldL}{popgold!20}
\colorlet{popredL}{popred!20}
\colorlet{popgrayL}{popgray!20}
% Teintes claires pour fonds de tableaux / zones
\colorlet{popblueL}{popblue!10!white}
\colorlet{poporangeL}{poporange!14!white}
\colorlet{popgreenL}{popgreen!12!white}
\colorlet{poppurpleL}{poppurple!12!white}
\colorlet{poppinkL}{poppink!12!white}
\colorlet{popgoldL}{popgold!14!white}

\pagecolor{popcream}

% ============================================================
% Typographie
% ============================================================
\defaultfontfeatures{Ligatures=TeX}
\setmainfont{PlusJakartaSans-Medium.ttf}[Path=../../../fonts/,BoldFont=PlusJakartaSans-Bold.ttf]
% Maths en sans-serif (Fira Math) avec chiffres et lettres droites de Plus
% Jakarta, pour que les formules se fondent dans le texte.
\usepackage[math-style=ISO,bold-style=ISO]{unicode-math}
\setmathfont{FiraMath-Regular.otf}[Path=../../../fonts/]
\setmathfont{PlusJakartaSans-Medium.ttf}[Path=../../../fonts/,range={up/num,up/{Latin,latin},\mathup}]
\usepackage{icomma} % virgule décimale sans espace en mode math
% Caractères mathématiques Unicode absents des polices de texte (Plus Jakarta,
% Archivo Black) : rendus par leur équivalent mathématique, en texte comme en
% formule — sinon ils s'affichent en carré vide (ex. « Arithmétique dans ℤ »).
\usepackage{newunicodechar}
\newunicodechar{ℝ}{\ensuremath{\mathbb{R}}}
\newunicodechar{ℤ}{\ensuremath{\mathbb{Z}}}
\newunicodechar{ℂ}{\ensuremath{\mathbb{C}}}
\newunicodechar{ℕ}{\ensuremath{\mathbb{N}}}
\newunicodechar{ℚ}{\ensuremath{\mathbb{Q}}}
\newunicodechar{𝔻}{\ensuremath{\mathbb{D}}}
\newunicodechar{α}{\ensuremath{\alpha}}
\newunicodechar{β}{\ensuremath{\beta}}
\newunicodechar{γ}{\ensuremath{\gamma}}
\newunicodechar{δ}{\ensuremath{\delta}}
\newunicodechar{ε}{\ensuremath{\varepsilon}}
\newunicodechar{θ}{\ensuremath{\theta}}
\newunicodechar{λ}{\ensuremath{\lambda}}
\newunicodechar{μ}{\ensuremath{\mu}}
\newunicodechar{σ}{\ensuremath{\sigma}}
\newunicodechar{τ}{\ensuremath{\tau}}
\newunicodechar{φ}{\ensuremath{\varphi}}
\newunicodechar{ψ}{\ensuremath{\psi}}
\newunicodechar{ω}{\ensuremath{\omega}}
\newunicodechar{Σ}{\ensuremath{\Sigma}}
% majuscules produites par \MakeUppercase dans les titres (α-aminés -> Α)
\newunicodechar{Α}{\ensuremath{\alpha}}
\newunicodechar{Β}{\ensuremath{\beta}}
\newunicodechar{Γ}{\ensuremath{\Gamma}}
\newunicodechar{Δ}{\ensuremath{\Delta}}
\newunicodechar{Ε}{\ensuremath{\varepsilon}}
\newunicodechar{Θ}{\ensuremath{\Theta}}
\newunicodechar{Λ}{\ensuremath{\Lambda}}
\newunicodechar{Μ}{\ensuremath{\mu}}
\newunicodechar{Τ}{\ensuremath{\tau}}
\newunicodechar{Φ}{\ensuremath{\Phi}}
\newunicodechar{Ψ}{\ensuremath{\Psi}}
\newunicodechar{Ω}{\ensuremath{\Omega}}
\newunicodechar{↦}{\ensuremath{\mapsto}}
\newunicodechar{∈}{\ensuremath{\in}}
\newunicodechar{∉}{\ensuremath{\notin}}
\newunicodechar{∧}{\ensuremath{\wedge}}
\newunicodechar{⇔}{\ensuremath{\Leftrightarrow}}
\newunicodechar{⟺}{\ensuremath{\Longleftrightarrow}}
\newunicodechar{⇒}{\ensuremath{\Rightarrow}}
\newunicodechar{⟹}{\ensuremath{\Longrightarrow}}
\newunicodechar{∘}{\ensuremath{\circ}}
\newunicodechar{∪}{\ensuremath{\cup}}
\newunicodechar{∩}{\ensuremath{\cap}}
\newunicodechar{≡}{\ensuremath{\equiv}}
\newunicodechar{⊂}{\ensuremath{\subset}}
\newunicodechar{⊥}{\ensuremath{\perp}}
\newunicodechar{∃}{\ensuremath{\exists}}
\newunicodechar{∀}{\ensuremath{\forall}}
\newunicodechar{∛}{\ensuremath{\sqrt[3]{\,}}}
\newunicodechar{∜}{\ensuremath{\sqrt[4]{\,}}}
\newunicodechar{⃗}{\ensuremath{{}^{\to}}}
\newunicodechar{ⁿ}{\ifmmode^{n}\else\textsuperscript{\ensuremath{n}}\fi}
\newunicodechar{ˣ}{\ifmmode^{x}\else\textsuperscript{\ensuremath{x}}\fi}
\newunicodechar{ᵇ}{\ifmmode^{b}\else\textsuperscript{\ensuremath{b}}\fi}
\newunicodechar{ʳ}{\ifmmode^{r}\else\textsuperscript{\ensuremath{r}}\fi}
\newunicodechar{ᵘ}{\ifmmode^{u}\else\textsuperscript{\ensuremath{u}}\fi}
\newunicodechar{ʸ}{\ifmmode^{y}\else\textsuperscript{\ensuremath{y}}\fi}
\newunicodechar{ᵃ}{\ifmmode^{a}\else\textsuperscript{\ensuremath{a}}\fi}
\newunicodechar{ᵉ}{\ifmmode^{e}\else\textsuperscript{\ensuremath{e}}\fi}
\newunicodechar{ᵏ}{\ifmmode^{k}\else\textsuperscript{\ensuremath{k}}\fi}
\newunicodechar{ᵖ}{\ifmmode^{p}\else\textsuperscript{\ensuremath{p}}\fi}
\newunicodechar{ᴺ}{\ifmmode^{N}\else\textsuperscript{\ensuremath{N}}\fi}
\newunicodechar{ᶜ}{\ifmmode^{c}\else\textsuperscript{\ensuremath{c}}\fi}
\newunicodechar{ʰ}{\ifmmode^{h}\else\textsuperscript{\ensuremath{h}}\fi}
\newunicodechar{ᵠ}{\ifmmode^{q}\else\textsuperscript{\ensuremath{q}}\fi}
\newunicodechar{ₙ}{\ifmmode_{n}\else\textsubscript{\ensuremath{n}}\fi}
\newunicodechar{ₐ}{\ifmmode_{a}\else\textsubscript{\ensuremath{a}}\fi}
\newunicodechar{ᵢ}{\ifmmode_{i}\else\textsubscript{\ensuremath{i}}\fi}
\newunicodechar{ᵣ}{\ifmmode_{r}\else\textsubscript{\ensuremath{r}}\fi}
\newunicodechar{ₖ}{\ifmmode_{k}\else\textsubscript{\ensuremath{k}}\fi}
\newunicodechar{ᵦ}{\ifmmode_{\beta}\else\textsubscript{\ensuremath{\beta}}\fi}
\newunicodechar{ₓ}{\ifmmode_{x}\else\textsubscript{\ensuremath{x}}\fi}
\newfontfamily\popdisplay{ArchivoBlack-Regular.ttf}[Path=../../../fonts/]
\newfontfamily\poplogo{SpaceGrotesk-Bold.ttf}[Path=../../../fonts/]
\newfontfamily\popsemi{PlusJakartaSans-SemiBold.ttf}[Path=../../../fonts/]
\newfontfamily\popxbold{PlusJakartaSans-ExtraBold.ttf}[Path=../../../fonts/]
\newfontfamily\popxboldls{PlusJakartaSans-ExtraBold.ttf}[Path=../../../fonts/,LetterSpace=3.0]

\setlist[enumerate]{leftmargin=1.7em,itemsep=5pt,topsep=4pt}
\setlist[itemize]{leftmargin=1.7em,itemsep=4pt,topsep=4pt,label={\color{poporange}\textbullet}}
\setlength{\parindent}{0pt}
\setlength{\parskip}{5pt}
\renewcommand{\arraystretch}{1.25}

% pgfplots : style Pop par défaut
\pgfplotsset{
  every axis/.append style={axis line style={popink,line width=1pt},tick style={popink},
    grid style={popline,line width=0.5pt},label style={font=\small},tick label style={font=\footnotesize}},
  popcurve/.style={poporange,line width=1.8pt,no markers,smooth},
}
% Même style disponible dans un \draw TikZ simple (hors pgfplots)
\tikzset{popcurve/.style={poporange,line width=1.8pt}}
\tikzset{popgrid/.style={popline,very thin}}
\colorlet{popcurve}{poporange} % le nom du style est parfois utilisé comme couleur

% ============================================================
% Pill / squiggle
% ============================================================
\newcommand{\poppill}[3]{%
  \tikz[baseline=-0.55ex]\node[draw=popink,line width=1.2pt,rounded corners=2.6pt,%
    fill=#2,inner xsep=6.5pt,inner ysep=3pt]{%
    \popxboldls\fontsize{7}{7}\selectfont\color{#3}\MakeUppercase{#1}};%
}
\newcommand{\popsquiggle}[1]{%
  \tikz[baseline=-0.4ex]{
    \draw[#1,line width=2.6pt,line cap=round]
      plot [smooth] coordinates {(0,0.09) (0.34,0.01) (0.68,0.21) (1.02,0.01) (1.36,0.10)};
  }%
}
% Mot-clé surligné (marqueur orange sous le texte)
\newcommand{\cle}[1]{{\popxbold\color{popdark}#1}}

% ============================================================
% En-tête / pied
% ============================================================
\pagestyle{fancy}
\fancyhf{}
\renewcommand{\headrulewidth}{0pt}
\renewcommand{\footrulewidth}{0pt}
\lhead{\raisebox{-0.15ex}{\poplogo\fontsize{13}{13}\selectfont\color{popink}OnBuch\color{poporange}+}}
\rhead{\poppill{\DOCMATIERE\ \textbullet\ \DOCNIVEAU\ \textbullet\ Leçon \DOCLECON}{white}{popink}}
\lfoot{\popxbold\fontsize{7.6}{7.6}\selectfont\color{popmuted}\MakeUppercase{Cours signé OnBuch+}\ \textbullet\ {\popsemi\DOCTITRE}}
\rfoot{\tikz[baseline=-0.6ex]\node[rounded corners=8.5pt,fill=popink,minimum height=17pt,minimum width=17pt,inner xsep=5pt,inner sep=0]{\popxbold\fontsize{8}{8}\selectfont\color{popcream}\thepage\,/\,\pageref*{LastPage}};}

% ============================================================
% Titres
% ============================================================
\newcommand{\chaptitre}[2]{%
  \begin{tcolorbox}[enhanced,colback=poporange,colframe=popink,boxrule=2.6pt,arc=11pt,
      drop shadow=popink,left=15pt,right=15pt,top=12pt,bottom=11pt,before skip=3pt,after skip=18pt]
    \poppill{#2}{popink}{popcream}\par\vspace{9pt}
    {\raggedright\hyphenpenalty=10000\exhyphenpenalty=10000\popdisplay\fontsize{22}{24}\selectfont\color{popcream}\MakeUppercase{#1}\par}
  \end{tcolorbox}
}
% Section numérotée : pastille orange avec le numéro + titre Archivo
\newcounter{popsec}
\newcommand{\coursec}[1]{%
  \stepcounter{popsec}\setcounter{popsub}{0}%
  \par\vspace{16pt}\noindent
  \begin{minipage}{\linewidth}
  \tikz[baseline=-0.7ex]\node[draw=popink,line width=1.6pt,fill=poporange,rounded corners=4pt,minimum size=22pt,inner sep=0]{\popdisplay\fontsize{12}{12}\selectfont\color{popcream}\thepopsec};%
  \hspace{8pt}{\popdisplay\fontsize{15}{17}\selectfont\color{popink}\MakeUppercase{#1}}\par
  \vspace{4pt}\noindent{\color{poporange}\rule{36pt}{3.6pt}}
  \end{minipage}\par\nopagebreak\vspace{8pt}
}
\newcounter{popsub}
\newcommand{\courssub}[1]{%
  \stepcounter{popsub}%
  \par\vspace{9pt}\noindent{\popxbold\fontsize{11.5}{13}\selectfont\color{popink}{\color{poporange}\thepopsec.\thepopsub}\hspace{6pt}#1}\par\nopagebreak\vspace{4pt}
}

% ============================================================
% Boîtes pédagogiques — même recette : fond blanc, bordure épaisse colorée,
% ombre dure encre, badge plein. Titre optionnel [#1].
% ============================================================
\tcbset{popbase/.style={breakable,enhanced,colback=white,boxrule=2.3pt,arc=9pt,
  shadow={3pt}{-3pt}{0pt}{popink},left=14pt,right=14pt,top=11pt,bottom=11pt,before skip=11pt,after skip=15pt,
  fontupper=\popsemi\fontsize{10.6}{14.5}\selectfont\color{popink},
  fontlower=\popsemi\fontsize{10.6}{14.5}\selectfont\color{popink}}}
% Coche dessinée (le glyphe ✓ n'existe pas dans les polices du document)
\newcommand{\popcheck}[1]{\tikz[baseline=-0.5ex]\draw[#1,line width=1.6pt,line cap=round,line join=round](-0.13,0.0)--(-0.04,-0.09)--(0.13,0.1);}
\providecommand{\coche}{\popcheck{popgreen}}
\providecommand{\resultat}[1]{\par\smallskip\noindent\fcolorbox{popgreen}{popgreenL}{\parbox{\dimexpr\linewidth-2\fboxsep-2\fboxrule\relax}{\textbf{Résultat.} #1}}\par\smallskip}
\newcommand{\popboxhead}[3]{\poppill{#1}{#2}{white}\ifx\relax#3\relax\else\hspace{6pt}{\popxbold\fontsize{10.6}{12}\selectfont\color{popink}#3}\fi\par\vspace{7pt}}

\newtcolorbox{aretenir}[1][]{popbase,colframe=popgreen,before upper={\popboxhead{À retenir}{popgreen}{#1}}}
\newtcolorbox{exemplebox}[1][]{popbase,colframe=poporange,before upper={\popboxhead{Exemple}{poporange}{#1}}}
\newtcolorbox{definition}[1][]{popbase,colframe=popblue,before upper={\popboxhead{Définition}{popblue}{#1}}}
\newtcolorbox{propriete}[1][]{popbase,colframe=poppurple,before upper={\popboxhead{Propriété}{poppurple}{#1}}}
\newtcolorbox{methode}[1][]{popbase,colframe=popgold,before upper={\popboxhead{Méthode}{popgold}{#1}}}
\newtcolorbox{attention}[1][]{popbase,colframe=poppink,before upper={\popboxhead{Attention}{poppink}{#1}}}
\newtcolorbox{experience}[1][]{popbase,colframe=popblue,colback=popblueL,before upper={\popboxhead{Expérience}{popblue}{#1}}}
% « Le savais-tu ? » : carte violette pleine (contexte, culture, Cameroun)
\newtcolorbox{savaistu}[1][]{popbase,colback=poppurple,colframe=popink,
  fontupper=\popsemi\fontsize{10.4}{14.2}\selectfont\color{white},
  before upper={\poppill{Le savais-tu\,?}{white}{poppurple}\ifx\relax#1\relax\else\hspace{6pt}{\popxbold\color{white}#1}\fi\par\vspace{7pt}}}
% Exercice résolu : énoncé en haut, \tcblower puis la solution sur fond crème
\newcounter{popexo}\newcounter{popexr}
\newtcolorbox{exoresolu}[1][]{popbase,colframe=poporange,colbacklower=popcream,
  segmentation style={popink,line width=1.2pt,dashed},
  before upper={\stepcounter{popexr}\popboxhead{Exercice résolu \thepopexr}{poporange}{#1}},
  before lower={\poppill{Solution}{popink}{popcream}\par\vspace{6pt}}}
% Exercice (non résolu en place) — la correction est dans la partie Corrigés
\newtcolorbox{exercice}[1][]{popbase,colframe=popink,
  before upper={\stepcounter{popexo}\popboxhead{Exercice \thepopexo}{popink}{#1}}}
% Corrigé
\newtcolorbox{corrige}[1][]{popbase,colframe=popgreen,colback=popgreenL,
  before upper={\popboxhead{Corrigé}{popgreen}{#1}}}
% Objectifs / prérequis / bilan
\newtcolorbox{objectifs}{popbase,colframe=popink,colback=poporangeL,
  before upper={\popboxhead{Objectifs de la leçon}{popink}{}}}
\newtcolorbox{prerequis}{popbase,colframe=popink,colback=popblueL,
  before upper={\popboxhead{Prérequis}{popblue}{}}}
\newtcolorbox{bilan}{popbase,colframe=popink,colback=white,boxrule=2.8pt,
  before upper={\popboxhead{Fiche bilan}{poporange}{L'essentiel à connaître pour l'examen}}}
% Figure encadrée avec légende
\newtcolorbox{popfigure}[1][]{enhanced,breakable=false,colback=white,colframe=popline,boxrule=1.4pt,arc=8pt,
  left=10pt,right=10pt,top=10pt,bottom=8pt,before skip=10pt,after skip=12pt,halign=center,
  lower separated=false,halign lower=center,
  fontlower=\popsemi\fontsize{9}{11.5}\selectfont\color{popmuted},
  #1}
\newcommand{\legende}[1]{\par\vspace{6pt}{\popsemi\fontsize{9}{11.5}\selectfont\color{popmuted}#1}}

% ============================================================
% Signature OnBuch+
% ============================================================
\newcommand{\poplogoinline}[1]{{\poplogo\fontsize{#1}{#1}\selectfont\color{popink}OnBuch\color{poporange}+}}
\newcommand{\signatureonbuch}{%
  \begin{tcolorbox}[enhanced,colback=popink,colframe=popink,boxrule=2.2pt,arc=10pt,
      drop shadow=poporange,left=14pt,right=14pt,top=10pt,bottom=10pt,before skip=0pt,after skip=0pt]
    \tikz[baseline=-0.7ex]\node[circle,fill=popgreen,draw=popcream,line width=1.3pt,minimum size=22pt,inner sep=0]{\popcheck{white}};%
    \hspace{9pt}\begin{minipage}[c]{0.62\linewidth}
      {\popxbold\fontsize{12}{13}\selectfont\color{popcream}Signé\ \textbullet\ {\poplogo OnBuch\color{poporange}+}}\par\vspace{2pt}
      {\raggedright\popsemi\fontsize{8}{10.5}\selectfont\color{popline}\MakeUppercase{Équipe pédagogique OnBuch+}\\\MakeUppercase{Conforme au programme officiel MINESEC}\par}
    \end{minipage}\hfill
    {\poplogo\fontsize{22}{22}\selectfont\color{popcream}OnBuch\color{poporange}+}
  \end{tcolorbox}%
}

% ============================================================
% Page de garde
% #1 titre · #2 matière · #3 niveau · #4 module · #5 n° leçon · #6 durée ·
% #7 description / objectifs (texte ou itemize)
% ============================================================
\newcommand{\pagegarde}[7]{%
  \thispagestyle{empty}
  \newgeometry{a4paper,margin=1.7cm,top=1.6cm,bottom=1.4cm}%
  \begin{tikzpicture}[remember picture,overlay]
    \fill[poporange!18] ([xshift=-3.2cm,yshift=-3.6cm]current page.north east) circle (4.6cm);
    \fill[poppurple!14] ([xshift=2.4cm,yshift=7.6cm]current page.south west) circle (3.8cm);
  \end{tikzpicture}%
  {\poplogo\fontsize{30}{30}\selectfont\color{popink}OnBuch\color{poporange}+}\hfill
  \raisebox{0.6ex}{\poppill{Cours signé \textbullet\ Premium}{popink}{popcream}}\par
  \vspace{1pt}\popsquiggle{poppurple}\par
  \vspace{8pt}
  \begin{tcolorbox}[enhanced,colback=poporange,colframe=popink,boxrule=2.8pt,arc=13pt,
      drop shadow=popink,left=18pt,right=18pt,top=12pt,bottom=13pt,after skip=10pt]
    \poppill{#2 \textbullet\ #3}{popink}{popcream}\hspace{5pt}\poppill{#4}{white}{popink}\par\vspace{8pt}
    {\popdisplay\fontsize{14}{14}\selectfont\color{popink}LEÇON #5}\par\vspace{4pt}
    {\raggedright\hyphenpenalty=10000\exhyphenpenalty=10000\popdisplay\fontsize{25}{27}\selectfont\color{popcream}\MakeUppercase{#1}\par}
    \vspace{8pt}
    {\popxbold\fontsize{9.5}{9.5}\selectfont\color{popink}\MakeUppercase{Durée conseillée : #6}}
  \end{tcolorbox}
  \begin{tcolorbox}[enhanced,colback=white,colframe=popink,boxrule=2.2pt,arc=10pt,
      drop shadow=popink,left=16pt,right=16pt,top=10pt,bottom=10pt,after skip=8pt]
    \poppill{Dans cette leçon}{poppurple}{white}\par\vspace{5pt}
    \popsemi\fontsize{9.3}{12.6}\selectfont\color{popink}#7
  \end{tcolorbox}
  \noindent\begin{minipage}[t]{0.487\linewidth}
    \begin{tcolorbox}[enhanced,colback=popgreen,colframe=popink,boxrule=2.2pt,arc=10pt,
        drop shadow=popink,left=11pt,right=11pt,top=10pt,bottom=10pt,height=3.1cm,valign=center]
      \tcbox[colback=white,colframe=popink,boxrule=1.3pt,arc=5pt,boxsep=0pt,nobeforeafter,
        left=3pt,right=3pt,top=3pt,bottom=3pt,valign=center]{\qrcode[height=1.9cm]{https://play.google.com/store/apps/details?id=com.luvvix.onbuch&hl=fr}}%
      \hspace{8pt}\begin{minipage}[c]{0.5\linewidth}
        {\popdisplay\fontsize{11}{12.5}\selectfont\color{white}\MakeUppercase{Télécharge OnBuch}}\par\vspace{4pt}
        {\raggedright\popsemi\fontsize{8.4}{11}\selectfont\color{white}Gratuit \textbullet\ Google Play\\Tuteur IA, annales, concours, résultats\par}
      \end{minipage}
    \end{tcolorbox}
  \end{minipage}\hfill
  \begin{minipage}[t]{0.487\linewidth}
    \begin{tcolorbox}[enhanced,colback=poppurple,colframe=popink,boxrule=2.2pt,arc=10pt,
        drop shadow=popink,left=11pt,right=11pt,top=10pt,bottom=10pt,height=3.1cm,valign=center]
      \tikz[baseline=-0.7ex]\node[circle,fill=white,draw=popink,line width=1.3pt,minimum size=1.6cm,inner sep=0]{\popdisplay\fontsize{24}{24}\selectfont\color{poppurple}+};%
      \hspace{8pt}\begin{minipage}[c]{0.6\linewidth}
        {\popdisplay\fontsize{11}{12.5}\selectfont\color{white}\MakeUppercase{Passe à OnBuch+}}\par\vspace{4pt}
        {\raggedright\popsemi\fontsize{8.4}{11}\selectfont\color{white}Tous les cours signés, Tuteur IA illimité, fiches PDF — onglet OnBuch+ de l'app\par}
      \end{minipage}
    \end{tcolorbox}
  \end{minipage}
  \vspace{8pt}
  \popcontenu
  \vfill
  \signatureonbuch
  \restoregeometry
  \newpage
}

% Rangée « Ce cours contient » (page de garde)
\newcommand{\poptile}[3]{% #1 couleur #2 gros texte #3 légende
  \begin{tcolorbox}[enhanced,colback=#1,colframe=popink,boxrule=2pt,arc=9pt,shadow={2.5pt}{-2.5pt}{0pt}{popink},
      left=6pt,right=6pt,top=9pt,bottom=9pt,halign=center,valign=center,height=1.75cm,width=\linewidth,nobeforeafter]
    {\popdisplay\fontsize{12.5}{13}\selectfont\color{white}\MakeUppercase{#2}}\par\vspace{3pt}
    {\centering\popsemi\fontsize{7.8}{9.5}\selectfont\color{white}#3\par}
  \end{tcolorbox}}
\newcommand{\popcontenu}{%
  \par\noindent{\popxboldls\fontsize{7.5}{7.5}\selectfont\color{popmuted}CE COURS SIGNÉ CONTIENT}\par\vspace{7pt}
  \noindent\begin{minipage}[t]{0.235\linewidth}\poptile{popblue}{Cours}{complet, illustré, pas à pas}\end{minipage}\hfill
  \begin{minipage}[t]{0.235\linewidth}\poptile{poporange}{Méthodes}{et exercices résolus}\end{minipage}\hfill
  \begin{minipage}[t]{0.235\linewidth}\poptile{poppink}{Exercices}{type examen + corrigés}\end{minipage}\hfill
  \begin{minipage}[t]{0.235\linewidth}\poptile{popgreen}{Fiche bilan}{l'essentiel à retenir}\end{minipage}\par}

% Sommaire
\newtcolorbox{sommaire}{enhanced,colback=white,colframe=popink,boxrule=2.2pt,arc=10pt,
  drop shadow=popink,left=18pt,right=18pt,top=13pt,bottom=13pt,before skip=6pt,after skip=20pt,
  before upper={\poppill{Sommaire}{poppurple}{white}\par\vspace{9pt}\popsemi\fontsize{10.6}{14.5}\selectfont\color{popink}}}

% Signature de fin de cours — poussée tout en bas de la dernière page
\newcommand{\findecours}{%
  \par\vfill\nopagebreak
  \begin{center}\popsquiggle{poporange}\end{center}\vspace{4pt}
  \signatureonbuch
}
\AtBeginDocument{\def\llbracket{\mathopen{[\mkern-3.5mu[}}\def\rrbracket{\mathclose{]\mkern-3.5mu]}}} % intervalles d'entiers (glyphe ⟦ absent de la police)
% Glyphes absents de Fira Math : redéfinis à partir de glyphes disponibles
\AtBeginDocument{%
  \def\setminus{\mathbin{\backslash}}\let\smallsetminus\setminus
  \def\vdots{\mathord{\vbox{\baselineskip3.2pt\lineskiplimit0pt\kern4pt\hbox{.}\hbox{.}\hbox{.}}}}%
  \def\ddots{\mathinner{\mkern1mu\raise6.5pt\hbox{.}\mkern2mu\raise3.6pt\hbox{.}\mkern2mu\raise0.7pt\hbox{.}\mkern1mu}}%
  \def\triangle{\mathord{\scalebox{0.9}{$\symup{\Delta}$}}}%
  \def\nabla{\mathord{\raisebox{\depth}{\scalebox{1}[-1]{$\symup{\Delta}$}}}}%
  \def\checkmark{\popcheck{popdark}}%
  \def\star{\mathbin{\ast}}\def\bigstar{\mathord{\ast}}%
  \def\diamond{\mathbin{\scalebox{0.6}{$\Diamond$}}}%
  \def\bigcup{\mathop{\mathchoice{\raisebox{-0.3em}{\scalebox{1.7}{$\cup$}}}{\scalebox{1.25}{$\cup$}}{\cup}{\cup}}\displaylimits}%
  \def\bigcap{\mathop{\mathchoice{\raisebox{-0.3em}{\scalebox{1.7}{$\cap$}}}{\scalebox{1.25}{$\cap$}}{\cap}{\cap}}\displaylimits}%
  \def\lesseqqgtr{\mathrel{\lessgtr}}\def\bigoplus{\mathop{\oplus}\displaylimits}%
}
\providecommand{\ssi}{si et seulement si} % abréviation fréquente
\AtBeginDocument{\providecommand{\wideparen}{\overparen}} % arc de cercle

\begin{document}

\pagegarde{Quelques machines simples}{PCT}{3ème}{Module 2 : Électricité}{N11}{2 séances (1 h chacune)}{Cette leçon présente les principales machines simples (levier, plan incliné, poulie, treuil), leurs définitions, leurs avantages et leurs inconvénients. À travers des situations concrètes de la vie quotidienne camerounaise, l'élève apprend à identifier ces dispositifs, à réaliser des expériences et des calculs simples, et à respecter les consignes de sécurité.
\vspace{4pt}
\begin{multicols}{2}\small
\begin{itemize}
\item À la fin de cette leçon, je sais définir une machine simple et citer ses rôles
\item identifier les machines simples présentes dans des objets et situations de la vie courante
\item définir le levier, distinguer ses trois éléments (point d'appui, effort, résistance) et citer des exemples
\item définir le plan incliné et expliquer pourquoi il facilite la montée d'une charge
\item définir la poulie (fixe et mobile) et le treuil, et décrire leur fonctionnement
\item réaliser et interpréter une expérience simple mettant en évidence l'avantage d'une machine simple
\end{itemize}\end{multicols}}

\begin{sommaire}
\begin{multicols}{2}
\begin{enumerate}
\item Notion de machine simple
\item Le levier
\item Le plan incliné
\item La poulie
\item Le treuil
\item Bilan comparatif et sécurité
\item Activité d'intégration
\item Méthodes et exercices résolus
\item Exercices
\item Corrigés des exercices
\item Fiche bilan
\end{enumerate}
\end{multicols}
\end{sommaire}

\begin{objectifs}
\begin{itemize}
\item À la fin de cette leçon, je sais définir une machine simple et citer ses rôles
\item identifier les machines simples présentes dans des objets et situations de la vie courante
\item définir le levier, distinguer ses trois éléments (point d'appui, effort, résistance) et citer des exemples
\item définir le plan incliné et expliquer pourquoi il facilite la montée d'une charge
\item définir la poulie (fixe et mobile) et le treuil, et décrire leur fonctionnement
\item réaliser et interpréter une expérience simple mettant en évidence l'avantage d'une machine simple
\item comparer les avantages et les inconvénients de chaque machine simple
\item effectuer un calcul simple sur un levier en équilibre (moments des forces)
\item appliquer les consignes de sécurité lors de l'utilisation des machines simples
\end{itemize}
\end{objectifs}

\begin{prerequis}
\begin{itemize}
\item Notion de force : effets d'une force, unité (newton), mesure au dynamomètre
\item Le poids d'un corps : relation P = m × g
\item Notion d'équilibre d'un solide soumis à des forces
\item Schématisation d'une force (point d'application, direction, sens, intensité)
\end{itemize}
\end{prerequis}

\newpage

\coursec{Notion de machine simple}

\textbf{Situation d'accroche.} Au marché Mokolo, à Yaoundé, une vendeuse soulève sans effort un sac de \SI{50}{kg} de manioc à l'aide d'une simple barre de fer posée sur une pierre. Sur un chantier de la même ville, des ouvriers hissent des seaux de béton au troisième étage grâce à une roue et une corde. Comment de simples objets permettent-ils de soulever des charges que nos muscles seuls ne pourraient déplacer ? Ces dispositifs, utilisés depuis des millénaires, sont appelés \cle{machines simples}. Cette leçon nous apprend à les identifier, à comprendre leur fonctionnement et à les utiliser en toute sécurité.

\textbf{Problématique.} Qu'est-ce qu'une machine simple ? Quels sont ses rôles ? Pourquoi dit-on qu'elle ne crée pas d'énergie ?

\courssub{Constat expérimental : l'effort diminue avec un dispositif}

\begin{experience}[Soulever une charge avec ou sans dispositif]
\textbf{Matériel :} un seau rempli d'eau (ou un sac de sable), un dynamomètre, une planche solide, une corde passant sur une barre fixe.

\textbf{Protocole :}
\begin{enumerate}
\item Accroche le seau au dynamomètre et soulève-le directement à la verticale. Note la force $F_1$ indiquée.
\item Pose la planche en appui sur une marche pour former une rampe, puis tire le seau le long de la planche avec le dynamomètre. Note la force $F_2$.
\item Passe la corde sur la barre fixe, attache le seau à une extrémité et tire l'autre extrémité vers le bas avec le dynamomètre. Note la force $F_3$.
\end{enumerate}

\textbf{Observations :} la force $F_1$ est grande (elle est égale au poids du seau). La force $F_2$ est nettement plus petite que $F_1$. La force $F_3$ est voisine de $F_1$, mais elle s'exerce vers le bas, ce qui est plus commode : on utilise le poids de son propre corps.

\textbf{Interprétation :} la planche et la corde n'ont pas changé le poids du seau. Elles ont changé les conditions dans lesquelles on applique la force : la planche \emph{réduit l'intensité} de l'effort, la corde passant sur la barre en \emph{change la direction}.

\textbf{Conclusion :} certains dispositifs mécaniques permettent d'accomplir le même travail (soulever la même charge à la même hauteur) avec un effort plus faible ou mieux orienté. Ce sont des machines simples.
\end{experience}

\begin{popfigure}
\begin{center}
\begin{tikzpicture}[scale=0.9, >=Stealth]
% --- Situation 1 : levée directe ---
\draw[popline, thick] (-1.6,0) -- (1.6,0);
\fill[poppurpleL] (-0.5,0) rectangle (0.5,0.9);
\node[popdark, font=\small] at (0,0.45) {charge};
\draw[->, poporange, line width=2.2pt] (0,0.9) -- (0,2.6);
\node[poporange, font=\small\bfseries, right] at (0.05,2.2) {$F_1$ : grand effort};
% bonhomme
\draw[popdark, thick] (1.15,0) -- (1.15,0.8);
\draw[popdark, thick] (1.15,0.8) -- (0.9,0.35);
\draw[popdark, thick] (1.15,0.8) -- (1.4,0.35);
\draw[popdark, thick] (1.15,0.8) -- (0.75,1.15);
\draw[popdark, thick] (1.15,0.8) -- (1.5,1.15);
\fill[popdark] (1.15,1.05) circle (0.16);
\node[popdark, font=\small\bfseries] at (0,-0.55) {Levée directe};

% --- Situation 2 : avec levier ---
\begin{scope}[xshift=8.2cm]
\draw[popline, thick] (-1.8,0) -- (2.6,0);
\fill[poppurpleL] (-1.4,0.35) rectangle (-0.4,0.95);
\node[popdark, font=\small] at (-0.9,0.65) {charge};
\fill[popgold] (0.5,0) -- (0.2,0.45) -- (0.8,0.45) -- cycle;
\draw[popblue, line width=3pt] (-1.7,0.95) -- (2.3,0.15);
\draw[->, popgreen, line width=2.2pt] (2.3,0.15) -- (2.3,-1.05);
\node[popgreen, font=\small\bfseries, right] at (2.35,-0.5) {$F_2$ : petit effort};
% bonhomme
\draw[popdark, thick] (3.1,0) -- (3.1,0.8);
\draw[popdark, thick] (3.1,0.8) -- (2.85,0.35);
\draw[popdark, thick] (3.1,0.8) -- (3.35,0.35);
\draw[popdark, thick] (3.1,0.8) -- (2.75,0.5);
\draw[popdark, thick] (3.1,0.8) -- (3.45,0.5);
\fill[popdark] (3.1,1.05) circle (0.16);
\node[popdark, font=\small\bfseries] at (0.4,-0.55) {Avec une machine simple (levier)};
\end{scope}
\end{tikzpicture}
\end{center}
\legende{Comparaison de l'effort à fournir : à gauche, soulever directement la charge exige une grande force $F_1$ ; à droite, avec une barre appuyée sur une pierre (levier), une petite force $F_2$ suffit.}
\end{popfigure}

\courssub{Définition d'une machine simple}

\begin{definition}[Machine simple]
Une \cle{machine simple} est un dispositif mécanique qui permet d'accomplir un travail en modifiant \textbf{l'intensité} ou \textbf{la direction} de la force à appliquer.
\end{definition}

Décomposons cette définition pour bien la comprendre :

\begin{itemize}
\item \textbf{« Dispositif mécanique »} : il s'agit d'un objet matériel (barre, planche, corde, roue, manivelle) qui transmet un mouvement ou une force. Il n'a pas de moteur et ne consomme ni carburant ni électricité.
\item \textbf{« Accomplir un travail »} : en physique, on dit qu'une force travaille lorsqu'elle déplace son point d'application. Soulever un sac de manioc, monter un seau de béton, enfoncer un clou : ce sont des travaux.
\item \textbf{« Modifier l'intensité de la force »} : la machine permet d'utiliser une force plus petite que le poids de la charge (c'est le cas du levier, du plan incliné, du treuil).
\item \textbf{« Modifier la direction de la force »} : la machine permet d'orienter l'effort dans une direction plus pratique (tirer vers le bas au lieu de pousser vers le haut, comme avec la poulie d'un puits).
\end{itemize}

\begin{attention}[Une machine simple n'est pas un moteur]
Ne confonds pas \emph{machine simple} et \emph{machine motrice}. Le moteur de la moto-taxi ou le groupe électrogène \emph{fournissent} de l'énergie. Une machine simple, elle, ne possède aucune source d'énergie : c'est toujours l'opérateur (ou un moteur) qui fournit l'effort. La machine se contente de le \emph{transformer}.
\end{attention}

\courssub{Les rôles d'une machine simple}

L'expérience du seau d'eau et l'observation des chantiers montrent qu'une machine simple peut jouer trois rôles, souvent combinés.

\begin{aretenir}[Rôles d'une machine simple]
\begin{enumerate}
\item \textbf{Réduire l'effort} : soulever ou déplacer une charge lourde avec une force plus petite que son poids (exemple : la barre à mine pour desceller une pierre).
\item \textbf{Changer la direction de la force} : appliquer l'effort dans la direction la plus commode (exemple : tirer vers le bas la corde du puits pour faire monter le seau).
\item \textbf{Déplacer une charge plus facilement} : adapter le mouvement à la situation (exemple : faire rouler un fût d'huile sur une planche inclinée pour le monter dans le camion).
\end{enumerate}
\end{aretenir}

\begin{exemplebox}[Autour de nous]
\begin{itemize}
\item La \textbf{barre de fer} de la vendeuse de Mokolo, posée sur une pierre, est un \emph{levier} : elle réduit l'effort.
\item La \textbf{planche de chargement} posée sur la ridelle d'un camion est un \emph{plan incliné} : on pousse les sacs ou on fait rouler les fûts avec moins d'effort.
\item La \textbf{roue et la corde} du chantier forment une \emph{poulie} : l'ouvrier tire vers le bas, le seau monte.
\item La \textbf{manivelle} du treuil de forage permet d'enrouler un câble et de remonter de lourds outils avec un effort modéré.
\end{itemize}
\end{exemplebox}

\courssub{Principe fondamental : on ne gagne rien sans rien}

Revenons à l'expérience de la planche inclinée. Avec la rampe, la force $F_2$ est plus petite que le poids du seau : c'est un vrai gain d'effort. Mais observe bien : au lieu de soulever le seau d'un coup sur une hauteur $h$, tu le tires sur toute la \emph{longueur} de la planche, qui est bien plus grande que $h$. Le chemin parcouru a augmenté dans la même proportion que la force a diminué.

C'est un résultat général, que nous admettons :

\begin{propriete}[Principe fondamental des machines simples (admis)]
Une machine simple \textbf{ne crée pas d'énergie}. Le travail fourni par l'opérateur est au moins égal au travail utile reçu par la charge. Par conséquent :
\[ \text{ce qu'on gagne en force, on le perd en distance (et inversement).} \]
Si une machine divise l'effort par 2, la force doit être appliquée sur une distance 2 fois plus grande.
\end{propriete}

\begin{exemplebox}[Vérification sur la planche]
Un sac de \SI{50}{kg} doit être monté à \SI{1}{m} de hauteur dans un camion.
\begin{itemize}
\item En le soulevant directement, on applique une grande force (environ \SI{500}{N}) sur \SI{1}{m}.
\item Avec une planche de \SI{2}{m} de long, la force à fournir n'est plus que d'environ \SI{250}{N} (moitié moins), mais on pousse le sac sur \SI{2}{m} (deux fois plus loin).
\end{itemize}
Le produit $\text{force} \times \text{distance}$ reste le même : $500 \times 1 = 250 \times 2$. La machine n'a rien « supprimé » : elle a réparti l'effort.
\end{exemplebox}

\begin{attention}[Erreur fréquente]
Croire qu'une machine simple \emph{supprime} l'effort. C'est faux : elle le \textbf{réduit} seulement, et toujours au prix d'un déplacement plus long. Dire au BEPC « la machine supprime la force » est une erreur de rédaction ; écris plutôt « la machine réduit l'intensité de la force à appliquer ».
\end{attention}

\courssub{Les quatre machines simples au programme}

Le programme de la classe de 3ème retient quatre machines simples, que nous étudierons une par une dans les sections suivantes.

\begin{definition}[Les quatre machines simples]
\begin{itemize}
\item Le \cle{levier} : une barre rigide qui tourne autour d'un point fixe appelé \emph{point d'appui} (barre à mine, balançoire, pince).
\item Le \cle{plan incliné} : une surface plane inclinée par rapport à l'horizontale (rampe de chargement, route en lacet, vis).
\item La \cle{poulie} : une roue à gorge sur laquelle passe une corde (poulie du puits, palan de chantier).
\item Le \cle{treuil} : un tambour muni d'une manivelle, sur lequel s'enroule une corde ou un câble (treuil de forage, manivelle de store).
\end{itemize}
\end{definition}

\begin{popfigure}
\begin{center}
\begin{tikzpicture}[scale=0.75, >=Stealth]
% Levier
\begin{scope}
\draw[popline, thick] (-1.5,0) -- (1.5,0);
\fill[popgold] (0,0) -- (-0.25,0.35) -- (0.25,0.35) -- cycle;
\draw[popblue, line width=2.5pt] (-1.3,0.75) -- (1.3,0.1);
\draw[->, popgreen, line width=1.8pt] (1.3,0.1) -- (1.3,-0.7);
\fill[poppurpleL] (-1.5,0.55) rectangle (-1.1,0.95);
\node[popdark, font=\small\bfseries] at (0,-1.1) {Le levier};
\end{scope}
% Plan incliné
\begin{scope}[xshift=4.2cm]
\draw[popline, thick] (-1.4,0) -- (1.4,0);
\fill[popblueL] (-1.4,0) -- (1.4,1.1) -- (1.4,0) -- cycle;
\draw[popblue, thick] (-1.4,0) -- (1.4,1.1);
\fill[poppurpleL] (0.1,0.62) rectangle (0.55,0.95);
\draw[->, popgreen, line width=1.8pt] (0.35,0.8) -- (1.15,1.42);
\node[popdark, font=\small\bfseries] at (0,-1.1) {Le plan incliné};
\end{scope}
% Poulie
\begin{scope}[xshift=8.4cm]
\draw[popline, thick] (-1.2,1.9) -- (1.2,1.9);
\draw[popdark, thick] (0,1.9) -- (0,1.5);
\draw[popblue, thick, fill=popblueL] (0,1.1) circle (0.4);
\fill[popdark] (0,1.1) circle (0.05);
\draw[popdark, thick] (-0.4,1.1) -- (-0.4,-0.5);
\draw[popdark, thick] (0.4,1.1) -- (0.4,-0.9);
\fill[poppurpleL] (-0.7,-0.9) rectangle (-0.1,-0.5);
\draw[->, popgreen, line width=1.8pt] (0.4,-0.9) -- (0.4,-1.5);
\node[popdark, font=\small\bfseries] at (0,-1.9) {La poulie};
\end{scope}
% Treuil
\begin{scope}[xshift=12.6cm]
\draw[popline, thick] (-1.2,0) -- (1.2,0);
\draw[popdark, thick] (-0.7,0) -- (-0.7,1.1);
\draw[popdark, thick] (0.7,0) -- (0.7,1.1);
\draw[popblue, thick, fill=popblueL] (-0.55,0.75) rectangle (0.55,1.05);
\draw[popdark, thick] (0.55,0.9) -- (1.0,0.9) -- (1.0,1.3);
\fill[popdark] (1.0,1.35) circle (0.06);
\draw[popdark, thick] (0,0.75) -- (0,-0.2);
\fill[poppurpleL] (-0.3,-0.6) rectangle (0.3,-0.2);
\node[popdark, font=\small\bfseries] at (0,-1.1) {Le treuil};
\end{scope}
\end{tikzpicture}
\end{center}
\legende{Les quatre machines simples au programme de 3ème : le levier (barre et point d'appui), le plan incliné (rampe), la poulie (roue et corde), le treuil (tambour et manivelle).}
\end{popfigure}

\begin{savaistu}[Des machines vieilles comme le monde]
Les pyramides d'Égypte, il y a plus de 4\,500 ans, ont été construites avec des leviers, des plans inclinés en briques crues et des rouleaux de bois pour traîner les blocs de plusieurs tonnes. Au Cameroun, les cases traditionnelles aux lourds toits de paille ont été montées à l'aide de perches utilisées comme leviers et de rampes de terre servant de plans inclinés. Bien avant la théorie, nos ancêtres maîtrisaient déjà la pratique des machines simples !
\end{savaistu}

\begin{exoresolu}[Identifier une machine simple et son rôle]
\textbf{Énoncé.} Sur le port de Douala, un docker fait rouler un fût d'huile de palme le long d'une planche posée entre le sol et la plateforme d'un camion, haute de \SI{1,2}{m}. La planche mesure \SI{3,6}{m}.
\begin{enumerate}
\item Quelle machine simple le docker utilise-t-il ?
\item Quel rôle joue cette machine ?
\item Le fût pèse environ \SI{900}{N}. Le docker pousse avec une force d'environ \SI{300}{N}. Vérifie que le principe fondamental des machines simples est respecté.
\end{enumerate}
\tcblower
\textbf{Solution détaillée.}
\begin{enumerate}
\item La planche posée entre le sol et la plateforme forme une surface plane inclinée : c'est un \textbf{plan incliné}.
\item Cette machine \textbf{réduit l'effort} : au lieu de soulever le fût de \SI{900}{N} à la verticale, le docker le pousse avec une force trois fois plus petite. Elle rend aussi le déplacement plus facile grâce au roulement.
\item Calculons le produit force $\times$ distance dans les deux cas :
\begin{align*}
\text{Sans machine : } & 900 \times \num{1,2} = \SI{1080}{J} \\
\text{Avec la planche : } & 300 \times \num{3,6} = \SI{1080}{J}
\end{align*}
Les deux produits sont égaux : le travail fourni est le même. La force a été divisée par 3, mais la distance a été multipliée par 3. Le principe fondamental est bien respecté : \emph{ce qu'on gagne en force, on le perd en distance}.
\end{enumerate}
\end{exoresolu}

\begin{exercice}[À toi de jouer]
Pour chaque situation, identifie la machine simple utilisée et précise son rôle principal :
\begin{enumerate}
\item Un mécanicien desserre un boulon coincé avec une longue clé.
\item Une maman puise l'eau du puits en tirant la corde qui passe sur une roue fixée au-dessus de l'ouverture.
\item Les motos-taxis gravissent une côte en suivant une route qui monte en zigzag.
\item Un foreur remonte ses outils du fond d'un puits en tournant une manivelle qui enroule un câble.
\end{enumerate}
\end{exercice}

\begin{corrige}[Exercice 1]
\begin{enumerate}
\item La longue clé agit comme un \textbf{levier} (la clé tourne autour de l'axe du boulon) : elle réduit l'effort grâce à la longueur du bras.
\item La roue et la corde forment une \textbf{poulie} : elle change la direction de la force (on tire vers le bas, le seau monte).
\item La route en lacet est un \textbf{plan incliné} : elle réduit l'effort nécessaire pour monter, au prix d'un trajet plus long.
\item La manivelle et le tambour forment un \textbf{treuil} : il réduit l'effort pour remonter la charge.
\end{enumerate}
\end{corrige}

\begin{aretenir}[L'essentiel de la section]
\begin{itemize}
\item Une \cle{machine simple} est un dispositif mécanique qui permet d'accomplir un travail en modifiant l'intensité ou la direction de la force à appliquer.
\item Ses rôles : \textbf{réduire l'effort}, \textbf{changer la direction de la force}, \textbf{déplacer une charge plus facilement}.
\item Principe fondamental (admis) : une machine simple \textbf{ne crée pas d'énergie} ; ce qu'on gagne en force, on le perd en distance.
\item Les quatre machines simples au programme : le \textbf{levier}, le \textbf{plan incliné}, la \textbf{poulie} et le \textbf{treuil}.
\end{itemize}
\end{aretenir}

\coursec{Notion de machine simple}

\textbf{Situation-problème.} Au marché de Mfoundi, une marchande doit charger des sacs de riz de \SI{50}{kg} sur une benne de camion haute de \SI{1,2}{m}. Elle ne peut pas les soulever verticalement. Elle utilise alors une longue planche en bois qu'elle pose sur le bord de la benne : elle pousse le sac le long de la planche. L'effort lui semble bien plus faible. Pourquoi cette planche inclinée facilite-t-elle le travail ?

\begin{definition}[Machine simple]
Une \cle{machine simple} est un \textbf{dispositif mécanique élémentaire} qui permet d'accomplir un travail en modifiant \textbf{soit l'intensité, soit la direction d'une force} (l'effort) que l'on doit fournir pour vaincre une \textbf{résistance} (la charge).
\end{definition}

Toute machine simple met en jeu quatre grandeurs fondamentales :

\begin{center}
\begin{tabularx}{\linewidth}{|l|X|X|}
\hline
\rowcolor{popblueL} \textbf{Grandeur} & \textbf{Notation} & \textbf{Rôle} \\
\hline
Effort (force motrice) & $F$ en \SI{}{N} & Force fournie par l'utilisateur \\
\hline
Déplacement de l'effort & $d_F$ en \SI{}{m} & Distance parcourue par le point d'application de $F$ \\
\hline
Résistance (charge) & $R$ en \SI{}{N} & Force à vaincre (poids, frottement, etc.) \\
\hline
Déplacement de la résistance & $d_R$ en \SI{}{m} & Distance parcourue par le point d'application de $R$ \\
\hline
\end{tabularx}
\end{center}

\begin{propriete}[Règle d'or des machines simples (idéal)]
Dans une machine simple \textbf{idéale} (sans frottements, sans masse propre), le \textbf{travail moteur} fourni par l'utilisateur est égal au \textbf{travail résistant} reçu par la charge :
\[ W_{\text{moteur}} = W_{\text{résistant}} \quad \Longrightarrow \quad F \times d_F = R \times d_R \]
\emph{Conséquence :} si la machine permet de diminuer l'effort ($F < R$), elle impose en contrepartie un déplacement plus grand pour l'effort ($d_F > d_R$). \textbf{« Ce qu'on gagne en force, on le perd en déplacement. »}
\end{propriete}

\begin{attention}[Rendement réel]
Dans la réalité, les frottements (dans les axes, sur les surfaces) et la masse de la machine elle-même « mangent » une partie du travail fourni. Le \cle{rendement} $\eta$ (eta) est toujours inférieur à 1 (ou \SI{100}{\%}) :
\[ \eta = \frac{W_{\text{résistant}}}{W_{\text{moteur}}} = \frac{R \times d_R}{F \times d_F} < 1 \]
Au BEPC, on raisonne souvent sur la machine \emph{idéale} ($\eta = 1$) sauf si l'énoncé donne un rendement.
\end{attention}

Les \textbf{quatre machines simples} au programme de 3ème sont :
\begin{enumerate}
\item Le \cle{levier} (barre rigide pivotante) ;
\item Le \cle{plan incliné} (surface plane oblique) ;
\item La \cle{poulie} (roue à gorge guidant un câble) ;
\item Le \cle{treuil} (cylindre tournant autour d'un axe).
\end{enumerate}
Nous les étudierons l'une après l'autre, en commençant par la plus ancienne : le levier.

\coursec{Le levier}

Dans la section précédente, nous avons découvert ce qu'est une \cle{machine simple} : un dispositif qui permet d'accomplir un travail en modifiant l'intensité ou la direction d'une force. Nous allons maintenant étudier la plus ancienne et la plus répandue de toutes les machines simples : le \cle{levier}. Tu l'utilises chaque jour sans le savoir : quand tu ouvres une bouteille avec un décapsuleur, quand tu joues à la balançoire, quand le menuisier arrache un clou avec le pied-de-biche, ou quand le maçon pousse sa brouette sur le chantier.

\courssub{Définition et éléments d'un levier}

\textbf{Situation-problème.} Sur un chantier de construction à Douala, un ouvrier doit déplacer une grosse pierre trop lourde pour être soulevée à mains nues. Il glisse une longue barre de fer sous la pierre, cale une pierre plus petite sous la barre, puis appuie fortement sur l'extrémité libre. La grosse pierre se soulève ! Comment une simple barre permet-elle de multiplier la force de l'ouvrier ?

\begin{definition}[Le levier]
Un \cle{levier} est une \textbf{barre rigide} qui peut tourner autour d'un point fixe appelé \cle{point d'appui} (ou \emph{pivot}). Il est soumis à deux forces :
\begin{itemize}
\item l'\cle{effort} $\vec{F}$ (ou force motrice) : la force que l'on exerce sur la barre ;
\item la \cle{résistance} $\vec{R}$ (ou charge) : la force à vaincre, exercée par l'objet que l'on veut déplacer.
\end{itemize}
\end{definition}

Tout levier possède donc \textbf{trois éléments essentiels} :

\begin{center}
\begin{tabularx}{\linewidth}{|l|X|X|}
\hline
\rowcolor{popblueL} \textbf{Élément} & \textbf{Rôle} & \textbf{Exemple (barre à mine)} \\
\hline
Point d'appui (pivot) & Point fixe autour duquel la barre tourne & La petite pierre calée sous la barre \\
\hline
Effort $\vec{F}$ & Force motrice appliquée par l'utilisateur & La poussée de l'ouvrier vers le bas \\
\hline
Résistance $\vec{R}$ & Charge à vaincre & Le poids de la grosse pierre \\
\hline
\end{tabularx}
\end{center}

On appelle \cle{bras de levier} d'une force la distance entre le point d'appui et le point d'application de cette force. On note :
\begin{itemize}
\item $d_1$ : le bras de levier de l'effort $\vec{F}$ ;
\item $d_2$ : le bras de levier de la résistance $\vec{R}$.
\end{itemize}

\begin{popfigure}
\begin{center}
\begin{tikzpicture}[scale=1,>=Stealth]
% Point d'appui triangulaire
\fill[poppurple!70] (0,-0.55) -- (-0.5,-1.3) -- (0.5,-1.3) -- cycle;
\draw[popdark,thick] (0,-0.55) -- (-0.5,-1.3) -- (0.5,-1.3) -- cycle;
\fill[pattern=north east lines,pattern color=popmuted] (-0.9,-1.3) rectangle (0.9,-1.5);
\draw[popdark] (-0.9,-1.3) -- (0.9,-1.3);
\node[below,popdark] at (0,-1.55) {\small Point d'appui (pivot)};
% Barre
\draw[line width=3pt,popblue] (-4.2,-0.35) -- (3.2,-0.35);
\fill[popdark] (0,-0.35) circle (2pt);
% Résistance (charge) à gauche
\draw[->,very thick,poporange] (-3.5,0.55) -- (-3.5,-0.3);
\node[above,poporange] at (-3.5,0.55) {\small $\vec{R}$ (résistance)};
\fill[poporange!40,draw=poporange] (-3.9,-0.35) rectangle (-3.1,0.15);
% Effort à droite
\draw[->,very thick,popgreen] (2.5,0.9) -- (2.5,-0.3);
\node[above,popgreen] at (2.5,0.9) {\small $\vec{F}$ (effort)};
% Distances
\draw[<->,popdark] (0,-0.75) -- (-3.5,-0.75) node[midway,below]{\small $d_2$};
\draw[<->,popdark] (0,-0.95) -- (2.5,-0.95) node[midway,below]{\small $d_1$};
\end{tikzpicture}
\end{center}
\legende{Schéma d'un levier en équilibre : la barre rigide (bleue) tourne autour du point d'appui triangulaire. La résistance $\vec{R}$ s'exerce à la distance $d_2$ du pivot, l'effort $\vec{F}$ à la distance $d_1$.}
\end{popfigure}

\courssub{Expérience : équilibre d'un levier}

Comment l'effort nécessaire dépend-il de la position du point d'appui ? Répondons-y par la démarche d'investigation.

\begin{experience}[Équilibrer une règle sur une gomme]
\textbf{Matériel :} une règle graduée rigide (en bois ou plastique), une gomme (ou un taille-crayon) servant de point d'appui, des masses marquées (\SI{50}{g}, \SI{100}{g}, \SI{200}{g}), du fil pour accrocher les masses.

\textbf{Protocole :}
\begin{enumerate}
\item Placer la gomme sous la règle, par exemple à la graduation \SI{10}{cm}.
\item Accrocher une masse $m_2 = \SI{200}{g}$ (la résistance) à la graduation \SI{5}{cm}, soit $d_2 = \SI{5}{cm}$ du point d'appui.
\item Chercher où placer une masse $m_1 = \SI{100}{g}$ (l'effort) de l'autre côté pour que la règle soit horizontale (équilibre). Noter la graduation et en déduire $d_1$.
\item Recommencer en changeant les masses et les positions du point d'appui. Compléter le tableau ci-dessous.
\end{enumerate}

\textbf{Résultats obtenus (exemple) :}
\begin{center}
\begin{tabular}{|c|c|c|c|c|}
\hline
\rowcolor{popblueL} Expérience & $F$ (N) & $d_1$ (cm) & $R$ (N) & $d_2$ (cm) \\
\hline
1 & 1,0 & 10 & 2,0 & 5 \\
\hline
2 & 1,0 & 20 & 2,0 & 10 \\
\hline
3 & 0,5 & 20 & 2,0 & 5 \\
\hline
4 & 2,0 & 10 & 1,0 & 20 \\
\hline
\end{tabular}
\end{center}
\emph{Rappel :} on utilise $g \approx \SI{10}{N/kg}$, donc \SI{100}{g} $\rightarrow$ \SI{1}{N}, \SI{200}{g} $\rightarrow$ \SI{2}{N}, \SI{50}{g} $\rightarrow$ \SI{0,5}{N}.

\textbf{Observations :} à chaque équilibre, le produit $F \times d_1$ est égal au produit $R \times d_2$ (par exemple, expérience 1 : $1,0 \times 10 = 2,0 \times 5 = 10$).

\textbf{Interprétation :} quand le bras de levier $d_1$ de l'effort augmente, l'effort $F$ nécessaire à l'équilibre diminue. C'est pourquoi l'ouvrier appuie le plus loin possible du point d'appui.

\textbf{Conclusion :} un levier est en équilibre lorsque les produits de chaque force par son bras de levier sont égaux.
\end{experience}

\begin{propriete}[Condition d'équilibre d'un levier]
Un levier est en \textbf{équilibre} lorsque le produit de l'effort par son bras de levier est égal au produit de la résistance par son bras de levier :
\[ F \times d_1 = R \times d_2 \]
où $F$ et $R$ sont en newtons (N) et $d_1$, $d_2$ dans la \textbf{même unité} de longueur (m ou cm).

\emph{Conséquence pratique :} plus le bras de levier $d_1$ de l'effort est long par rapport à $d_2$, plus l'effort $F$ nécessaire est faible. Le levier permet donc de \textbf{démultiplier une force}. Cette propriété est établie expérimentalement ; savoir l'utiliser dans un calcul est exigible au BEPC.
\end{propriete}

\begin{methode}[Calculer l'effort nécessaire sur un levier]
\begin{enumerate}
\item \textbf{Identifier} le point d'appui, l'effort $F$ et la résistance $R$ sur le schéma ou l'énoncé.
\item \textbf{Mesurer} (ou relever) les bras de levier $d_1$ et $d_2$, distances de chaque force au point d'appui, dans la \emph{même unité}.
\item \textbf{Écrire} la condition d'équilibre : $F \times d_1 = R \times d_2$.
\item \textbf{Isoler} l'inconnue : $F = \dfrac{R \times d_2}{d_1}$, puis calculer et conclure avec l'unité (newton).
\end{enumerate}
\end{methode}

\begin{exemplebox}[Soulever une charge avec une barre]
Un ouvrier veut soulever une pierre dont le poids est $R = \SI{200}{N}$. La pierre est à $d_2 = \SI{0,5}{m}$ du point d'appui et l'ouvrier appuie à $d_1 = \SI{2}{m}$ du point d'appui. L'effort nécessaire est :
\[ F = \frac{R \times d_2}{d_1} = \frac{200 \times 0,5}{2} = \SI{50}{N}. \]
L'effort est \textbf{quatre fois plus faible} que la charge : le levier a démultiplié la force de l'ouvrier par 4.
\end{exemplebox}

\begin{attention}[Erreur fréquente : distance au point d'appui ou longueur de la barre ?]
Les distances $d_1$ et $d_2$ sont les distances entre chaque force et le \textbf{point d'appui}, et \emph{non} la longueur totale de la barre ! Si le point d'appui n'est pas au milieu, $d_1 + d_2$ est la longueur de la barre, mais chaque distance se mesure séparément \emph{à partir du pivot}. Vérifie aussi que $d_1$ et $d_2$ sont exprimées dans la \textbf{même unité} avant de calculer.
\end{attention}

\courssub{Les trois types de leviers}

Selon la \textbf{position relative} du point d'appui, de l'effort et de la résistance, on distingue trois types de leviers.

\begin{definition}[Types de leviers]
\begin{itemize}
\item \cle{Levier inter-appui} : le \textbf{point d'appui} est \emph{entre} l'effort et la résistance. Exemples : la balançoire, les ciseaux, la barre à mine, le pied-de-biche.
\item \cle{Levier inter-résistant} : la \textbf{résistance} est \emph{entre} le point d'appui et l'effort. Exemples : la brouette, le casse-noix, l'ouvre-bouteille (décapsuleur).
\item \cle{Levier inter-effort} : l'\textbf{effort} est \emph{entre} le point d'appui et la résistance. Exemples : la pince à épiler, la pédale de frein, la pagaie de pirogue.
\end{itemize}
\end{definition}

\begin{popfigure}
\begin{center}
\begin{tikzpicture}[scale=0.85,>=Stealth]
% --- Levier inter-appui ---
\begin{scope}[shift={(0,0)}]
\fill[poppurple!70] (0,-0.3) -- (-0.35,-0.85) -- (0.35,-0.85) -- cycle;
\draw[popdark] (0,-0.3) -- (-0.35,-0.85) -- (0.35,-0.85) -- cycle;
\draw[line width=2.5pt,popblue] (-2.2,-0.15) -- (2.2,-0.15);
\draw[->,very thick,poporange] (-1.8,0.45) -- (-1.8,-0.1);
\node[above,poporange] at (-1.8,0.45) {\small $R$};
\draw[->,very thick,popgreen] (1.8,0.45) -- (1.8,-0.1);
\node[above,popgreen] at (1.8,0.45) {\small $F$};
\node[below,popdark] at (0,-0.9) {\small appui};
\node[below,align=center,popdark] at (0,-1.4) {\small \textbf{Inter-appui}\\ \scriptsize balançoire, ciseaux};
\end{scope}
% --- Levier inter-résistant ---
\begin{scope}[shift={(5.0,0)}]
\fill[poppurple!70] (-1.8,-0.3) -- (-2.15,-0.85) -- (-1.45,-0.85) -- cycle;
\draw[popdark] (-1.8,-0.3) -- (-2.15,-0.85) -- (-1.45,-0.85) -- cycle;
\draw[line width=2.5pt,popblue] (-2.2,-0.15) -- (2.2,-0.15);
\draw[->,very thick,poporange] (-0.4,0.45) -- (-0.4,-0.1);
\node[above,poporange] at (-0.4,0.45) {\small $R$};
\draw[->,very thick,popgreen] (1.8,0.45) -- (1.8,-0.1);
\node[above,popgreen] at (1.8,0.45) {\small $F$};
\end{scope}
\end{tikzpicture}
\legende{Leviers inter-appui (balançoire, ciseaux) et inter-résistant : $R$ résistance, $F$ effort.}
\end{center}
\end{popfigure}


\coursec{Le plan incliné}

Dans la section précédente, nous avons découvert le \cle{levier} : une barre rigide qui tourne autour d'un point d'appui et qui permet de soulever de lourdes charges avec un effort réduit. Mais le levier a une limite : il ne peut soulever une charge que d'une faible hauteur, car le bras ne tourne que d'un petit angle. Comment faire alors pour monter une charge lourde de plusieurs mètres, par exemple pour charger des sacs de sable sur la benne d'un camion ? Observe les manœuvres au bord de la route : ils posent une longue planche entre le sol et la benne, puis ils font rouler ou glisser les sacs le long de la planche. Cette planche, c'est une deuxième machine simple : le \cle{plan incliné}.

\courssub{Définition du plan incliné}

\begin{definition}[Le plan incliné]
Un \cle{plan incliné} est une surface plane qui forme un \cle{angle} $\alpha$ avec l'horizontale et qui est utilisée pour \textbf{élever} ou \textbf{descendre} des charges avec un effort réduit. Plus l'angle $\alpha$ est petit, plus la pente est dite \emph{douce} ; plus l'angle est grand, plus la pente est dite \emph{forte}.
\end{definition}

L'idée intuitive est simple : au lieu de soulever une charge \emph{verticalement} (il faut alors vaincre tout son poids d'un coup), on la fait monter \emph{progressivement} le long d'une surface oblique. L'effort est étalé sur une plus grande distance, donc il est plus faible à chaque instant. C'est exactement ce que tu ressens quand tu montes une colline par un sentier en pente douce plutôt que par un raccourci très raide : le raccourci est plus court, mais il fatigue beaucoup plus.

\begin{popfigure}
\begin{center}
\begin{tikzpicture}[scale=1.0]
  % Sol horizontal
  \draw[popline, thick] (-0.5,0) -- (8.5,0);
  % Plan incliné (triangle rectangle)
  \fill[popblueL] (0,0) -- (7,0) -- (7,3) -- cycle;
  \draw[popblue, very thick] (0,0) -- (7,3);
  \draw[popdark, thick] (0,0) -- (7,0) -- (7,3);
  % Hachures du sol
  \foreach \x in {-0.3,0.3,...,8.3} \draw[popmuted] (\x,0) -- (\x-0.25,-0.25);
  % Angle alpha
  \draw[poporange, thick] (1.6,0) arc (0:23.2:1.6);
  \node[poporange] at (2.15,0.32) {$\alpha$};
  % Hauteur h
  \draw[<->, popdark] (7.35,0) -- (7.35,3) node[midway, right] {$h$};
  % Longueur L (base horizontale)
  \draw[<->, popdark] (0.15,-0.45) -- (6.85,-0.45) node[midway, below] {longueur du trajet horizontal};
  % Caisse sur la pente
  \begin{scope}[rotate around={23.2:(3.5,1.5)}]
    \fill[popgoldL] (3.0,1.5) rectangle (4.0,2.3);
    \draw[popdark, thick] (3.0,1.5) rectangle (4.0,2.3);
    \node[popdark] at (3.5,1.9) {\small charge};
  \end{scope}
  % Force de traction F le long de la pente
  \draw[->, popgreen, very thick] (3.85,2.05) -- ++(23.2:2.2) node[above right, popgreen] {$\vec{F}$ (traction)};
  % Poids P vertical
  \draw[->, poppink, very thick] (3.5,1.9) -- (3.5,0.35) node[below, poppink] {$\vec{P}$ (poids)};
\end{tikzpicture}
\end{center}
\legende{Schéma du plan incliné : la charge monte le long de la pente sous l'action de la force de traction $\vec{F}$, parallèle au plan. Son poids $\vec{P}$ reste vertical, dirigé vers le bas. L'angle $\alpha$ caractérise l'inclinaison et $h$ est la hauteur à atteindre.}
\end{popfigure}

\begin{aretenir}[Les éléments d'un plan incliné]
\begin{itemize}
  \item La \textbf{surface} (la planche, la rampe, la route) : c'est le support sur lequel la charge glisse ou roule.
  \item L'\textbf{angle d'inclinaison} $\alpha$ : angle entre le plan et l'horizontale.
  \item La \textbf{hauteur} $h$ : dénivelé entre le bas et le haut du plan.
  \item La \textbf{longueur} $L$ du plan : distance réellement parcourue par la charge. On a toujours $L > h$.
\end{itemize}
\end{aretenir}

\courssub{Expérience de classe : comparer les forces}

\begin{experience}[Tirer une charge verticalement puis sur un plan incliné]
\textbf{Matériel} : une masse marquée (ou un petit sac rempli de sable), un dynamomètre, une planche lisse, des cales pour incliner la planche, un rapporteur, une règle.

\textbf{Protocole} :
\begin{enumerate}
  \item Accrocher la masse au dynamomètre et la soulever \textbf{verticalement}, lentement. Noter la force indiquée : c'est le poids $P$ de l'objet.
  \item Poser la planche en pente douce (angle $\alpha$ faible) et tirer la masse \textbf{le long de la planche}, à vitesse constante, avec le dynamomètre parallèle à la pente. Noter la force $F$.
  \item Recommencer en augmentant progressivement l'inclinaison (pente moyenne, puis forte). Noter $F$ à chaque fois.
\end{enumerate}

\textbf{Observations} (résultats typiques pour une masse de poids $P = \SI{10}{N}$, frottements faibles) :
\begin{center}
\begin{tabularx}{\linewidth}{|l|X|X|X|X|}
\hline
\rowcolor{popblueL} Situation & Vertical ($\alpha = 90^\circ$) & Pente douce ($\alpha \approx 15^\circ$) & Pente moyenne ($\alpha \approx 30^\circ$) & Pente forte ($\alpha \approx 60^\circ$) \\
\hline
Force mesurée & \SI{10,0}{N} & environ \SI{2,6}{N} & environ \SI{5,0}{N} & environ \SI{8,7}{N} \\
\hline
\end{tabularx}
\end{center}

\textbf{Interprétation} : sur le plan incliné, la force à fournir est \textbf{toujours inférieure au poids} de l'objet. De plus, plus la pente est douce, plus la force est faible. Quand le plan devient vertical ($\alpha = 90^\circ$), on retrouve $F = P$ : le plan incliné n'apporte alors plus aucun avantage.

\textbf{Conclusion} : le plan incliné permet d'élever une charge avec un effort \textbf{plus faible} que son poids, mais la charge parcourt une distance \textbf{plus grande} que la hauteur gagnée.
\end{experience}

\begin{popfigure}
\begin{center}
\begin{tikzpicture}
\begin{axis}[
  width=0.85\linewidth, height=6.5cm,
  xlabel={Angle d'inclinaison $\alpha$ (degrés)},
  ylabel={Force $F$ (N)},
  xmin=0, xmax=90, ymin=0, ymax=12,
  xtick={0,15,30,45,60,75,90},
  ytick={0,2,4,6,8,10,12},
  grid=both, grid style={popline!40},
  axis line style={popdark},
  label style={popdark},
]
\addplot[popcurve, domain=0:90, samples=100] {10*sin(x)};
\addplot[only marks, mark=*, mark size=2pt, poporange] coordinates {(0,0) (15,2.6) (30,5) (45,7.1) (60,8.7) (75,9.7) (90,10)};
\draw[dashed, popmuted] (axis cs:30,0) -- (axis cs:30,5) -- (axis cs:0,5);
\node[popdark, anchor=south west, font=\small] at (axis cs:32,5.3) {pente moyenne : $F \approx \SI{5}{N}$};
\end{axis}
\end{tikzpicture}
\end{center}
\legende{Force nécessaire $F$ en fonction de l'angle d'inclinaison $\alpha$, pour une charge de poids $P = \SI{10}{N}$ (frottements négligés). La courbe est croissante : plus le plan est incliné, plus l'effort est grand. Pour $\alpha = 90^\circ$ (montée verticale), on retrouve $F = P = \SI{10}{N}$.}
\end{popfigure}

\courssub{Interprétation : pourquoi l'effort diminue-t-il ?}

Quand on soulève une charge verticalement, la force à exercer doit compenser \textbf{la totalité} du poids $\vec{P}$. Sur un plan incliné, le plan \textbf{porte une partie du poids} : la charge appuie sur la surface, qui la soutient. Il ne reste plus à l'opérateur qu'à vaincre la partie du poids qui « tire » la charge vers le bas \emph{le long de la pente}. Or cette partie est d'autant plus petite que la pente est douce.

\begin{propriete}[Effort sur un plan incliné — établie expérimentalement]
Sur un plan incliné, l'effort $F$ nécessaire pour monter une charge à vitesse constante est :
\begin{itemize}
  \item \textbf{toujours inférieur au poids} de la charge : $F < P$ (si les frottements sont négligés) ;
  \item \textbf{d'autant plus faible que la pente est douce} : quand l'angle $\alpha$ diminue, $F$ diminue ; quand $\alpha$ augmente, $F$ augmente.
\end{itemize}
En contrepartie, la distance parcourue $L$ est d'autant plus grande que la pente est douce. \emph{On gagne en force ce que l'on perd en distance.}
\end{propriete}

\begin{exemplebox}[Mesures au dynamomètre]
Un objet de poids $P = \SI{10}{N}$ est tiré au dynamomètre :
\begin{itemize}
  \item \textbf{verticalement} : le dynamomètre indique $\SI{10}{N}$ ;
  \item \textbf{sur un plan incliné doux} ($\alpha \approx 30^\circ$) : le dynamomètre indique environ $\SI{5}{N}$.
\end{itemize}
L'effort a été divisé par deux, mais pour s'élever de $h = \SI{1}{m}$, l'objet parcourt $L = \SI{2}{m}$ le long de la pente : la distance a été multipliée par deux.
\end{exemplebox}

\begin{attention}[Erreur fréquente : « le plan incliné rend l'objet plus léger »]
Non ! Le \textbf{poids} $\vec{P}$ de l'objet ne change \textbf{jamais} : il dépend de la masse de l'objet et de l'attraction de la Terre, pas du chemin suivi. Ce qui diminue, c'est \textbf{l'effort $F$ que l'opérateur doit fournir}. Au BEPC, écris toujours : « le plan incliné réduit \emph{l'effort} (la force à exercer) », et jamais « le plan incliné réduit le poids ».
\end{attention}

\courssub{Avantages et inconvénients du plan incliné}

\begin{center}
\begin{tabularx}{\linewidth}{|X|X|}
\hline
\rowcolor{popgreenL} \textbf{Avantages} & \textbf{Inconvénients} \\
\hline
Réduction de l'effort : $F < P$ & Distance parcourue plus grande ($L > h$) \\
\hline
Permet d'élever de lourdes charges sans moteur ni machine complexe & Les frottements entre la charge et le plan augmentent l'effort réel \\
\hline
Dispositif simple, bon marché, facile à fabriquer (une planche suffit) & La montée est plus lente : le déplacement prend plus de temps \\
\hline
Utilisable partout, même sans électricité & Il faut de la place pour installer un plan long \\
\hline
\end{tabularx}
\end{center}

\courssub{Le plan incliné dans la vie quotidienne}

\begin{itemize}
  \item \textbf{La planche de chargement} : pour charger des sacs de sable, de ciment ou des bidons d'huile sur un camion, les manœuvres glissent les charges le long d'une planche inclinée au lieu de les soulever à bout de bras.
  \item \textbf{La rampe d'accès pour personnes à mobilité réduite} : à l'entrée des écoles, des hôpitaux et des administrations, une rampe en pente douce permet à un fauteuil roulant de monter sans escalier. Plus la rampe est longue et douce, moins l'effort est grand.
  \item \textbf{Les routes en lacets en montagne} : dans les monts Bamboutos ou sur les routes de l'Ouest, la route ne monte jamais en ligne droite vers le sommet : elle serpente en zigzag. Chaque lacet est un plan incliné à pente douce.
  \item \textbf{Les escaliers} : un escalier est un plan incliné « en marches » : on monte progressivement, avec moins d'effort qu'en grimpant à la verticale.
\end{itemize}

\begin{savaistu}[Les routes de montagne en zigzag sont des plans inclinés]
Dans les régions montagneuses du Cameroun, comme les monts Bamboutos, une route qui monterait tout droit serait si raide qu'aucune moto-taxi ni aucun camion chargé ne pourrait la gravir. Les ingénieurs tracent donc la route en \textbf{lacets} : le trajet devient beaucoup plus long, mais la pente devient douce et la montée possible. C'est le prix à payer du plan incliné : \emph{on échange de la distance contre de la force}. De même, la \textbf{vis} est un plan incliné \emph{enroulé} autour d'un cylindre (c'est pour cela qu'elle s'enfonce facilement en tournant), et le \textbf{coin} (la lame de la machette, le clou) est un \emph{double} plan incliné qui fend le bois. Ce sont des compléments utiles à connaître, mais non exigibles au BEPC.
\end{savaistu}

\begin{exoresolu}[Chargement d'un camion de sable]
\textbf{Énoncé.} Des manœuvres doivent charger un sac de poids $P = \SI{600}{N}$ sur une benne située à $h = \SI{1,5}{m}$ du sol. Ils disposent d'une planche de longueur $L = \SI{4,5}{m}$.
\begin{enumerate}
  \item Quelle force faudrait-il pour soulever le sac verticalement ?
  \item Les mesures au dynamomètre donnent $F = \SI{200}{N}$ sur la planche. Comparer $F$ et $P$. Conclure.
  \item Comparer la distance parcourue sur la planche à la hauteur gagnée. Conclure.
  \item Que se passerait-il si l'on utilisait une planche plus longue, de $\SI{9}{m}$, pour la même hauteur ?
\end{enumerate}
\tcblower
\textbf{Solution détaillée.}
\begin{enumerate}
  \item Pour soulever le sac verticalement, il faut exercer une force au moins égale à son poids : $F_{\text{vertical}} = P = \SI{600}{N}$.
  \item Sur la planche, $F = \SI{200}{N} < P = \SI{600}{N}$. L'effort est trois fois plus faible : le plan incliné a bien \textbf{réduit l'effort} (le poids du sac, lui, reste $\SI{600}{N}$).
  \item Sur la planche, le sac parcourt $L = \SI{4,5}{m}$ pour ne s'élever que de $h = \SI{1,5}{m}$ : la distance est trois fois plus grande. \textbf{Conclusion} : on a gagné en force (effort divisé par 3) ce que l'on a perdu en distance (trajet multiplié par 3).
  \item Avec une planche de $\SI{9}{m}$ pour la même hauteur, la pente serait plus douce (angle $\alpha$ plus petit). L'effort serait donc \textbf{encore plus faible} (environ $\SI{100}{N}$), mais le trajet serait plus long et le chargement prendrait plus de temps.
\end{enumerate}
\end{exoresolu}

\begin{exercice}[Rampe d'accès d'un hôpital]
Pour franchir un dénivelé de $h = \SI{0,8}{m}$, on hésite entre deux rampes : une rampe courte de $\SI{2}{m}$ et une rampe longue de $\SI{6}{m}$.
\begin{enumerate}
  \item Laquelle des deux rampes a la pente la plus douce ? Justifier.
  \item Laquelle demandera le moins d'effort pour monter un fauteuil roulant ?
  \item Citer un inconvénient de la rampe la plus longue.
\end{enumerate}
\end{exercice}

\begin{corrige}[Exercice : rampe d'accès d'un hôpital]
\begin{enumerate}
  \item Pour la même hauteur $h = \SI{0,8}{m}$, la rampe de $\SI{6}{m}$ est plus longue, donc son angle $\alpha$ avec l'horizontale est plus petit : c'est elle qui a la pente la plus douce.
  \item C'est aussi elle qui demande le moins d'effort, car l'effort sur un plan incliné est d'autant plus faible que la pente est douce ($F < P$ dans les deux cas, mais $F$ est plus petit pour la rampe longue).
  \item Inconvénients possibles : elle occupe plus de place, le trajet est plus long, la montée prend plus de temps.
\end{enumerate}
\end{corrige}

\begin{aretenir}[L'essentiel sur le plan incliné]
\begin{itemize}
  \item Le plan incliné est une \textbf{surface plane inclinée} d'un angle $\alpha$ par rapport à l'horizontale, utilisée pour élever ou descendre des charges.
  \item L'effort à fournir est \textbf{plus faible que le poids} : $F < P$ (frottements négligés).
  \item Plus la pente est \textbf{douce} (plan long, angle faible), plus l'effort est \textbf{faible}, mais plus la distance parcourue est \textbf{grande}.
  \item Le \textbf{poids} de la charge ne change pas : c'est l'effort qui diminue.
  \item Exemples : planche de chargement, rampe d'accès, routes en lacets, escaliers ; la vis et le coin sont des plans inclinés déguisés.
\end{itemize}
\end{aretenir}

Le plan incliné réduit donc l'effort, mais il impose de déplacer la charge sur toute la longueur de la pente. Existe-t-il une machine simple qui permette de soulever une charge verticalement, sans la déplacer sur une longue distance, tout en réduisant l'effort ? Oui : c'est la \cle{poulie}, que nous étudions dans la section suivante.

\coursec{La poulie}

Dans la section précédente, nous avons vu que le \cle{plan incliné} permet de soulever une charge avec un effort réduit, au prix d'un trajet plus long. Mais sur un chantier de construction à Douala ou à Yaoundé, on ne peut pas toujours installer une rampe : les sacs de ciment doivent monter verticalement, parfois jusqu'au troisième étage. Comment faire ? Les maçons utilisent alors un objet très simple : une petite roue à gorge fixée en haut du mur, sur laquelle passe une corde. Cette roue, c'est la \cle{poulie}, héroïne discrète des puits, des chantiers et des garages. Dans cette section, nous allons définir la poulie, distinguer ses deux grands modes d'utilisation (fixe et mobile), établir expérimentalement les lois qui les gouvernent, puis comparer leurs avantages et leurs inconvénients.

\courssub{Définition et description de la poulie}

\begin{definition}[La poulie]
Une \cle{poulie} est une \textbf{roue à gorge} qui tourne librement autour d'un \textbf{axe}. Une \textbf{corde} (ou une courroie, une chaîne) passe dans la gorge de la roue. Elle sert à \textbf{transmettre une force} et à soulever ou déplacer une charge.
\end{definition}

Une poulie comporte donc trois éléments essentiels :
\begin{itemize}
\item la \textbf{roue à gorge} : la gorge (rainure circulaire) empêche la corde de glisser sur le côté ;
\item l'\textbf{axe} : c'est autour de lui que la roue tourne ; c'est lui qui peut être fixe ou mobile ;
\item la \textbf{corde} : elle relie l'opérateur à la charge en passant dans la gorge.
\end{itemize}

\begin{savaistu}[La poulie, une invention millénaire]
On attribue à Archimède (III\textsuperscript{e} siècle avant J.-C.) l'étude des systèmes de poulies. Aujourd'hui encore, au Cameroun, la poulie est partout : au-dessus des puits villageois pour remonter les seaux d'eau, sur les chantiers de construction pour hisser le sable et le ciment, dans les garages des mécaniciens de quartier pour sortir un moteur de voiture grâce à un palan, et même sur les grues du port de Douala.
\end{savaistu}

\courssub{La poulie fixe}

\begin{definition}[Poulie fixe]
Une poulie est dite \cle{fixe} lorsque son \textbf{axe reste immobile} pendant l'utilisation : la poulie est accrochée à un support (poutre, branche, crochet) et seule la roue tourne.
\end{definition}

Tirons sur la corde d'une poulie fixe : la charge monte pendant que notre main descend. La poulie fixe \textbf{change la direction de la force} : au lieu de tirer la charge vers le haut (position fatigante), on tire la corde \textbf{vers le bas}, en profitant du poids de son propre corps. C'est beaucoup plus confortable. Mais l'effort est-il réduit ? C'est ce que l'expérience va nous montrer.

\begin{experience}[Comparaison des forces avec et sans poulies]
\textbf{Matériel :} une charge (sac ou poids) de résistance $R = \SI{10}{N}$, un dynamomètre, une poulie fixe, une poulie mobile, une corde, un support, une règle graduée.

\textbf{Protocole :}
\begin{enumerate}
\item \textbf{Montage 1 (direct) :} accrocher la charge au dynamomètre et la soulever lentement. Noter la force $F_1$.
\item \textbf{Montage 2 (poulie fixe) :} fixer la poulie au support, passer la corde sur la gorge, accrocher la charge à une extrémité et le dynamomètre à l'autre. Tirer vers le bas. Noter $F_2$.
\item \textbf{Montage 3 (poulie mobile) :} fixer une extrémité de la corde au support, passer la corde sous la poulie qui porte la charge, et tirer l'autre extrémité vers le haut avec le dynamomètre. Noter $F_3$. Mesurer aussi la longueur de corde tirée et la hauteur de montée de la charge.
\end{enumerate}

\textbf{Observations :}
\begin{center}
\begin{tabularx}{\linewidth}{|l|X|X|}
\hline
\rowcolor{popblueL} \textbf{Montage} & \textbf{Force mesurée} & \textbf{Corde tirée pour monter de \SI{20}{cm}} \\
\hline
Direct & $F_1 = \SI{10}{N}$ & \SI{20}{cm} \\
\hline
Poulie fixe & $F_2 = \SI{10}{N}$ & \SI{20}{cm} \\
\hline
Poulie mobile & $F_3 = \SI{5}{N}$ & \SI{40}{cm} \\
\hline
\end{tabularx}
\end{center}

\textbf{Interprétation :} avec la poulie fixe, la force est la même qu'en direct ($F_2 = R$) : seule la direction a changé. Avec la poulie mobile, la force est divisée par deux ($F_3 = R/2$), mais il faut tirer deux fois plus de corde.

\textbf{Conclusion :} la poulie fixe ne réduit pas l'effort ; la poulie mobile le divise par deux, au prix d'une course de corde doublée.
\end{experience}

\begin{propriete}[Poulie fixe]
Pour soulever une charge de résistance $R$ avec une \textbf{poulie fixe}, la force motrice nécessaire est égale à la résistance :
\[ F = R \]
La poulie fixe \textbf{ne réduit pas l'effort} ; elle \textbf{change seulement la direction} de la force, ce qui rend le travail plus commode. La longueur de corde tirée est égale à la hauteur de montée de la charge.
\end{propriete}

\courssub{La poulie mobile}

\begin{definition}[Poulie mobile]
Une poulie est dite \cle{mobile} lorsque son \textbf{axe se déplace avec la charge} : la poulie est accrochée au fardeau lui-même. Une extrémité de la corde est fixée à un support, l'autre est tirée par l'opérateur.
\end{definition}

Pourquoi l'effort est-il divisé par deux ? Regardons de près le montage : la charge, suspendue à l'axe de la poulie mobile, est soutenue par \textbf{deux brins de corde} (le brin fixé au support et le brin que l'on tire). Chaque brin supporte donc la moitié du poids. En tirant sur un seul brin, on ne fournit que la moitié de l'effort total.

\begin{propriete}[Poulie mobile]
Pour soulever une charge de résistance $R$ avec une \textbf{poulie mobile}, la force motrice nécessaire est égale à la moitié de la résistance :
\[ F = \dfrac{R}{2} \]
En contrepartie, il faut tirer une longueur de corde \textbf{deux fois plus grande} que la hauteur de montée : pour élever la charge de \SI{1}{m}, on tire \SI{2}{m} de corde.
\end{propriete}

\begin{attention}[Ce qu'on gagne en force, on le perd en distance]
La poulie mobile illustre une loi générale des machines simples : \textbf{on ne gagne jamais sur les deux tableaux}. Si l'effort est divisé par deux, la distance à parcourir est doublée. Le travail (force $\times$ distance) reste le même (dans l'idéal, sans frottement). Ne dis jamais qu'une machine simple « supprime » l'effort : elle le \emph{réduit} en échange d'un déplacement plus long.
\end{attention}

\begin{popfigure}
\begin{tikzpicture}[scale=0.9, >=Stealth]
% ============ POULIE FIXE ============
% Support (poutre)
\fill[pattern=north east lines, pattern color=popmuted] (-1.8,4.6) rectangle (1.8,4.95);
\draw[popdark, thick] (-1.8,4.6) -- (1.8,4.6);
% Axe fixe
\draw[popdark, thick] (0,4.6) -- (0,4.05);
\fill[popdark] (0,3.55) circle (0.07);
% Roue à gorge (cercle extérieur + cercle intérieur pour la gorge)
\draw[popblue, very thick, fill=popblueL] (0,3.55) circle (0.5);
\draw[popblue, thick] (0,3.55) circle (0.35);
% Corde : brin gauche (charge) -> descend de la gorge gauche
\draw[popdark, thick] (-0.5,3.55) -- (-0.5,4.05) arc[start angle=180, end angle=0, radius=0.5] -- (0.5,1.4);
% Brin droit (effort) -> tiré vers le bas
\draw[popdark, thick] (-0.5,4.05) -- (-0.5,1.6);
% Charge (rectangle)
\draw[fill=poporangeL, draw=poporange, thick] (-0.85,1.0) rectangle (-0.15,1.6);
\node[popdark] at (-0.5,1.3) {\small $R$};
% Flèche force F
\draw[popgreen, very thick, ->] (0.5,1.4) -- (0.5,0.4) node[midway, right, popdark]{\small $F = R$};
% Titre
\node[popblue, font=\bfseries] at (0,5.35) {Poulie fixe};
\node[popdark, font=\small] at (0,-0.1) {L'axe est fixe : $F = R$};

% ============ POULIE MOBILE ============
\begin{scope}[xshift=7.5cm]
% Support
\fill[pattern=north east lines, pattern color=popmuted] (-1.8,4.6) rectangle (1.8,4.95);
\draw[popdark, thick] (-1.8,4.6) -- (1.8,4.6);
% Point d'attache fixe de la corde (gauche)
\fill[popdark] (-0.5,4.6) circle (0.07);
% Corde : du point fixe, descend, passe SOUS la poulie mobile, remonte vers la droite
\draw[popdark, thick] (-0.5,4.6) -- (-0.5,2.55) arc[start angle=180, end angle=360, radius=0.5] -- (0.5,4.6);
% Poulie mobile (centre à (0,2.05))
\draw[popgreen, very thick, fill=popgreenL] (0,2.05) circle (0.5);
\draw[popgreen, thick] (0,2.05) circle (0.35);
\fill[popdark] (0,2.05) circle (0.07); % axe mobile
% Lien axe-charge
\draw[popdark, thick] (0,1.55) -- (0,1.3);
% Charge
\draw[fill=poporangeL, draw=poporange, thick] (-0.35,0.7) rectangle (0.35,1.3);
\node[popdark] at (0,1.0) {\small $R$};
% Flèche force F (vers le haut)
\draw[popgreen, very thick, ->] (0.5,4.6) -- (0.5,5.5) node[midway, right, popdark]{\small $F = \dfrac{R}{2}$};
% Titre
\node[popgreen, font=\bfseries] at (0,5.9) {Poulie mobile};
\node[popdark, font=\small] at (0,0.2) {L'axe monte avec la charge : $F = R/2$};
\end{scope}
\end{tikzpicture}
\legende{Comparaison des deux montages. À gauche, la poulie \textbf{fixe} : l'axe est attaché au support, la force $F$ est égale à la résistance $R$, seule la direction change. À droite, la poulie \textbf{mobile} : l'axe est solidaire de la charge, qui est soutenue par deux brins de corde, d'où $F = R/2$.}
\end{popfigure}

\courssub{Exemples de la vie courante}

\begin{exemplebox}[La poulie du puits villageois]
Dans de nombreux villages du Cameroun, une poulie est fixée au-dessus du puits. Le seau plein est lourd : le tirer directement vers le haut, bras tendus au-dessus du vide, serait pénible et dangereux. Grâce à la poulie fixe, on tire la corde \textbf{vers le bas} ou horizontalement, en reculant : le geste est plus sûr et plus confortable. L'effort n'est pas réduit ($F = R$), mais la direction de la force est bien meilleure.
\end{exemplebox}

\begin{exemplebox}[Autres applications]
\begin{itemize}
\item \textbf{Chantiers de construction :} les maçons hissent seaux de mortier et sacs de ciment avec une poulie fixe en haut de l'échafaudage.
\item \textbf{Palan du mécanicien :} pour sortir un moteur de voiture, le mécanicien utilise un palan, c'est-à-dire une \textbf{association de poulies} fixes et mobiles qui divise l'effort par quatre, six ou davantage.
\item \textbf{Grues de chantier :} les grues combinent des poulies (pour les câbles) avec un treuil motorisé.
\end{itemize}
\end{exemplebox}

\begin{popfigure}
\begin{tikzpicture}[scale=0.95, >=Stealth]
% Poteaux du puits
\draw[popdark, very thick] (-1.8,0) -- (-1.8,3.2);
\draw[popdark, very thick] (1.8,0) -- (1.8,3.2);
% Barre transversale
\draw[popdark, very thick] (-2.1,3.2) -- (2.1,3.2);
% Axe poulie
\draw[popdark, thick] (0,3.2) -- (0,2.8);
\fill[popdark] (0,2.4) circle (0.06);
% Poulie
\draw[popblue, very thick, fill=popblueL] (0,2.4) circle (0.4);
\draw[popblue, thick] (0,2.4) circle (0.3);
% Corde : passe sur la poulie
\draw[popdark, thick] (-0.4,2.8) arc[start angle=180, end angle=0, radius=0.4] -- (0.4,1.1);
\draw[popdark, thick] (-0.4,2.8) -- (-0.4,-1.1);
% Seau dans le puits (trapèze)
\draw[fill=popgoldL, draw=popgold, thick] (-0.75,-1.1) -- (-0.6,-1.7) -- (-0.2,-1.7) -- (-0.05,-1.1) -- cycle;
\node[popdark, font=\small] at (-0.4,-1.4) {seau};
% Margelle du puits
\draw[fill=popmuted!40, draw=popdark, thick] (-1.5,-0.4) rectangle (0.9,0);
\draw[popdark, thick] (-1.5,-0.4) -- (-1.5,0);
% Eau (ligne ondulée)
\draw[popblue!60, thick, decorate, decoration={snake, amplitude=0.6pt, segment length=4pt}] (-1.4,-1.9) -- (0.8,-1.9);
\node[popblue, font=\small] at (1.35,-1.9) {eau};
% Flèche force (utilisateur tire vers le bas)
\draw[popgreen, very thick, ->] (0.4,1.1) -- (0.4,0.35) node[midway, right, popdark]{\small $F$};
% Légendes fléchées
\node[popdark, align=left, font=\small] at (3.6,2.4) {1 : poulie fixe};
\draw[popmuted, ->] (2.6,2.4) -- (0.45,2.45);
\node[popdark, align=left, font=\small] at (3.6,1.5) {2 : corde};
\draw[popmuted, ->] (2.9,1.5) -- (0.48,1.6);
\node[popdark, align=left, font=\small] at (-3.9,0.9) {3 : puits};
\draw[popmuted, ->] (-3.3,0.85) -- (-1.15,-0.1);
\end{tikzpicture}
\legende{Puits traditionnel équipé d'une poulie fixe. L'utilisateur tire la corde vers le bas (force $F$) et le seau monte : la poulie change la direction de l'effort, sans le réduire ($F = R$).}
\end{popfigure}

\courssub{Avantages et inconvénients}

\begin{center}
\begin{tabularx}{\linewidth}{|l|X|X|}
\hline
\rowcolor{popblueL} \textbf{Type de poulie} & \textbf{Avantages} & \textbf{Inconvénients} \\
\hline
Poulie fixe & Change la direction de la force : geste confortable et sûr (on tire vers le bas) ; montage simple. & Ne réduit pas l'effort : $F = R$. \\
\hline
Poulie mobile & Réduit l'effort de moitié : $F = R/2$ ; idéale pour les charges lourdes. & Course de corde doublée (2 m de corde pour 1 m de montée) ; montage plus complexe ; la direction de traction n'est pas améliorée (on tire vers le haut). \\
\hline
\end{tabularx}
\end{center}

\begin{savaistu}[Et si l'on combinait les deux ?]
En associant une poulie fixe et une poulie mobile, on obtient le meilleur des deux mondes : l'effort est divisé par deux \emph{et} on tire vers le bas. En multipliant les poulies, on construit un \textbf{palan} : avec deux poulies mobiles, l'effort est divisé par quatre ! C'est le principe des appareils de levage des garages et des grues. Cette association de poulies sera étudiée comme complément en classe de seconde.
\end{savaistu}

\begin{methode}[Distinguer une poulie fixe d'une poulie mobile]
Pour identifier le type de poulie dans un montage ou un schéma :
\begin{enumerate}
\item Repérer l'\textbf{axe} de la poulie.
\item Observer ce qui se passe quand la charge monte : l'axe \textbf{reste immobile} ? C'est une poulie \textbf{fixe} ($F = R$). L'axe \textbf{monte avec la charge} ? C'est une poulie \textbf{mobile} ($F = R/2$).
\item Vérifier le nombre de brins de corde qui soutiennent la charge : un seul brin $\rightarrow$ effort entier ; deux brins $\rightarrow$ effort divisé par deux.
\end{enumerate}
\end{methode}

\begin{attention}[Erreur fréquente : la poulie fixe ne réduit pas l'effort]
Beaucoup d'élèves écrivent au BEPC que la poulie fixe « facilite le levage en réduisant la force ». C'est \textbf{faux} : avec une poulie fixe, $F = R$, l'effort est le même qu'en soulevant directement la charge. Son seul avantage est de \textbf{changer la direction} de la force, ce qui rend le geste plus commode et plus sûr. Seule la poulie \textbf{mobile} (ou une association de poulies) réduit l'effort.
\end{attention}

\begin{exoresolu}[Le seau d'eau du puits]
Un seau rempli d'eau a une résistance (poids) $R = \SI{100}{N}$.\\
\textbf{1.} On le remonte avec une poulie fixe. Quelle force $F$ faut-il exercer ? Dans quel sens ?\\
\textbf{2.} On utilise maintenant une poulie mobile. Quelle force faut-il exercer ?\\
\textbf{3.} Avec la poulie mobile, quelle longueur de corde faut-il tirer pour que le seau monte de \SI{1}{m} ?
\tcblower
\textbf{1.} Avec une poulie fixe, la force est égale à la résistance :
\[ F = R = \SI{100}{N} \]
On tire la corde \textbf{vers le bas} : la poulie fixe change la direction de la force, ce qui est plus confortable, mais ne la réduit pas.

\textbf{2.} Avec une poulie mobile, la charge est soutenue par deux brins de corde, donc :
\[ F = \dfrac{R}{2} = \dfrac{100}{2} = \SI{50}{N} \]

\textbf{3.} Ce qu'on gagne en force, on le perd en distance : la course de corde est doublée. Pour monter le seau de \SI{1}{m}, il faut tirer :
\[ 2 \times 1 = \SI{2}{m} \text{ de corde.} \]
\end{exoresolu}

\begin{exercice}[Sur le chantier]
Un maçon doit hisser un sac de ciment de résistance $R = \SI{500}{N}$ jusqu'au premier étage d'un immeuble, à \SI{3}{m} de hauteur.\\
\textbf{1.} Il utilise une poulie fixe. Quelle force doit-il exercer ? Quelle longueur de corde tire-t-il ?\\
\textbf{2.} Son collègue lui conseille d'ajouter une poulie mobile. Quelle force doit-il alors exercer ? Quelle longueur de corde doit-il tirer ?\\
\textbf{3.} Citer un avantage et un inconvénient de chaque montage.
\end{exercice}

\begin{corrige}[Exercice Sur le chantier]
\textbf{1.} Poulie fixe : $F = R = \SI{500}{N}$. La longueur de corde tirée est égale à la hauteur : \SI{3}{m}.\\
\textbf{2.} Poulie mobile : $F = R/2 = 500/2 = \SI{250}{N}$. Longueur de corde : $2 \times 3 = \SI{6}{m}$.\\
\textbf{3.} Poulie fixe : avantage, on tire vers le bas (confort, sécurité) ; inconvénient, l'effort n'est pas réduit. Poulie mobile : avantage, effort divisé par deux ; inconvénient, il faut tirer deux fois plus de corde.
\end{corrige}

\begin{aretenir}[L'essentiel sur la poulie]
\begin{itemize}
\item Une \textbf{poulie} est une roue à gorge tournant autour d'un axe, sur laquelle passe une corde.
\item \textbf{Poulie fixe} : l'axe est immobile ; elle change la direction de la force mais ne réduit pas l'effort : $F = R$.
\item \textbf{Poulie mobile} : l'axe se déplace avec la charge ; la charge est soutenue par deux brins de corde, donc l'effort est divisé par deux : $F = R/2$ ; en contrepartie, on tire deux fois plus de corde.
\item Pour distinguer les deux : regarder si l'axe de la poulie bouge avec la charge.
\item Applications : puits, chantiers, palans des mécaniciens, grues.
\end{itemize}
\end{aretenir}

\coursec{Le treuil}

Dans la section précédente, nous avons vu que la poulie permet de soulever une charge en tirant sur une corde, et que la poulie mobile divise l'effort par deux. Mais imagine maintenant le pêcheur de Kribi qui doit remonter un filet lourd de poissons, ou le chauffeur de moto-taxi dont la remorque est embourbée : tirer à la main, même avec une poulie, ne suffit plus. Il faut un dispositif capable de multiplier considérablement la force musculaire tout en gardant un contrôle précis du mouvement. Ce dispositif existe depuis des millénaires : c'est le \cle{treuil}. Tu en as peut-être déjà vu un au-dessus d'un puits, pour remonter le seau d'eau. Dans cette section, nous allons définir le treuil, comprendre son fonctionnement en le reliant au levier et à la poulie, établir la relation qui permet de calculer l'effort, puis peser ses avantages et ses inconvénients.

\courssub{Qu'est-ce qu'un treuil ?}

\textbf{Une situation concrète pour commencer.}\par

Dans de nombreux quartiers de Yaoundé ou de Douala, et dans la plupart des villages, on puise l'eau dans un puits. Le seau plein peut peser l'équivalent de \SI{150}{N} à \SI{200}{N}. Le soulever à bout de bras, en tirant directement sur la corde, est vite épuisant, et la corde peut glisser des mains. Alors on installe au-dessus du puits un cylindre horizontal, monté sur deux supports, avec une manivelle sur le côté. En tournant la manivelle, la corde s'enroule autour du cylindre et le seau monte doucement, sans à-coup, avec un effort étonnamment faible. Ce cylindre muni d'une manivelle, c'est un treuil.

\begin{definition}[Le treuil]
Un \cle{treuil} est une machine simple constituée d'un \cle{tambour} (cylindre) autour duquel s'enroule une corde ou un câble, et d'une \cle{manivelle} qui permet de mettre le tambour en rotation. En tournant la manivelle, on enroule la corde et l'on \textbf{soulève} ou l'on \textbf{tire} une charge fixée à son extrémité.
\end{definition}

\textbf{Description des parties d'un treuil.}\par

Un treuil comprend toujours les éléments suivants :

\begin{itemize}
\item le \cle{tambour} : cylindre horizontal, de rayon $r$, autour duquel la corde s'enroule tour après tour ;
\item la \cle{manivelle} : bras de longueur $L$, fixé à l'axe du tambour, sur lequel l'opérateur exerce la force motrice $F$ ;
\item la \cle{corde} ou le \cle{câble} : il relie la charge au tambour ; un tour de manivelle enroule une longueur de corde égale au périmètre du tambour ;
\item le \cle{support} (ou bâti) : il maintient l'axe du tambour et supporte l'ensemble.
\end{itemize}

\begin{popfigure}
\begin{center}
\begin{tikzpicture}[scale=1.0, every node/.style={font=\small}]
% Supports
\draw[fill=popmuted!40, draw=popdark] (-2.6,-0.2) rectangle (-2.2,1.6);
\draw[fill=popmuted!40, draw=popdark] (2.2,-0.2) rectangle (2.6,1.6);
% Tambour (cylindre vu de face)
\draw[fill=poporangeL, draw=popdark, thick] (-2.2,0.5) rectangle (2.2,1.3);
\draw[popdark, thick] (0,0.9) circle (0.12);
\fill[popdark] (0,0.9) circle (0.05);
% Axe
\draw[popdark, thick] (-2.6,0.9) -- (3.1,0.9);
% Manivelle
\draw[popblue, very thick] (3.1,0.9) -- (3.1,2.3);
\draw[popblue, very thick] (3.1,2.3) -- (3.9,2.3);
\draw[fill=popblue, draw=popblue] (3.9,2.3) circle (0.09);
% Annotation L
\draw[<->, popblue, thick] (3.55,0.9) -- (3.55,2.3) node[midway, right, popblue]{$L$};
% Annotation r sur le tambour
\draw[<->, poppink, thick] (0,0.9) -- (0,1.3) node[midway, right, poppink]{$r$};
% Corde enroulée
\draw[popgreen, thick] (-0.9,1.3) arc (90:270:0.12 and 0.4);
\draw[popgreen, thick] (-0.6,1.3) arc (90:270:0.12 and 0.4);
\draw[popgreen, thick] (-0.3,1.3) arc (90:270:0.12 and 0.4);
% Corde qui descend
\draw[popgreen, thick] (-0.3,0.5) -- (-0.3,-1.6);
% Charge
\draw[fill=poppurpleL, draw=popdark, thick] (-0.75,-2.2) rectangle (0.15,-1.6);
\node[popdark] at (-0.3,-1.9) {$R$};
% Labels
\node[popdark, above] at (0,1.35) {tambour};
\node[popblue, right] at (3.95,2.3) {manivelle};
\node[popgreen, left] at (-0.95,0.0) {corde};
\node[popdark, below] at (-2.4,-0.25) {support};
% Force F sur la manivelle
\draw[->, popblue, very thick] (3.9,2.3) -- (3.9,1.7) node[midway, right, popblue]{$F$};
\end{tikzpicture}
\end{center}
\legende{Schéma d'un treuil. La force motrice $F$ est appliquée au bout de la manivelle de longueur $L$ ; la charge $R$ est suspendue à la corde enroulée autour du tambour de rayon $r$.}
\end{popfigure}

\courssub{Principe de fonctionnement : le treuil, un levier combiné à une poulie}

\textbf{Intuition.}\par

Reprenons le seau du puits. Si tu tires la corde directement à la main, tu dois fournir toute la force nécessaire pour équilibrer le poids du seau. Si tu utilises le treuil, ta main n'agit plus sur la corde : elle agit au bout de la manivelle, loin de l'axe de rotation. Or nous avons appris avec le levier que plus le bras de levier est long, plus l'effort est faible. La manivelle joue exactement le rôle d'un bras de levier qui tourne autour de l'axe du tambour.

\begin{propriete}[Principe du treuil]
Le treuil est une \textbf{combinaison du levier et de la poulie} :
\begin{itemize}
\item la \textbf{manivelle} fonctionne comme un \cle{levier} dont le point d'appui est l'axe du tambour : le bras de levier moteur a pour longueur $L$ ;
\item le \textbf{tambour} fonctionne comme une \cle{poulie} de rayon $r$ qui guide et enroule la corde : le bras de levier résistant a pour longueur $r$.
\end{itemize}
L'effort est d'autant plus \textbf{réduit} que la manivelle est \textbf{longue} par rapport au rayon du tambour.
\end{propriete}

\begin{experience}[Mesure de l'effort sur un treuil de laboratoire]
\textbf{Matériel} : un treuil de laboratoire (tambour de rayon $r$ mesurable, manivelle de longueur $L$), une masse marquée (charge $R$), un dynamomètre, une règle graduée, un support.

\textbf{Protocole} :
\begin{enumerate}
\item Mesurer le rayon $r$ du tambour et la longueur $L$ de la manivelle.
\item Suspendre la charge $R$ à la corde enroulée sur le tambour.
\item Accrocher le dynamomètre à l'extrémité de la manivelle et tirer perpendiculairement à la manivelle, juste assez pour maintenir la charge en équilibre.
\item Noter la force $F$ indiquée par le dynamomètre.
\item Recommencer avec une manivelle plus longue, puis avec un tambour de rayon différent.
\end{enumerate}

\textbf{Observations} : la force $F$ mesurée est nettement inférieure à la charge $R$. Quand on allonge la manivelle, $F$ diminue ; quand on utilise un tambour plus gros, $F$ augmente. Dans tous les cas, le produit $F \times L$ reste voisin du produit $R \times r$.

\textbf{Interprétation} : à l'équilibre, les effets de rotation des deux forces autour de l'axe se compensent. L'effet de la force motrice dépend du produit $F \times L$ ; celui de la charge dépend du produit $R \times r$. L'équilibre est réalisé lorsque ces deux produits sont égaux.
\end{experience}

\begin{propriete}[Relation d'équilibre du treuil (admise)]
Lorsque le treuil est en équilibre (charge immobile ou montant à vitesse constante), la force motrice $F$ appliquée au bout de la manivelle et la charge $R$ suspendue à la corde vérifient :
\[ F \times L = R \times r \]
où $L$ est la longueur de la manivelle et $r$ le rayon du tambour, exprimés dans la \textbf{même unité} de longueur. On en déduit :
\[ F = \dfrac{R \times r}{L} \]
Cette relation est admise en classe de 3ème : elle n'est pas à démontrer au BEPC, mais tu dois savoir l'utiliser dans un calcul simple.
\end{propriete}

\begin{attention}[Ne confonds pas $L$ et $r$ !]
L'erreur la plus fréquente au BEPC est de \textbf{confondre le rayon $r$ du tambour avec la longueur $L$ de la manivelle} dans la formule. Retiens bien :
\begin{itemize}
\item $L$ est la longueur de la \textbf{manivelle} (le grand bras, celui de ta main) ;
\item $r$ est le rayon du \textbf{tambour} (le petit bras, celui de la corde).
\end{itemize}
Comme $L$ est toujours plus grand que $r$, la force $F$ est toujours \textbf{plus petite} que la charge $R$. Si ton calcul donne $F > R$, c'est que tu as interverti $L$ et $r$ ! Autre piège : $L$ et $r$ doivent être dans la \textbf{même unité} (par exemple les deux en centimètres) avant de calculer.
\end{attention}

\begin{exoresolu}[Calcul de l'effort sur un treuil de puits]
\textbf{Énoncé.} Pour remonter un seau d'eau, Mamadou utilise un treuil dont le tambour a un rayon $r = \SI{10}{cm}$ et dont la manivelle a une longueur $L = \SI{50}{cm}$. Le seau plein exerce sur la corde une charge $R = \SI{300}{N}$.
\begin{enumerate}
\item Calculer la force $F$ que Mamadou doit exercer au bout de la manivelle pour maintenir le seau en équilibre.
\item Comparer $F$ à $R$ et conclure.
\end{enumerate}
\tcblower
\textbf{Solution détaillée.}
\begin{enumerate}
\item À l'équilibre, la relation du treuil s'écrit :
\[ F \times L = R \times r \quad \text{donc} \quad F = \dfrac{R \times r}{L}. \]
Les deux longueurs sont déjà dans la même unité (le centimètre). Application numérique :
\begin{align*}
F &= \dfrac{300 \times 10}{50} \\
F &= \dfrac{3000}{50} \\
F &= \SI{60}{N}.
\end{align*}
\item Comparaison : $F = \SI{60}{N}$ et $R = \SI{300}{N}$, donc $F = \dfrac{R}{5}$. L'effort est \textbf{cinq fois plus faible} que la charge : le treuil a bien démultiplié la force de Mamadou. Le rapport $\dfrac{L}{r} = \dfrac{50}{10} = 5$ donne directement le facteur de réduction de l'effort.
\end{enumerate}
\end{exoresolu}

\begin{exemplebox}[Influence de la longueur de la manivelle]
Avec la même charge $R = \SI{300}{N}$ et le même tambour $r = \SI{10}{cm}$ :
\begin{itemize}
\item manivelle courte $L = \SI{25}{cm}$ : $F = \dfrac{300 \times 10}{25} = \SI{120}{N}$ ;
\item manivelle moyenne $L = \SI{50}{cm}$ : $F = \SI{60}{N}$ ;
\item manivelle longue $L = \SI{100}{cm}$ : $F = \dfrac{300 \times 10}{100} = \SI{30}{N}$.
\end{itemize}
\textbf{Conclusion} : en doublant la longueur de la manivelle, on divise l'effort par deux. C'est exactement le même comportement que le levier.
\end{exemplebox}

\courssub{Le treuil de puits : un exemple complet}

Le treuil de puits est l'application la plus familière du treuil dans la vie quotidienne camerounaise. Observons-le de près pour identifier chaque partie et vérifier qu'il correspond bien à notre définition.

\begin{popfigure}
\begin{center}
\begin{tikzpicture}[scale=0.95, every node/.style={font=\small}]
% Mur du puits
\draw[fill=popmuted!30, draw=popdark] (-3.2,-2.6) rectangle (-2.6,0.4);
\draw[fill=popmuted!30, draw=popdark] (2.6,-2.6) rectangle (3.2,0.4);
% Poteaux
\draw[fill=popgoldL, draw=popdark, thick] (-2.9,0.4) rectangle (-2.5,2.2);
\draw[fill=popgoldL, draw=popdark, thick] (2.5,0.4) rectangle (2.9,2.2);
% Toit simple
\draw[fill=poporangeL, draw=popdark, thick] (-3.3,2.2) -- (0,3.1) -- (3.3,2.2) -- cycle;
% Tambour
\draw[fill=poporangeL, draw=popdark, thick] (-2.5,1.35) rectangle (2.5,1.85);
% Axe
\draw[popdark, thick] (-2.9,1.6) -- (3.4,1.6);
% Manivelle
\draw[popblue, very thick] (3.4,1.6) -- (3.4,2.6);
\draw[popblue, very thick] (3.4,2.6) -- (4.1,2.6);
\draw[fill=popblue] (4.1,2.6) circle (0.08);
% Corde enroulée
\draw[popgreen, thick] (-0.25,1.85) arc (90:270:0.1 and 0.25);
\draw[popgreen, thick] (0.0,1.85) arc (90:270:0.1 and 0.25);
\draw[popgreen, thick] (0.25,1.85) arc (90:270:0.1 and 0.25);
% Corde verticale
\draw[popgreen, thick] (0.0,1.35) -- (0.0,-1.2);
% Seau
\draw[fill=poppurpleL, draw=popdark, thick] (-0.45,-1.2) -- (-0.3,-1.95) -- (0.3,-1.95) -- (0.45,-1.2) -- cycle;
\draw[popdark] (-0.45,-1.2) arc (180:360:0.45 and 0.12);
% Eau du puits
\draw[fill=popblueL, draw=none] (-2.6,-2.3) rectangle (2.6,-2.6);
\draw[popblue] (-2.6,-2.3) -- (2.6,-2.3);
% Légendes fléchées
\node[popdark] (tamb) at (-3.6,1.6) {\textbf{1}};
\draw[->, popmuted] (-3.35,1.6) -- (-2.55,1.6);
\node[popblue] (man) at (4.55,2.6) {\textbf{2}};
\node[popgreen] (cor) at (-1.5,-0.3) {\textbf{3}};
\draw[->, popmuted] (-1.25,-0.25) -- (-0.1,-0.1);
\node[poppurple] (sea) at (1.6,-1.6) {\textbf{4}};
\draw[->, popmuted] (1.35,-1.6) -- (0.4,-1.6);
\end{tikzpicture}
\end{center}
\legende{Treuil de puits. \textbf{1} : tambour autour duquel s'enroule la corde ; \textbf{2} : manivelle actionnée par l'opérateur ; \textbf{3} : corde ; \textbf{4} : seau (la charge). Le toit protège le puits des feuilles et des saletés.}
\end{popfigure}

\textbf{Fonctionnement pas à pas.}\par

\begin{enumerate}
\item L'opérateur tourne la manivelle dans le sens qui enroule la corde autour du tambour.
\item Chaque tour de manivelle enroule une longueur de corde égale au périmètre du tambour, soit $2 \times \pi \times r$ : le seau monte de cette hauteur à chaque tour.
\item Arrivé en haut, l'opérateur retient la manivelle (ou un cliquet bloque le tambour) : le seau ne retombe pas.
\item Pour redescendre le seau vide, on tourne la manivelle en sens inverse, en contrôlant la vitesse.
\end{enumerate}

\begin{savaistu}[Les treuils des bâtisseurs]
Les grandes cathédrales du Moyen Âge, avec leurs pierres de plusieurs centaines de kilogrammes, ont été bâties grâce à d'immenses treuils actionnés par des hommes qui marchaient à l'intérieur d'une grande roue, comme des écureuils dans une cage : leur poids faisait tourner le tambour qui hissait les pierres. Au Cameroun, les cases à étages et les greniers traditionnels étaient également construits en hissant matériaux et récoltes à l'aide de treuils et de poulies. Aujourd'hui encore, les pêcheurs de Kribi et de Limbe utilisent des enrouleurs — de véritables treuils — pour remonter leurs filets lourds de poissons sans se blesser les mains, et les dépanneuses qui sortent les camions embourbés de la route en saison des pluies sont équipées de puissants treuils, manuels ou motorisés.
\end{savaistu}

\courssub{Avantages et inconvénients du treuil}

\begin{aretenir}[Avantages du treuil]
\begin{itemize}
\item \textbf{Effort fortement réduit} : la force motrice est divisée par le rapport $\dfrac{L}{r}$, qui peut être grand (5, 10 ou plus). Une seule personne peut soulever des charges très lourdes.
\item \textbf{Contrôle précis} : la montée et la descente de la charge sont lentes et régulières ; on peut arrêter la charge à n'importe quelle hauteur, ce qui est précieux pour les manœuvres délicates (dépannage d'un véhicule, remontée d'un filet).
\item \textbf{Sécurité} : les mains ne touchent pas directement la corde tendue ; un cliquet de blocage peut empêcher la charge de retomber.
\end{itemize}
\end{aretenir}

\begin{attention}[Inconvénients du treuil]
\begin{itemize}
\item \textbf{Lenteur} : chaque tour de manivelle ne fait monter la charge que d'un périmètre de tambour ($2\pi r$). Pour remonter un seau de \SI{10}{m} avec un tambour de rayon \SI{10}{cm} (périmètre environ \SI{63}{cm}), il faut environ 16 tours de manivelle ! L'effort est réduit, mais le travail se paie en \textbf{nombre de tours} : on ne gagne jamais en travail, on gagne seulement en force.
\item \textbf{Encombrement} : le tambour, la manivelle et le support occupent une place importante, surtout si la manivelle est longue.
\item \textbf{Usure et entretien} : la corde ou le câble s'use en s'enroulant ; un câble effiloché peut casser. Il faut inspecter régulièrement le câble et ne jamais se placer sous la charge suspendue.
\end{itemize}
\end{attention}

\begin{methode}[Résoudre un exercice sur le treuil]
\begin{enumerate}
\item \textbf{Repérer} dans l'énoncé la charge $R$ (en newtons), la longueur $L$ de la manivelle et le rayon $r$ du tambour. Attention : si l'on donne le \emph{diamètre} du tambour, divise-le par deux pour obtenir $r$.
\item \textbf{Convertir} $L$ et $r$ dans la même unité de longueur (par exemple en centimètres).
\item \textbf{Écrire} la relation d'équilibre : $F \times L = R \times r$.
\item \textbf{Isoler} la grandeur cherchée, par exemple $F = \dfrac{R \times r}{L}$.
\item \textbf{Calculer} puis \textbf{vérifier} : la force $F$ doit être plus petite que la charge $R$. Conclure par une phrase avec l'unité.
\end{enumerate}
\end{methode}

\begin{exercice}[La remorque embourbée]
Pour sortir sa remorque de la boue, un chauffeur utilise un treuil dont le tambour a un diamètre de \SI{16}{cm} et dont la manivelle mesure \SI{40}{cm}. La résistance à vaincre est $R = \SI{500}{N}$.
\begin{enumerate}
\item Quel est le rayon $r$ du tambour ?
\item Calculer la force $F$ que le chauffeur doit exercer sur la manivelle.
\item Le chauffeur remplace la manivelle par une autre, deux fois plus longue. Que devient la force $F$ ?
\end{enumerate}
\end{exercice}

\begin{corrige}[Exercice : la remorque embourbée]
\begin{enumerate}
\item Le rayon est la moitié du diamètre : $r = \dfrac{16}{2} = \SI{8}{cm}$.
\item À l'équilibre : $F \times L = R \times r$, donc $F = \dfrac{R \times r}{L} = \dfrac{500 \times 8}{40} = \dfrac{4000}{40} = \SI{100}{N}$.
\item La nouvelle manivelle mesure $L' = \SI{80}{cm}$. Alors $F' = \dfrac{500 \times 8}{80} = \SI{50}{N}$. En doublant la longueur de la manivelle, l'effort est divisé par deux, ce qui confirme que le treuil se comporte comme un levier.
\end{enumerate}
\end{corrige}

\begin{aretenir}[L'essentiel sur le treuil]
\begin{itemize}
\item Le treuil est une machine simple formée d'un \textbf{tambour} (rayon $r$) et d'une \textbf{manivelle} (longueur $L$) ; une corde s'enroule sur le tambour pour soulever ou tirer une charge.
\item Il combine le \textbf{levier} (la manivelle) et la \textbf{poulie} (le tambour).
\item À l'équilibre : $F \times L = R \times r$, soit $F = \dfrac{R \times r}{L}$. Plus la manivelle est longue et le tambour fin, plus l'effort est faible.
\item Avantages : effort fortement réduit, contrôle précis, sécurité. Inconvénients : lenteur (beaucoup de tours de manivelle) et encombrement.
\item Exemples : treuil de puits, manivelle de remorque, treuil de dépannage, enrouleur de filets des pêcheurs de Kribi.
\end{itemize}
\end{aretenir}

Le treuil achève notre tour d'horizon des machines simples : levier, plan incliné, poulie et treuil obéissent tous à la même grande idée — on ne supprime jamais le travail à fournir, on le rend simplement plus facile en réduisant la force au prix d'un déplacement plus grand. Il nous reste à dresser le bilan comparatif de ces machines et à dégager les règles de sécurité qui accompagnent leur usage : ce sera l'objet de la dernière section de cette leçon.

\coursec{Bilan comparatif et sécurité}

Nous venons d'étudier le \cle{treuil}, cette machine qui enroule une corde autour d'un tambour grâce à une manivelle. Avec lui s'achève notre tour des quatre machines simples du programme : le \cle{levier}, le \cle{plan incliné}, la \cle{poulie} et le \cle{treuil}. Il est temps de faire le bilan : comparer ces machines, comprendre ce qu'elles ont en commun, apprendre à choisir la bonne machine selon la situation, et surtout retenir les règles de sécurité qui sauvent des vies sur les chantiers, dans les ateliers et même à la maison.

\courssub{Tableau comparatif des quatre machines simples}

Commençons par rassembler dans un seul tableau tout ce que nous avons appris. Lis-le ligne par ligne, puis colonne par colonne : c'est un excellent outil de révision pour le BEPC.

\begin{center}
\begin{tabularx}{\linewidth}{|l|X|X|X|X|}
\hline
\rowcolor{popblueL} \textbf{Machine} & \textbf{Définition et principe} & \textbf{Avantage principal} & \textbf{Inconvénient principal} & \textbf{Exemple camerounais} \\
\hline
\textbf{Levier} & Tige rigide tournant autour d'un point d'appui ; l'effort appliqué loin du point d'appui soulève une charge proche de celui-ci. & Forte réduction de l'effort possible (rapport des bras de levier). & La charge ne se déplace que sur une courte distance ; il faut un point d'appui solide. & Soulever une pierre avec une barre de fer ; décapsuler une bouteille ; la brouette au marché Mokolo. \\
\hline
\textbf{Plan incliné} & Surface plane en pente : on déplace la charge le long de la pente au lieu de la soulever verticalement. & Effort réduit d'autant plus que la pente est douce. & Le chemin parcouru est plus long ; il faut de la place ; frottements parfois importants. & Rampe en planches pour charger des sacs de maïs dans un camion ; route en lacets vers Dschang. \\
\hline
\textbf{Poulie} & Roue à gorge dans laquelle passe une corde ; fixe, elle change la direction de l'effort ; mobile, elle le divise par deux. & Confort (on tire vers le bas) et, avec la poulie mobile, effort divisé par deux. & La poulie fixe ne réduit pas l'effort ; la poulie mobile exige de tirer deux fois plus de corde. & Puits du village pour remonter l'eau ; palan du mécanicien pour sortir un moteur. \\
\hline
\textbf{Treuil} & Tambour sur lequel s'enroule une corde, actionné par une manivelle ; le rapport rayon de manivelle / rayon du tambour fixe le gain. & Très forte réduction de l'effort ; la charge peut être bloquée en position. & Manœuvre lente : beaucoup de tours de manivelle pour peu de hauteur. & Treuil de puits profond ; treuil de dépannage pour sortir un véhicule embourbé en saison des pluies. \\
\hline
\end{tabularx}
\end{center}

\begin{attention}[Bien lire le tableau]
Au BEPC, ne confonds pas les rôles : la \cle{poulie fixe} ne fait que \emph{changer la direction} de l'effort (l'effort reste égal au poids de la charge), alors que la \cle{poulie mobile} \emph{divise l'effort par deux}. De même, le plan incliné ne supprime pas l'effort, il le \emph{répartit sur une plus longue distance}.
\end{attention}

\courssub{Le point commun de toutes les machines simples}

Regarde à nouveau le tableau. Que remarques-tu ? Aucune machine ne fait disparaître l'effort gratuitement. Quand l'effort diminue, c'est toujours parce que \emph{quelque chose d'autre augmente} : la longueur du bras de levier, la longueur de la rampe, la longueur de corde à tirer, le nombre de tours de manivelle.

\begin{experience}[Observation : le prix du confort]
Reprenons la rampe de chargement du camion. Un sac de \SI{50}{kg} pèse environ \SI{500}{N}. Pour le monter directement à \SI{1,5}{m} du sol, il faut fournir un effort de \SI{500}{N} sur \SI{1,5}{m}. Avec une rampe de \SI{4,5}{m} (trois fois plus longue), l'effort tombe à environ \SI{170}{N}, mais il faut pousser le sac sur trois fois plus de distance, donc trois fois plus de temps. \textbf{Interprétation :} ce qu'on gagne en force, on le perd en distance (et en temps). Le produit « force $\times$ distance », c'est-à-dire le \emph{travail}, reste le même \emph{dans l'idéal (sans frottements)}.
\end{experience}

\begin{propriete}[Synthèse : la loi commune des machines simples]
Toute machine simple \cle{transforme l'effort} : elle peut en réduire l'intensité ou en changer la direction, mais jamais gratuitement.
\begin{itemize}
\item \textbf{Gain en force = perte en distance.} Si l'effort est divisé par $n$, la distance à parcourir est multipliée par $n$.
\item \textbf{Aucune machine ne crée d'énergie.} Le travail fourni par l'opérateur est au mieux égal (sans frottements) au travail reçu par la charge ; en réalité, il est toujours un peu supérieur à cause des frottements.
\end{itemize}
C'est pourquoi on dit qu'une machine simple est un \emph{transformateur d'effort}, pas un « multiplicateur de force » magique.
\end{propriete}

\begin{exemplebox}[Soulever une charge de \SI{200}{N} : comparatif chiffré]
Un jerrican rempli d'eau pèse \SI{200}{N}. Comparons les quatre façons de le soulever :
\begin{itemize}
\item \textbf{Directement (sans machine) :} effort = \SI{200}{N}.
\item \textbf{Levier} avec un rapport de bras de levier de 4 : effort = $200 \div 4 = \SI{50}{N}$.
\item \textbf{Poulie mobile} (deux brins de corde) : effort = $200 \div 2 = \SI{100}{N}$, mais il faut tirer \SI{2}{m} de corde pour monter le jerrican de \SI{1}{m}.
\item \textbf{Treuil} avec un rapport manivelle/tambour de 5 : effort = $200 \div 5 = \SI{40}{N}$, mais il faut tourner la manivelle sur une longueur cinq fois plus grande.
\end{itemize}
Conclusion : le treuil est le plus « économique » en effort, mais le plus lent. Le levier est rapide mais demande un bon point d'appui. À chaque situation sa machine !
\end{exemplebox}

\courssub{Comment choisir la bonne machine simple ?}

Sur un chantier, à la ferme ou à l'atelier, le choix d'une machine simple n'est pas au hasard. Voici les questions à se poser, dans l'ordre :

\begin{methode}[Choisir une machine simple selon la situation]
\begin{enumerate}
\item \textbf{Évaluer la charge.} Est-elle lourde (plusieurs centaines de newtons) ? Alors il faut un fort rapport de réduction : treuil ou palan (poulies combinées).
\item \textbf{Mesurer la hauteur ou la distance.} Pour soulever haut (puits, étage), poulie ou treuil. Pour déplacer sur une courte hauteur (chargement d'un camion), plan incliné.
\item \textbf{Observer l'espace disponible.} Peu de place ? Le levier est compact. Beaucoup de place ? La rampe douce devient possible et très efficace.
\item \textbf{Tenir compte de la rapidité souhaitée.} Besoin de vitesse : levier ou poulie fixe. La vitesse importe peu mais l'effort doit être minimal : treuil.
\item \textbf{Vérifier les points d'ancrage.} Pas de support solide au-dessus ? Oublie la poulie fixe. Pas de point d'appui ? Oublie le levier.
\end{enumerate}
\end{methode}

\begin{exoresolu}[Quelle machine pour quel chantier ?]
\textbf{Énoncé.} Dans chaque cas, choisis la machine simple la plus adaptée et justifie ton choix :
\begin{enumerate}
\item Un maçon doit monter des seaux de mortier de \SI{150}{N} à un échafaudage situé à \SI{4}{m}, avec une poutre solide au-dessus.
\item Un chauffeur doit charger des sacs de cacao de \SI{650}{N} dans une remorque haute de \SI{1,2}{m}, dans une grande cour dégagée.
\item Un mécanicien doit extraire un moteur très lourd d'une vieille voiture, sans se presser, seul dans son atelier.
\end{enumerate}
\tcblower
\textbf{Solution détaillée.}
\begin{enumerate}
\item \textbf{Poulie fixe} accrochée à la poutre : le maçon tire vers le bas (position confortable et sûre, il peut utiliser son poids), l'effort reste de \SI{150}{N}, ce qui est acceptable pour un seau. Le point d'ancrage existe, la hauteur est grande : la poulie fixe s'impose.
\item \textbf{Plan incliné} (solides planches entre le sol et la remorque) : la cour est dégagée, donc la place ne manque pas ; les sacs sont lourds mais la hauteur est faible. En roulant ou faisant glisser les sacs sur la rampe, l'effort est fortement réduit. Une poulie serait difficile à ancrer ici.
\item \textbf{Treuil} (ou palan) : la charge est très lourde, le mécanicien est seul et n'est pas pressé. Le fort rapport de réduction du treuil lui permet de soulever le moteur avec un petit effort, quitte à tourner longtemps la manivelle.
\end{enumerate}
\textbf{À retenir :} le bon choix dépend toujours de quatre critères : la charge, la hauteur, l'espace et la rapidité souhaitée.
\end{exoresolu}

\begin{methode}[Identifier une machine simple cachée dans un objet]
Beaucoup d'objets de la vie courante cachent une machine simple. Pour la reconnaître, cherche l'élément caractéristique :
\begin{enumerate}
\item Un \textbf{point d'appui} autour duquel tourne une tige $\Rightarrow$ c'est un \cle{levier} (ciseaux, pince, brouette, pédale).
\item Une \textbf{pente} le long de laquelle on déplace la charge $\Rightarrow$ c'est un \cle{plan incliné} (rampe, escalier, vis dont le filet est un plan incliné enroulé).
\item Une \textbf{roue à gorge} avec une corde $\Rightarrow$ c'est une \cle{poulie} (poulie de puits, dérailleur de vélo).
\item Un \textbf{tambour} qui enroule une corde, actionné par une \textbf{manivelle} $\Rightarrow$ c'est un \cle{treuil} (moulinet de pêche, manivelle de store).
\end{enumerate}
\end{methode}

\courssub{La sécurité : des règles qui protègent la vie}

Une machine simple mal utilisée transforme un outil utile en danger. Une corde usée qui casse, une charge qui retombe, un levier qui échappe : les accidents sont fréquents sur les chantiers et dans les ateliers. Voici les consignes à connaître et à appliquer sans exception.

\begin{attention}[Consignes de sécurité essentielles]
\begin{itemize}
\item \textbf{Avant l'usage}, vérifie toujours l'état des \cle{cordes}, \cle{câbles}, crochets et supports : aucune fibre effilochée, aucune rouille profonde, aucune fissure.
\item \textbf{Ne dépasse jamais la charge maximale} indiquée sur l'appareil (palan, treuil, sangle).
\item \textbf{Ne te place jamais sous une charge suspendue}, même une seconde, même pour ramasser un outil. Si la corde casse, la charge tombe verticalement.
\item \textbf{Porte des équipements de protection} : gants contre les brûlures de corde et les coupures, chaussures fermées (idéalement à embout de fer) contre les chutes d'objets.
\item \textbf{Garde une posture correcte} : dos droit, jambes fléchies pour pousser ou tirer ; ne te penche pas au-dessus d'un levier sous tension.
\item \textbf{Communique} : sur un chantier, annonce toujours à voix haute avant de soulever ou de descendre une charge.
\end{itemize}
\end{attention}

\begin{popfigure}
\begin{center}
\begin{tikzpicture}[scale=0.9, font=\small]
% Potence
\draw[line width=2.5pt, popdark] (0,0) -- (0,4.2) -- (3.2,4.2);
\draw[line width=1.5pt, popdark] (0,2.6) -- (1.6,4.2);
% Poulie
\draw[fill=popblueL, draw=popblue, line width=1pt] (2.8,4.2) circle (0.3);
\fill[popblue] (2.8,4.2) circle (0.05);
% Corde et charge
\draw[line width=1.2pt, popdark] (2.8,4.2) -- (2.8,2.2);
\draw[fill=popgoldL, draw=popgold, line width=1pt] (2.2,1.4) rectangle (3.4,2.2);
\node at (2.8,1.8) {\textbf{CHARGE}};
% Zone interdite (hachurée)
\fill[pattern=north east lines, pattern color=poporange] (1.6,0) rectangle (4.0,1.4);
\draw[poporange, line width=1.2pt, dashed] (1.6,0) rectangle (4.0,1.4);
\node[poporange, font=\small\bfseries] at (2.8,0.7) {ZONE INTERDITE};
% Personnage écarté
\draw[fill=popgreenL, draw=popgreen, line width=1pt] (5.6,0.9) circle (0.22);
\draw[line width=1.4pt, popgreen] (5.6,0.68) -- (5.6,0.1);
\draw[line width=1.4pt, popgreen] (5.6,0.5) -- (5.15,0.25);
\draw[line width=1.4pt, popgreen] (5.6,0.5) -- (6.05,0.25);
\draw[line width=1.4pt, popgreen] (5.6,0.1) -- (5.3,-0.35);
\draw[line width=1.4pt, popgreen] (5.6,0.1) -- (5.9,-0.35);
\node[popgreen, font=\small] at (5.6,-0.65) {Opérateur à distance};
% Flèche de chute
\draw[-{Stealth[length=3mm]}, line width=1.5pt, poppink] (4.6,2.0) -- (4.6,0.6);
\node[poppink, font=\small, align=center] at (5.35,1.55) {chute\\possible};
% Sol
\draw[line width=1.5pt, popdark] (-0.6,-0.35) -- (6.6,-0.35);
\end{tikzpicture}
\end{center}
\legende{Schéma de signalisation : la zone située sous une charge suspendue (hachurée en orange) est strictement interdite. L'opérateur se tient à distance, jamais dans l'axe de chute de la charge.}
\end{popfigure}

\courssub{Entretenir les machines simples}

Une machine bien entretenue est une machine sûre et durable. Deux gestes simples suffisent dans la plupart des cas :

\begin{itemize}
\item \textbf{Le graissage des axes.} L'axe de la poulie, le tambour du treuil, la charnière du levier doivent être graissés régulièrement (huile ou graisse). Le graissage réduit les \cle{frottements} : l'effort à fournir diminue, les pièces s'usent moins vite et ne chauffent pas. C'est la même raison pour laquelle on huile la chaîne d'une moto-taxi.
\item \textbf{Le remplacement des pièces usées.} Une corde effilochée, un câble dont des brins sont rompus, une manivelle fissurée doivent être remplacés \emph{immédiatement}, pas « plus tard ». Une corde qui a perdu un tiers de ses fibres peut casser sous une charge pourtant normale.
\end{itemize}

\begin{exemplebox}[Le puits du quartier]
Dans beaucoup de quartiers et de villages, le puits est équipé d'une poulie ou d'un treuil. Un comité d'entretien vérifie chaque mois la corde, graisse l'axe et remplace la corde dès les premiers signes d'usure. Résultat : aucun seau perdu au fond du puits et aucun accident depuis des années. L'entretien préventif coûte bien moins cher que l'accident.
\end{exemplebox}

\courssub{Des machines simples aux machines complexes}

Terminons par une idée importante : les machines simples ne sont pas des curiosités de manuel scolaire. Ce sont les \emph{briques de base} de toutes les machines modernes.

\begin{itemize}
\item La \textbf{grue} de chantier = un \cle{treuil} (moteur + tambour) associé à des \cle{poulies} (le palan) montées sur une flèche.
\item La \textbf{pelle mécanique} = un enchaînement de \cle{leviers} articulés (flèche, balancier, godet) actionnés par des vérins.
\item Le \textbf{vélo} = pédalier (manivelle de treuil), chaîne et pignons, freins et guidon (leviers).
\item La \textbf{vis} et l'\textbf{escalier} = des plans inclinés enroulés ou répétés.
\end{itemize}

\begin{savaistu}[Ton corps est une machine à leviers]
Le corps humain utilise des leviers en permanence ! Quand tu soulèves un cahier, ton avant-bras tourne autour du coude (le point d'appui), le poids du cahier est la charge, et le muscle \emph{biceps} fournit l'effort en tirant sur l'avant-bras près du coude. Comme l'effort est appliqué \emph{entre} le point d'appui et la charge, on parle de \textbf{levier inter-effort}. Ce levier ne réduit pas l'effort — le biceps doit tirer fort — mais il permet un mouvement rapide et ample de la main. Mâchoire, pied, nuque : tout ton squelette fonctionne ainsi, avec des muscles pour moteurs et des articulations pour points d'appui.
\end{savaistu}

\begin{aretenir}[L'essentiel de la leçon]
\begin{itemize}
\item Les quatre machines simples — \cle{levier}, \cle{plan incliné}, \cle{poulie}, \cle{treuil} — réduisent l'effort ou en changent la direction.
\item \textbf{Gain en force = perte en distance} : aucune machine ne crée d'énergie, le travail se conserve (aux frottements près).
\item Le choix d'une machine dépend de la charge, de la hauteur, de l'espace disponible et de la rapidité souhaitée.
\item Sécurité : vérifier cordes et supports, respecter la charge maximale, ne jamais se placer sous une charge suspendue, porter gants et chaussures de protection.
\item Entretien : graisser les axes, remplacer sans délai les cordes et câbles usés.
\item Toutes les machines complexes (grue, pelle mécanique, vélo) sont des combinaisons de machines simples.
\end{itemize}
\end{aretenir}

\begin{exercice}[Pour t'entraîner]
\begin{enumerate}
\item Recopie et complète : « Ce qu'une machine simple fait gagner en \ldots{}, elle le fait perdre en \ldots{} ».
\item Pour soulever une charge de \SI{300}{N}, on utilise un treuil dont le rapport manivelle/tambour vaut 6. Quel effort faut-il fournir ?
\item Cite trois consignes de sécurité à respecter avant d'utiliser un palan.
\item Dans une paire de ciseaux, identifie le point d'appui, la charge et l'effort. De quelle machine simple s'agit-il ?
\end{enumerate}
\end{exercice}

\begin{corrige}[Exercice n°1]
\begin{enumerate}
\item « Ce qu'une machine simple fait gagner en \textbf{force}, elle le fait perdre en \textbf{distance} (et en temps). »
\item Effort = $300 \div 6 = \SI{50}{N}$.
\item Vérifier l'état des cordes et crochets ; ne pas dépasser la charge maximale ; ne jamais se placer sous la charge suspendue (porter gants et chaussures de protection est aussi accepté).
\item Le point d'appui est la vis centrale ; la charge est la résistance du matériau coupé ; l'effort est appliqué par les doigts sur les anneaux. C'est un \cle{levier} (double levier inter-appui).
\end{enumerate}
\end{corrige}

\newpage

\coursec{Activité d'intégration}

\coursec{Leçon 11 : Quelques machines simples}

\courssub{Objectifs de l'activité}

À la fin de cette activité, tu seras capable de :
\begin{itemize}
\item \cle{identifier} une machine simple (levier, plan incliné, poulie, treuil) dans un dispositif concret de chantier ou de la vie courante ;
\item \cle{exploiter} des mesures expérimentales pour vérifier qu'un dispositif permet de rester en dessous d'un effort maximal imposé ;
\item \cle{appliquer} les relations d'équilibre des machines simples : $F \times d_1 = R \times d_2$ (levier, treuil) et $F \times L = R \times h$ (plan incliné) ;
\item \cle{comparer} les avantages et les inconvénients de plusieurs dispositifs (effort, distance parcourue, rapidité, sécurité, coût) ;
\item \cle{rédiger} une conclusion argumentée et des consignes de sécurité adaptées à un chantier réel.
\end{itemize}

\courssub{Situation déclenchante : Le chantier de l'école}

\textbf{Contexte.} L'école publique de ton quartier construit une nouvelle salle de classe au premier étage. Les maçons doivent monter des sacs de ciment \textbf{CIMENCAM} de \SI{50}{kg} chacun jusqu'au plancher, situé à une hauteur $h = \SI{3}{m}$ au-dessus du sol. Le poids d'un sac est $R = P = m \times g = 50 \times 10 = \SI{500}{N}$.

\textbf{Problème.} Le chef de chantier constate qu'un ouvrier, même costaud, ne peut fournir en continu qu'un \cle{effort maximal} $F_{\max} = \SI{150}{N}$ sans risquer de se blesser au dos (lombalgies). Porter les sacs à la force des bras est donc impossible. Il faut utiliser une \cle{machine simple}.

\textbf{Trois dispositifs sont proposés par les ouvriers :}
\begin{itemize}
\item \textbf{Dispositif A} : une planche solide de longueur $L = \SI{6}{m}$, posée en pente entre le sol et le plancher (plan incliné). On tire le sac avec une corde parallèle à la pente.
\item \textbf{Dispositif B} : une \cle{poulie mobile} accrochée au sac, avec une corde dont une extrémité est fixée en haut (au plancher) et l'autre tirée par l'ouvrier vers le bas.
\item \textbf{Dispositif C} : un \cle{treuil} fixé au plancher : un cylindre de rayon $r = \SI{10}{cm}$ sur lequel s'enroule la corde, actionné par une manivelle de longueur $d_1 = \SI{40}{cm}$.
\end{itemize}

Pour trancher, le professeur de PCT a fait réaliser par les élèves de 3ème une série de mesures au dynamomètre (le sac est soulevé ou déplacé à vitesse constante, frottements négligés) :

\begin{center}
\begin{tabularx}{\linewidth}{|l|X|X|X|}
\hline
\rowcolor{popblueL} \textbf{Dispositif} & \textbf{Force mesurée $F$} & \textbf{Distance parcourue par la force} & \textbf{Hauteur gagnée par le sac} \\
\hline
A : plan incliné ($L = \SI{6}{m}$) & \SI{250}{N} (trop fort !) & \SI{6}{m} & \SI{3}{m} \\
\hline
A' : plan incliné allongé ($L = \SI{12}{m}$) & \SI{125}{N} & \SI{12}{m} & \SI{3}{m} \\
\hline
B : poulie mobile (1 ouvrier) & \SI{250}{N} (trop fort !) & \SI{6}{m} de corde tirée & \SI{3}{m} \\
\hline
B' : poulie mobile (2 ouvriers) & \SI{125}{N} par ouvrier & \SI{6}{m} de corde & \SI{3}{m} \\
\hline
C : treuil & à calculer & manivelle : tours à compter & \SI{3}{m} \\
\hline
\end{tabularx}
\end{center}

Le chef de chantier te demande, à toi élève de 3ème, de l'aider à \cle{choisir} le dispositif le plus adapté, de le \cle{justifier par des calculs} et de \cle{rédiger les consignes de sécurité} à afficher sur le chantier.

\courssub{Consignes de l'activité}

\begin{enumerate}
\item Identifier la machine simple utilisée dans chacun des dispositifs A, B et C. Pour chacune, citer un exemple d'utilisation dans la vie courante au Cameroun.
\item Rappeler la valeur du poids $R$ d'un sac de ciment et l'effort maximal $F_{\max}$ qu'un ouvrier peut fournir. Calculer le rapport $R/F_{\max}$. Que doit donc faire la machine simple pour être acceptable ?
\item \textbf{Plan incliné} : en utilisant la relation $F \times L = R \times h$, vérifier que la mesure du dispositif A ($F = \SI{250}{N}$) est cohérente. Calculer ensuite la longueur $L$ minimale de la planche pour que l'effort descende à \SI{150}{N}. La mesure du dispositif A' confirme-t-elle ton calcul ?
\item \textbf{Poulie mobile} : rappeler pourquoi une poulie mobile divise l'effort par 2. Vérifier que la mesure du dispositif B est cohérente. Expliquer pourquoi la solution B' (deux ouvriers tirant ensemble) permet de respecter la limite de \SI{150}{N}.
\item \textbf{Treuil} : la manivelle joue le rôle de levier. En appliquant $F \times d_1 = R \times r$, calculer l'effort $F$ à exercer sur la manivelle pour soulever le sac. Comparer à $F_{\max}$.
\item \textbf{Comparaison} : compléter un tableau comparatif des trois solutions valables (A', B', C) selon les critères : effort fourni, distance parcourue par la force, rapidité, sécurité, coût et simplicité d'installation, main-d'œuvre nécessaire.
\item \textbf{Conclusion} : rédiger un court texte (5 à 8 lignes) recommandant au chef de chantier un dispositif, avec au moins deux arguments chiffrés.
\item \textbf{Sécurité} : proposer quatre consignes de sécurité à afficher sur le chantier.
\end{enumerate}

\courssub{Ressources à mobiliser}

\begin{definition}[Machine simple]
Un dispositif mécanique élémentaire qui permet de modifier l'intensité, la direction ou le point d'application d'une force. Il ne crée pas d'énergie : il échange de la force contre de la distance. Les machines simples de base sont : le levier, le plan incliné, la poulie, le treuil, le coin, la vis.
\end{definition}

\begin{propriete}[Règle d'or des machines simples (frottements négligés)]
Le travail moteur $W_m = F \times d_F$ fourni par l'utilisateur est égal au travail résistant $W_r = R \times d_R$ effectué sur la charge.
\[ F \times d_F = R \times d_R \quad \Longrightarrow \quad \frac{F}{R} = \frac{d_R}{d_F} \]
On \textbf{gagne en force} ($F < R$) mais on \textbf{perd en distance} ($d_F > d_R$).
\end{propriete}

\begin{aretenir}[Formules d'équilibre des machines simples du programme]
\begin{itemize}
\item \textbf{Poids d'un corps} : $P = m \times g$, avec $g = \SI{10}{N/kg}$ (valeur conventionnelle au collège).
\item \textbf{Levier} (et treuil) : à l'équilibre, $F \times d_1 = R \times d_2$.
  \begin{itemize}
  \item $F$ : effort moteur (en N), $d_1$ : bras de levier de la force (distance au point d'appui).
  \item $R$ : résistance (poids de la charge en N), $d_2$ : bras de levier de la résistance.
  \end{itemize}
\item \textbf{Plan incliné} (sans frottement, force parallèle à la pente) : $F \times L = R \times h$.
  \begin{itemize}
  \item $L$ : longueur de la pente (m), $h$ : hauteur verticale gagnée (m).
  \item Rapport de démultiplication : $\dfrac{L}{h} = \dfrac{1}{\sin\alpha}$.
  \end{itemize}
\item \textbf{Poulie mobile} : l'effort est divisé par 2 ($F = R/2$), mais la longueur de corde à tirer est doublée ($d_F = 2 \times h$).
\item \textbf{Treuil} : cas particulier du levier en rotation. $d_1 = \text{longueur manivelle}$, $d_2 = \text{rayon du tambour } r$.
\end{itemize}
\end{aretenir}

\begin{savaistu}[Au Cameroun, les machines simples sont partout]
\begin{itemize}
\item \textbf{Levier} : la brouette (roue = appui), le décapsuleur, la barre à mine pour soulever un bloc de béton.
\item \textbf{Plan incliné} : la rampe pour charger une moto-taxi dans un camion, la route en lacets pour monter à Dschang ou à Buéa, le toboggan.
\item \textbf{Poulie mobile / Palan} : le treuil du puits de village (poulie fixe), les palans des garages pour descendre un moteur, les grues de chantier (poulies composées).
\item \textbf{Treuil} : le treuil de la benne à ordures, le cric de voiture (vis + levier), le treuil de remorquage des 4x4.
\end{itemize}
\end{savaistu}

\courssub{Démarche d'investigation et corrigé commenté}

\begin{exoresolu}[Le chantier de l'école : corrigé complet et commenté]
Identifier chaque machine, vérifier les efforts par le calcul, comparer les dispositifs valables, recommander une solution et rédiger les consignes de sécurité.
\tcblower

\textbf{1. Identification des machines simples et exemples locaux.}
\begin{itemize}
\item \textbf{Dispositif A (et A')} : la planche en pente est un \cle{plan incliné}. \emph{Exemple camerounais :} la rampe en bois ou en tôle pour monter une moto-taxi \emph{« bendskin »} sur un camion de transport.
\item \textbf{Dispositif B (et B')} : c'est une \cle{poulie mobile} (la poulie est solidaire de la charge). \emph{Exemple camerounais :} le palan à chaîne utilisé dans les garages (Makepe, Bassa) pour sortir un moteur de voiture ; le système de levage des seaux d'eau au puits avec une poulie mobile (rare, souvent poulie fixe).
\item \textbf{Dispositif C} : c'est un \cle{treuil}. La manivelle et le cylindre forment un \cle{levier de 1ère espèce} en rotation continue. \emph{Exemple camerounais :} le treuil manuel du puits traditionnel pour remonter le seau ; le treuil de la benne à ordures ménagères ; le cric de voiture (vis sans fin + levier).
\end{itemize}

\textbf{2. Données de départ et critère de choix.}
\begin{itemize}
\item Poids du sac (résistance) : $R = m \times g = 50 \times 10 = \SI{500}{N}$.
\item Effort maximal admissible par ouvrier : $F_{\max} = \SI{150}{N}$.
\item Rapport charge/effort maximal : $\dfrac{R}{F_{\max}} = \dfrac{500}{150} \approx 3{,}33$.
\end{itemize}
\emph{Interprétation :} La machine simple doit avoir un \cle{rapport de démultiplication} (ou avantage mécanique) \textbf{supérieur ou égal à 3,33} pour que l'effort $F$ reste $\le \SI{150}{N}$. C'est le critère principal de sélection.

\textbf{3. Étude du plan incliné (Dispositifs A et A').}
\emph{Relation :} $F \times L = R \times h \quad \Leftrightarrow \quad F = \dfrac{R \times h}{L}$.

\begin{itemize}
\item \textbf{Vérification Dispositif A ($L = \SI{6}{m}$) :}
\[ F_A = \frac{500 \times 3}{6} = \frac{1500}{6} = \SI{250}{N}. \]
La mesure (\SI{250}{N}) est \textbf{cohérente}. Mais $250 > 150$ : \textbf{Dispositif A REFUSÉ} (danger lombaire).
\item \textbf{Recherche de $L_{\min}$ pour $F = \SI{150}{N}$ :}
\[ L_{\min} = \frac{R \times h}{F_{\max}} = \frac{500 \times 3}{150} = \frac{1500}{150} = \SI{10}{m}. \]
Il faut une planche d'au moins \SI{10}{m} de long.
\item \textbf{Vérification Dispositif A' ($L = \SI{12}{m}$) :}
\[ F_{A'} = \frac{500 \times 3}{12} = \frac{1500}{12} = \SI{125}{N}. \]
La mesure (\SI{125}{N}) est \textbf{confirmée}. $125 < 150$ : \textbf{Dispositif A' ACCEPTABLE}.
\item \textbf{Distance :} Pour monter de \SI{3}{m}, l'ouvrier tire le sac sur \SI{12}{m} (ou \SI{6}{m} pour A). Travail : $W = 125 \times 12 = \SI{1500}{J} = 500 \times 3$. Règle d'or respectée.
\end{itemize}

\textbf{4. Étude de la poulie mobile (Dispositifs B et B').}
\emph{Principe :} Le poids $R$ est soutenu par \textbf{deux brins} de corde verticaux. La tension $F$ est la même dans tout le fil (poulie idéale, sans frottement). Équilibre de la poulie mobile : $2F = R \Rightarrow F = R/2$.

\begin{itemize}
\item \textbf{Vérification Dispositif B (1 ouvrier) :}
\[ F_B = \frac{R}{2} = \frac{500}{2} = \SI{250}{N}. \]
La mesure (\SI{250}{N}) est \textbf{cohérente}. Mais $250 > 150$ : \textbf{Dispositif B REFUSÉ} pour un seul ouvrier.
\item \textbf{Dispositif B' (2 ouvriers) :} Les deux ouvriers tirent sur le même brin libre (ou sur deux brins si poulie double, ici on suppose traction commune). L'effort total $F_{\text{total}} = \SI{250}{N}$ est partagé : $F_{\text{par ouvrier}} = \SI{125}{N}$.
$125 < 150$ : \textbf{Dispositif B' ACCEPTABLE}.
\item \textbf{Distance :} Pour monter le sac de $h = \SI{3}{m}$, les deux brins se raccourcissent de \SI{3}{m} chacun $\Rightarrow$ il faut tirer \SI{6}{m} de corde. Travail : $125 \times 6 = \SI{750}{J}$ par ouvrier, total \SI{1500}{J}. Règle d'or respectée.
\end{itemize}

\textbf{5. Étude du treuil (Dispositif C).}
\emph{Modélisation :} Le treuil est un levier de 1ère espèce en rotation. Le point d'appui est l'axe du cylindre.
\begin{itemize}
\item Bras de levier de la force motrice : $d_1 = \text{longueur manivelle} = \SI{40}{cm}$.
\item Bras de levier de la résistance : $d_2 = \text{rayon du tambour } r = \SI{10}{cm}$.
\item \emph{Attention :} Les unités sont homogènes (cm), le rapport est sans dimension.
\end{itemize}

\emph{Relation d'équilibre :} $F \times d_1 = R \times r$.

\begin{align*}
F_C \times 40 &= 500 \times 10 \\
F_C &= \frac{500 \times 10}{40} = \frac{5000}{40} = \SI{125}{N}.
\end{align*}

$125 < 150$ : \textbf{Dispositif C ACCEPTABLE}. Un \textbf{seul} ouvrier suffit.

\emph{Vérification distance (optionnel mais instructif) :}
\begin{itemize}
\item 1 tour de manivelle $\Rightarrow$ force parcourt $2\pi d_1 \approx 2 \times 3,14 \times 40 \approx \SI{251}{cm} = \SI{2,51}{m}$.
\item Corde enroulée : $2\pi r \approx 2 \times 3,14 \times 10 \approx \SI{62,8}{cm} = \SI{0,628}{m}$.
\item Tours nécessaires pour $h = \SI{3}{m}$ : $N = \dfrac{3}{0,628} \approx 4,78 \approx 5$ tours.
\item Distance totale main : $5 \times 2,51 \approx \SI{12,55}{m}$.
\item Travail : $125 \times 12,55 \approx \SI{1569}{J} \approx \SI{1500}{J}$ (arrondis). Règle d'or respectée.
\end{itemize}

\textbf{6. Tableau comparatif des solutions valables (A', B', C).}

\begin{center}
\begin{tabularx}{\linewidth}{|l|X|X|X|}
\hline
\rowcolor{popblueL} \textbf{Critère de comparaison} & \textbf{A' : Plan incliné \SI{12}{m}} & \textbf{B' : Poulie mobile + 2 ouvriers} & \textbf{C : Treuil (manivelle 40 cm)} \\
\hline
Effort par ouvrier & \SI{125}{N} & \SI{125}{N} & \SI{125}{N} \\
\hline
Distance parcourue par la force & \SI{12}{m} (poussée/glissement) & \SI{6}{m} de corde tirée (verticale) & $\approx \SI{12,5}{m}$ (trajectoire circulaire main) \\
\hline
Main-d'œuvre nécessaire & 1 ouvrier & \textbf{2 ouvriers} obligatoires & 1 ouvrier \\
\hline
Sécurité (risques) & Sac qui peut glisser/retomber ; planche glissante (pluie) ; effort horizontal/diagonal fatigant & Corde tirée au-dessus de la tête ; risque si corde casse ; coordination 2 personnes & \textbf{Charge guidée verticalement} ; effort vertical descendant (naturel) ; manivelle controllable ; \textbf{meilleure stabilité} \\
\hline
Rapidité (temps de montée) & Lent (frottements, glissement) & Rapide si coordination bonne & Régulier, contrôlable \\
\hline
Coût / Installation & \textbf{Très faible} (planche bois dispo) & Faible (poulie + corde) mais point d'ancrage haut solide requis & Plus élevé (achat treuil métallique) ; fixation au plancher nécessaire \\
\hline
Encombrement au sol & \textbf{Grand} (\SI{12}{m} de long) & Faible (vertical) & Faible (au plancher) \\
\hline
\end{tabularx}
\end{center}

\textbf{7. Conclusion rédigée (modèle pour le chef de chantier).}

\emph{« Monsieur le Chef de chantier,}
\emph{Les trois dispositifs A' (planche \SI{12}{m}), B' (poulie mobile à 2 ouvriers) et C (treuil) ramènent l'effort à \SI{125}{N}, soit en dessous de la limite de sécurité de \SI{150}{N}. Nous vous recommandons le \textbf{treuil (Dispositif C)} pour deux raisons principales :}
\emph{1) Il ne nécessite qu'\textbf{un seul ouvrier} (calcul : $F = \frac{500 \times 10}{40} = \SI{125}{N}$), ce qui libère de la main-d'œuvre pour d'autres tâches.}
\emph{2) Il offre la \textbf{meilleure sécurité} : la charge monte verticalement de façon guidée et stable, sans risque de glissement sur une pente mouillée ni de corde tendue au-dessus des têtes.}
\emph{Si le budget ne permet pas l'achat immédiat d'un treuil, la planche de \SI{12}{m} (Dispositif A') est une solution de rechange acceptable, gratuite et simple à mettre en œuvre, à condition de bien la fixer et de prévoir des tasseaux antidérapants. »}

\textbf{8. Consignes de sécurité à afficher sur le chantier (obligatoires).}
\begin{enumerate}
\item \textbf{Vérification systématique :} Avant chaque levage, contrôler l'état des cordes (pas d'effilochage), de la manivelle, des axes, des fixations du treuil ou de la poulie, et de la planche (solidité, propreté).
\item \textbf{Zone interdite :} \textbf{Ne jamais se placer sous la charge suspendue} (sac en l'air) ni dans l'axe de la corde tendue. Port du \textbf{casque de chantier} obligatoire pour tous les présents.
\item \textbf{Geste professionnel :} Soulever \textbf{lentement et à vitesse constante}, sans à-coups. Ne \textbf{jamais lâcher la manivelle} ou la corde en cours de montée (risque de chute libre du sac). Utiliser le frein/rochet du treuil pour toute pause.
\item \textbf{Limites de charge et personnel :} Ne pas dépasser \SI{50}{kg} par sac (poids nominal CIMENCAM). Interdire l'accès de la zone de levage aux élèves et aux personnes non autorisées. Un seul ouvrier sur la manivelle (treuil) ; coordination stricte si poulie à deux.
\end{enumerate}

\courssub{Bilan de la leçon : Ce qu'il faut retenir pour le BEPC}

\begin{itemize}
\item Une \cle{machine simple} ne supprime pas le travail à fournir : elle \textbf{échange de la force contre de la distance}. Pour un même sac de \SI{500}{N} monté de \SI{3}{m}, le travail utile est toujours $W = R \times h = \SI{1500}{J}$, quelle que soit la machine (planche, poulie, treuil).
\item Chaque machine a son propre \cle{rapport de démultiplication} (avantages mécanique) :
  \begin{itemize}
  \item Plan incliné : $\dfrac{L}{h}$ (plus la pente est longue, moins l'effort est grand).
  \item Treuil (levier rotatif) : $\dfrac{d_1}{r} = \dfrac{\text{longueur manivelle}}{\text{rayon tambour}}$.
  \item Poulie mobile : 2 (fixe par construction simple).
  \end{itemize}
  On peut donc \textbf{calculer à l'avance} l'effort nécessaire et vérifier qu'il reste sous une limite de sécurité ($F \le F_{\max}$).
\item Le choix technique ne dépend pas seulement de l'effort : il faut comparer \cle{distance parcourue}, \cle{rapidité}, \cle{coût}, \cle{main-d'œuvre}, \cle{encombrement} et surtout \cle{sécurité} (stabilité de la charge, posture de l'opérateur).
\item Les formules $F \times d_1 = R \times d_2$ et $F \times L = R \times h$ sont des outils de \textbf{prévision} (dimensionner : $L \ge \SI{10}{m}$) et de \textbf{validation} de mesures expérimentales.
\end{itemize}

\end{exoresolu}

\begin{attention}[Erreurs fréquentes à éviter au BEPC]
\begin{itemize}
\item \textbf{Unités non homogènes :} Ne pas convertir les longueurs dans la même unité avant d'appliquer $F \times d_1 = R \times d_2$ (ex : $d_1 = \SI{40}{cm}$ et $r = \SI{0,1}{m}$ $\rightarrow$ erreur). \emph{Astuce :} gardez les cm si les deux sont en cm, le rapport est sans dimension.
\item \textbf{Confusion Masse / Poids :} La charge est une masse (\SI{50}{kg}), la résistance est un poids (\SI{500}{N}). Dans les formules des machines simples, on utilise \textbf{toujours les forces en newtons (N)}.
\item \textbf{Distance poulie mobile :} Oublier qu'avec une poulie mobile, la longueur de corde à tirer est le \textbf{double} de la hauteur gagnée ($d_F = 2h$).
\item \textbf{Efficacité énergétique :} Croire que la machine « crée » de la force ou de l'énergie. $W_{\text{moteur}} = W_{\text{résistant}}$ (frottements négligés). Si frottements, $W_{\text{moteur}} > W_{\text{résistant}}$ (rendement $< 100\%$).
\item \textbf{Sens de la force sur plan incliné :} La formule $F \times L = R \times h$ suppose que la force $F$ est \textbf{parallèle à la pente}. Si on tire horizontalement, la formule change (non au programme de 3ème).
\end{itemize}
\end{attention}

\newpage

\coursec{Méthodes et exercices résolus}

\coursec{Leçon 11 : Quelques machines simples}

\begin{definition}[Machine simple]
Une \cle{machine simple} est un dispositif mécanique élémentaire, sans moteur, qui permet de modifier l'intensité, la direction ou le point d'application d'une force (l'\cle{effort} $F$) pour vaincre une autre force (la \cle{résistance} $R$, souvent le poids d'une charge). Les quatre machines simples au programme sont : le \cle{levier}, le \cle{plan incliné}, la \cle{poulie} et le \cle{treuil}.
\end{definition}

\begin{popfigure}
\begin{tikzpicture}[scale=0.9, >=Stealth]
% Levier
\draw[popdark, thick] (-2.5,-0.2) -- (2.5,-0.2) node[midway, below=2mm] {Barre rigide};
\draw[popblue, thick] (0,0) -- (0,0.6) node[above] {$O$ (Point d'appui)};
\draw[poporange, thick, ->] (-2,-0.2) -- (-2,-1) node[below] {$R$ (Résistance)};
\draw[popgreen, thick, ->] (2,-0.2) -- (2,-1) node[below] {$F$ (Effort)};
\draw[dashed, popmuted] (0,-0.2) -- (0,-1.2);
\draw[<->, popdark] (0,-0.8) -- (-2,-0.8) node[midway, below] {$d_R$};
\draw[<->, popdark] (0,-0.8) -- (2,-0.8) node[midway, below] {$d_F$};
% Plan incliné
\begin{scope}[xshift=7cm]
\draw[popdark, thick] (0,0) -- (3,1.5) -- (3,1.2) -- (0,-0.3) -- cycle;
\draw[poporange, thick, ->] (1.5,1.35) -- (1.5,2.2) node[above] {$P$};
\draw[popgreen, thick, ->] (1.5,1.35) -- (2.5,2.0) node[right] {$F$};
\draw[dashed, popmuted] (1.5,1.35) -- (1.5,0);
\draw[<->, popdark] (3.2,0) -- (3.2,1.5) node[midway, right] {$h$};
\draw[<->, popdark] (0,-0.5) -- (3,-0.5) node[midway, below] {$L$};
\end{scope}
% Poulie fixe
\begin{scope}[yshift=-4.5cm]
\draw[popdark, thick] (-1,1) -- (1,1); % Support
\draw[popblue, thick] (0,1) circle (0.5); % Poulie
\draw[poporange, thick] (0,0.5) -- (0,-1.5) node[below] {$P$};
\draw[popgreen, thick, ->] (0.5,1) arc (0:-90:0.5) -- (0.5,-1.5) node[below] {$F$};
\node at (0,1.7) {Poulie fixe};
\end{scope}
% Poulie mobile
\begin{scope}[xshift=7cm, yshift=-4.5cm]
\draw[popdark, thick] (-1,1) -- (1,1); % Support
\draw[popblue, thick] (0,0.2) circle (0.5); % Poulie mobile
\draw[poporange, thick] (0,-0.3) -- (0,-1.5) node[below] {$P$};
\draw[popgreen, thick] (0,1) -- (0,0.7); % Brin fixe
\draw[popgreen, thick, ->] (0,-0.7) -- (0,-1.5) node[below] {$F$};
\draw[<->, popmuted] (-0.8,0.2) -- (-0.8,-0.3) node[midway, left] {$2$ brins};
\node at (0,1.7) {Poulie mobile};
\end{scope}
% Treuil
\begin{scope}[yshift=-9cm]
\draw[popblue, thick] (0,0) circle (0.6); % Tambour
\draw[popdark, thick] (0,0) -- (1.5,0) -- (1.5,-0.2) -- (0,-0.2) -- cycle; % Manivelle simplifiée
\draw[popgreen, thick, ->] (1.5,0) -- (2.2,0) node[right] {$F$};
\draw[poporange, thick] (0,-0.6) -- (0,-2) node[below] {$P$};
\draw[<->, popdark] (0.7,0) -- (0.7,-0.6) node[midway, right] {$r$};
\draw[<->, popdark] (0,0.2) -- (1.5,0.2) node[midway, above] {$L$};
\node at (0,1) {Treuil};
\end{scope}
\end{tikzpicture}
\legende{Synthèse visuelle des quatre machines simples : levier, plan incliné, poulie (fixe et mobile), treuil.}
\end{popfigure}

\courssub{Méthode 1 : Identifier et décrire une machine simple}

\begin{methode}[Reconnaître une machine simple dans une situation concrète]
Pour identifier une machine simple dans une situation de la vie courante, procède toujours ainsi :
\begin{enumerate}
\item Repère l'\cle{effort} (force motrice $F$) exercé par l'utilisateur et la \cle{résistance} (force résistante $R$, souvent le poids de la charge).
\item Cherche le dispositif qui transmet ou transforme cet effort : barre rigide tournant autour d'un point (levier), surface en pente (plan incliné), roue à gorge avec corde (poulie), cylindre + manivelle (treuil).
\item Identifie les éléments caractéristiques : point d'appui, bras de levier, hauteur et longueur du plan, nombre de brins de corde, rayon du tambour et longueur de la manivelle.
\item Conclus en nommant la machine et en précisant son rôle : \emph{démultiplier} l'effort ou \emph{changer la direction} de la force.
\end{enumerate}
Réflexe de rédaction : au BEPC, cite toujours les trois éléments d'un levier (point d'appui, effort, résistance) pour gagner tous les points.
\end{methode}

\begin{experience}[Vérification de la règle du levier]
\textbf{Matériel :} Règle graduée (50 cm), pivot (crayon ou triangle), deux porte-masses (10 g et 50 g), masses (10 g, 20 g, 50 g).
\textbf{Protocole :}
\begin{enumerate}
\item Place la règle sur le pivot à la graduation 25 cm (point d'appui $O$).
\item Accroche la masse de 50 g (résistance $R$) à la graduation 15 cm ($d_R = 10$ cm).
\item Équilibre la règle en plaçant la masse de 10 g (effort $F$) du côté opposé. Note sa position $d_F$.
\item Vérifie l'égalité $F \times d_F = R \times d_R$ (avec $F = m_F g$, $R = m_R g$).
\item Répète en changeant $d_R$ (ex. 5 cm, 20 cm).
\end{enumerate}
\textbf{Observations :} Le produit $m \times d$ (proportionnel au moment) est constant à l'équilibre.
\textbf{Interprétation :} La condition d'équilibre $F \times d_F = R \times d_R$ est vérifiée expérimentalement. Plus $d_F$ est grand, plus la masse $m_F$ nécessaire est petite.
\textbf{Conclusion :} Le levier permet de soulever une charge lourde avec un effort faible si le bras de l'effort est plus long que celui de la résistance.
\end{experience}

\begin{exoresolu}[Le déménageur et le rocher]
Pour déplacer un gros rocher devant sa maison à Yaoundé, un maçon glisse une barre de fer sous le rocher et cale une pierre sous la barre, près du rocher. Il appuie ensuite sur l'extrémité libre de la barre.
\begin{enumerate}
\item Quelle machine simple utilise-t-il ?
\item Nomme les trois éléments caractéristiques de cette machine dans cette situation.
\item Pourquoi la pierre est-elle placée près du rocher ?
\end{enumerate}
\tcblower
\textbf{1. Identification de la machine.} La barre de fer est une tige rigide qui tourne autour d'un point fixe (la pierre) sous l'action de deux forces : c'est un \cle{levier}.

\textbf{2. Les trois éléments caractéristiques.}
\begin{itemize}
\item Le \cle{point d'appui} : la pierre calée sous la barre ;
\item la \cle{résistance} $R$ : le poids du rocher, qui s'exerce sur la barre ;
\item l'\cle{effort} $F$ : la force exercée par le maçon sur l'extrémité libre.
\end{itemize}

\textbf{3. Position de la pierre.} La pierre est placée près du rocher pour que le bras de levier de l'effort (distance pierre--main) soit \emph{grand} et celui de la résistance (distance pierre--rocher) soit \emph{petit}. D'après la condition d'équilibre du levier $F \times d_F = R \times d_R$, plus $d_F$ est grand devant $d_R$, plus l'effort $F$ nécessaire est faible : le levier \emph{démultiplie} la force du maçon.

\textbf{Conclusion :} le maçon utilise un levier dont le point d'appui est la pierre ; placer celle-ci près du rocher réduit l'effort à fournir.
\end{exoresolu}

\courssub{Méthode 2 : Exploiter l'équilibre d'un levier}

\begin{methode}[Calculer avec la condition d'équilibre du levier]
Un levier est en équilibre quand les moments de l'effort et de la résistance par rapport au point d'appui sont égaux :
\[ F \times d_F = R \times d_R \]
où $d_F$ est le bras de levier de l'effort et $d_R$ celui de la résistance (distances mesurées depuis le point d'appui).
\begin{enumerate}
\item Fais un schéma du levier : place le point d'appui $O$, l'effort $F$, la résistance $R$.
\item Mesure ou calcule les bras de levier $d_F$ et $d_R$ \emph{à partir du point d'appui} (convertis tout en mètres).
\item Écris l'égalité des moments : $F \times d_F = R \times d_R$.
\item Isole la grandeur inconnue : $F = \dfrac{R \times d_R}{d_F}$.
\item Vérifie la cohérence : si $d_F > d_R$, alors $F < R$ (levier avantageux).
\end{enumerate}
\end{methode}

\begin{exoresolu}[La brouette du chantier]
Une brouette chargée de sable constitue un levier dont le point d'appui est l'axe de la roue. La charge a un poids $R = \SI{600}{N}$ appliqué à $d_R = \SI{0,40}{m}$ de l'axe. Les poignées sont à $d_F = \SI{1,20}{m}$ de l'axe.
\begin{enumerate}
\item Calcule l'effort $F$ que le maçon doit exercer pour soulever les poignées.
\item Compare $F$ et $R$ et conclus sur l'intérêt de la brouette.
\end{enumerate}
\tcblower
\textbf{1. Calcul de l'effort.} La brouette est un levier en équilibre autour de l'axe de la roue. La condition d'équilibre s'écrit :
\[ F \times d_F = R \times d_R \]
On isole $F$ :
\[ F = \frac{R \times d_R}{d_F} = \frac{600 \times 0{,}40}{1{,}20} = \frac{240}{1{,}20} = \SI{200}{N} \]

\textbf{2. Comparaison.} $F = \SI{200}{N} < R = \SI{600}{N}$ : l'effort est trois fois plus petit que la charge car le bras de levier de l'effort ($\SI{1,20}{m}$) est trois fois plus grand que celui de la résistance ($\SI{0,40}{m}$).

\textbf{Conclusion :} le maçon n'exerce que $\SI{200}{N}$ pour soulever une charge de $\SI{600}{N}$ : la brouette est un levier qui démultiplie l'effort, d'où son intérêt sur les chantiers.
\end{exoresolu}

\courssub{Méthode 3 : Exploiter le plan incliné}

\begin{methode}[Calculer l'effort sur un plan incliné]
Pour monter une charge de poids $P$ d'une hauteur $h$ le long d'un plan incliné de longueur $L$ (frottements négligés) :
\[ F = \frac{P \times h}{L} \]
\begin{enumerate}
\item Repère la hauteur verticale $h$ (ce qu'on veut gagner) et la longueur $L$ du plan (le chemin parcouru).
\item Convertis toutes les longueurs dans la même unité (le mètre).
\item Applique la formule $F = \dfrac{P \times h}{L}$.
\item Compare avec le poids : comme $L > h$, on a toujours $F < P$ : le plan incliné réduit l'effort, mais augmente la distance parcourue.
\end{enumerate}
Réflexe : plus le plan est long et peu pentu, plus l'effort est faible (rampe d'accès, route en lacets dans les montagnes de l'Ouest).
\end{methode}

\begin{experience}[Effort sur plan incliné : frottements et angle]
\textbf{Matériel :} Plan incliné réglable, chariot ou bloc en bois, dynamomètre, masses, rapporteur, règle.
\textbf{Protocole :}
\begin{enumerate}
\item Fixe le plan à un angle $\alpha = 30^\circ$. Mesure $h$ et $L$.
\item Attache le dynamomètre au chariot (masse $m = 500$ g, $P \approx \SI{5}{N}$) et tire parallèlement à la pente à vitesse constante. Note $F_{\text{mes}}$.
\item Calcule $F_{\text{th}} = P \times h / L$.
\item Compare $F_{\text{mes}}$ et $F_{\text{th}}$. Recommence pour $\alpha = 15^\circ$ et $45^\circ$.
\end{enumerate}
\textbf{Observations :} $F_{\text{mes}} > F_{\text{th}}$ à cause des frottements. L'écart augmente avec l'angle.
\textbf{Interprétation :} La formule $F = P h / L$ suppose un plan \emph{idéal} (sans frottements). En réel, il faut ajouter la force de frottement.
\textbf{Conclusion :} Le plan incliné réduit l'effort théorique, mais les frottements limitent ce gain ; il faut lubrifier ou utiliser des rouleaux (brouette, chariot).
\end{experience}

\begin{exoresolu}[Chargement du camion à Douala]
Pour charger un fût d'huile de poids $P = \SI{900}{N}$ dans un camion dont le plateau est à $h = \SI{1,2}{m}$ du sol, des dockers utilisent une planche de longueur $L = \SI{3}{m}$ posée en rampe (frottements négligés).
\begin{enumerate}
\item Calcule l'effort $F$ nécessaire pour faire monter le fût le long de la planche.
\item Compare avec l'effort qu'il faudrait pour soulever le fût verticalement.
\end{enumerate}
\tcblower
\textbf{1. Effort sur le plan incliné.} La planche est un plan incliné de hauteur $h = \SI{1,2}{m}$ et de longueur $L = \SI{3}{m}$. L'équilibre des travaux des forces donne :
\[ F = \frac{P \times h}{L} = \frac{900 \times 1{,}2}{3} = \frac{1080}{3} = \SI{360}{N} \]

\textbf{2. Comparaison.} Soulever le fût verticalement exigerait un effort égal au poids, soit $\SI{900}{N}$. Avec la rampe, l'effort n'est que de $\SI{360}{N}$, soit $2{,}5$ fois moins (car $\dfrac{L}{h} = \dfrac{3}{1{,}2} = 2{,}5$).

\textbf{Conclusion :} le plan incliné réduit l'effort à fournir ($\SI{360}{N}$ au lieu de $\SI{900}{N}$), au prix d'un déplacement plus long ($\SI{3}{m}$ au lieu de $\SI{1,2}{m}$) : on gagne en force ce qu'on perd en distance.
\end{exoresolu}

\courssub{Méthode 4 : Étudier une poulie}

\begin{methode}[Distinguer poulie fixe et poulie mobile]
\begin{enumerate}
\item \cle{Poulie fixe} (accrochée à un support) : elle \emph{change la direction} de la force mais ne la réduit pas : $F = P$. Tirer vers le bas est plus commode et plus sûr que soulever.
\item \cle{Poulie mobile} (solidaire de la charge) : la charge est soutenue par \emph{deux brins} de corde, donc l'effort est divisé par deux : $F = \dfrac{P}{2}$. En revanche, il faut tirer une longueur de corde double du déplacement de la charge.
\item Pour compter : l'effort est le poids divisé par le nombre $n$ de brins qui soutiennent la charge : $F = \dfrac{P}{n}$.
\end{enumerate}
Réflexe : vérifie si la poulie est fixée au plafond (fixe) ou attachée à la charge (mobile) avant tout calcul.
\end{methode}

\begin{experience}[Comparaison poulie fixe / poulie mobile]
\textbf{Matériel :} Deux poulies, support, corde, masse de 200 g ($P \approx \SI{2}{N}$), dynamomètre.
\textbf{Protocole :}
\begin{enumerate}
\item Monte une poulie fixe. Accroche la masse. Tire au dynamomètre pour lever la masse à vitesse constante. Note $F_1$.
\item Monte une poulie mobile (corde fixée au support, passant dans la poulie mobile attachée à la masse, effort vers le haut). Tire pour lever la masse. Note $F_2$.
\item Mesure la longueur de corde tirée $\Delta l$ pour un déplacement $\Delta h = \SI{10}{cm}$ de la masse dans les deux cas.
\end{enumerate}
\textbf{Observations :} $F_1 \approx \SI{2}{N}$ ; $F_2 \approx \SI{1}{N}$. Pour $\Delta h = \SI{10}{cm}$, $\Delta l_1 = \SI{10}{cm}$, $\Delta l_2 = \SI{20}{cm}$.
\textbf{Interprétation :} Poulie fixe : $F=P$, $\Delta l = \Delta h$. Poulie mobile : $F=P/2$, $\Delta l = 2 \Delta h$. Le travail $F \times \Delta l = P \times \Delta h$ est conservé (frottements négligés).
\textbf{Conclusion :} La poulie mobile démultiplie l'effort (gain de force) mais double la course de la corde (perte de distance).
\end{experience}

\begin{exoresolu}[Le seau du puits et le palan du garage]
\begin{enumerate}
\item Au village, Amina remonte un seau d'eau de poids $P = \SI{120}{N}$ à l'aide d'une poulie fixe au-dessus du puits. Quel effort doit-elle exercer ? Quel est l'avantage de cette poulie ?
\item Dans un garage de Bafoussam, un mécanicien soulève un moteur de poids $P = \SI{800}{N}$ avec une poulie mobile. Quel effort exerce-t-il ? De quelle longueur de corde doit-il tirer pour élever le moteur de $\SI{0,5}{m}$ ?
\end{enumerate}
\tcblower
\textbf{1. Poulie fixe.} Une poulie fixe ne démultiplie pas l'effort : $F = P = \SI{120}{N}$. Son avantage est de \emph{changer la direction} de la force : Amina tire la corde vers le bas (ou horizontalement), position plus commode et plus sûre que de soulever le seau au-dessus du puits.

\textbf{2. Poulie mobile.} La charge est soutenue par deux brins de corde :
\[ F = \frac{P}{2} = \frac{800}{2} = \SI{400}{N} \]
L'effort est divisé par deux. Mais la conservation du travail impose que la longueur de corde tirée soit double du déplacement de la charge :
\[ L = 2 \times 0{,}5 = \SI{1}{m} \]

\textbf{Conclusion :} la poulie fixe (effort $\SI{120}{N}$) rend le geste plus commode ; la poulie mobile réduit l'effort à $\SI{400}{N}$ mais oblige à tirer $\SI{1}{m}$ de corde pour élever le moteur de $\SI{0,5}{m}$.
\end{exoresolu}

\courssub{Méthode 5 : Étudier un treuil}

\begin{methode}[Calculer l'effort d'un treuil]
Un treuil est formé d'un tambour (cylindre) de rayon $r$ sur lequel s'enroule une corde, et d'une manivelle de longueur $L$ sur laquelle on exerce l'effort $F$. À l'équilibre, les moments sont égaux :
\[ F \times L = P \times r \qquad \text{donc} \qquad F = \frac{P \times r}{L} \]
\begin{enumerate}
\item Repère le rayon $r$ du tambour et la longueur $L$ de la manivelle (bras de levier de l'effort).
\item Convertis les longueurs dans la même unité.
\item Applique $F = \dfrac{P \times r}{L}$.
\item Conclus : comme $L > r$, l'effort est plus petit que la charge ; plus la manivelle est longue, plus l'effort est faible.
\end{enumerate}
\end{methode}

\begin{exoresolu}[Le treuil du puits de Maroua]
Pour puiser l'eau d'un puits, on utilise un treuil : le tambour a un rayon $r = \SI{10}{cm}$ et la manivelle a une longueur $L = \SI{40}{cm}$. Le seau plein pèse $P = \SI{150}{N}$.
\begin{enumerate}
\item Calcule l'effort $F$ à exercer sur la manivelle.
\item Que devient cet effort si on remplace la manivelle par une manivelle deux fois plus longue ?
\end{enumerate}
\tcblower
\textbf{1. Calcul de l'effort.} À l'équilibre du treuil, les moments par rapport à l'axe sont égaux :
\[ F \times L = P \times r \quad \Longrightarrow \quad F = \frac{P \times r}{L} = \frac{150 \times 10}{40} = \frac{1500}{40} = \SI{37,5}{N} \]
(Les unités cm se simplifient car c'est un rapport de longueurs).

\textbf{2. Manivelle deux fois plus longue.} Avec $L' = 2 \times 40 = \SI{80}{cm}$ :
\[ F' = \frac{P \times r}{L'} = \frac{150 \times 10}{80} = \SI{18,75}{N} \]
L'effort est divisé par deux : il est proportionnel à l'inverse de la longueur de la manivelle.

\textbf{Conclusion :} l'effort vaut $\SI{37,5}{N}$ avec la manivelle de $\SI{40}{cm}$ et seulement $\SI{18,75}{N}$ avec une manivelle double : allonger la manivelle diminue l'effort, mais augmente le nombre de tours à effectuer.
\end{exoresolu}

\courssub{Méthode 6 : Comparer les machines simples et appliquer les consignes de sécurité}

\begin{methode}[Choisir la machine adaptée et travailler en sécurité]
\begin{enumerate}
\item Définis le besoin : soulever verticalement (poulie, treuil), faire rouler ou glisser vers le haut (plan incliné), soulever ou écarter par rotation (levier).
\item Compare effort et déplacement : toute machine simple qui réduit l'effort \emph{augmente} le déplacement (le travail ne change pas, frottements négligés).
\item Rappelle les inconvénients : usure, frottements, encombrement, longueur de la manoeuvre.
\item Cite les consignes de sécurité : vérifier l'état des cordes, câbles et barres ; ne jamais dépasser la charge maximale ; ne jamais se placer sous une charge suspendue ; porter des gants ; caler les charges sur un plan incliné.
\end{enumerate}
\end{methode}

\begin{propriete}[Règle d'or des machines simples (Travail)]
Dans une machine simple idéale (sans frottements), le travail de l'effort (moteur) est égal au travail de la résistance (résistant) :
\[ W_F = W_R \quad \Longleftrightarrow \quad F \times \Delta l_F = R \times \Delta l_R \]
\textbf{Conséquence :} Si la machine démultiplie la force ($F < R$), elle multiplie le déplacement ($\Delta l_F > \Delta l_R$). \emph{On gagne en force ce qu'on perd en distance.}
\end{propriete}

\begin{exoresolu}[Le chantier de l'immeuble à Bonabéri]
Sur un chantier, un chef d'équipe doit : (a) soulever des sacs de ciment de $\SI{500}{N}$ jusqu'au deuxième étage ; (b) faire monter une brouette lourde sur une dalle surélevée de $\SI{50}{cm}$ ; (c) décoincer une poutre coincée dans la terre.
\begin{enumerate}
\item Propose une machine simple adaptée à chaque tâche en justifiant.
\item Cite deux consignes de sécurité à respecter pour la tâche (a).
\end{enumerate}
\tcblower
\textbf{1. Choix des machines.}
\begin{itemize}
\item (a) \emph{Sacs de ciment} : une \cle{poulie} (fixe ou mobile) fixée en hauteur. Elle permet de soulever verticalement ; avec une poulie mobile, l'effort est divisé par deux ($F = \SI{250}{N}$).
\item (b) \emph{Brouette} : un \cle{plan incliné} (planche en rampe). La brouette roule : il suffit d'un effort $F = \dfrac{P \times h}{L}$, d'autant plus faible que la planche est longue.
\item (c) \emph{Poutre coincée} : un \cle{levier} (barre de fer avec point d'appui proche de la poutre) pour démultiplier l'effort de soulèvement.
\end{itemize}

\textbf{2. Sécurité pour la tâche (a).} Par exemple : vérifier que la corde et le crochet sont en bon état et adaptés à la charge de $\SI{500}{N}$ ; interdire le passage et la présence de toute personne sous la charge suspendue pendant la montée.

\textbf{Conclusion :} chaque machine simple répond à un besoin précis ; leur utilisation exige le respect de consignes de sécurité pour éviter les accidents de chantier.
\end{exoresolu}

\begin{savaistu}[Machines simples au Cameroun : tradition et modernité]
\begin{itemize}
\item \textbf{Le levier ancestral :} Dans les chefferies de l'Ouest (Bamoun, Bamiléké), les portes lourdes des cases de notables s'ouvrent grâce à un système de levier en bois (pivot en bas, effort en haut) : un enfant peut ouvrir une porte de 200 kg.
\item \textbf{Les routes en lacets :} La route qui monte à Dschang ou à Bafoussam utilise le principe du plan incliné : on allonge le tracé (grande $L$) pour réduire la pente ($h/L$ petit), permettant aux camions chargés de monter sans patiner.
\item \textbf{Le treuil du forage :} Dans le Nord (Maroua, Garoua), les forages à pompe manuelle (type India Mark II) sont des treuils modernisés : grande manivelle ($L \approx \SI{60}{cm}$), petit tambour ($r \approx \SI{5}{cm}$), gain de force $\approx 12$ fois.
\item \textbf{Le palan de chantier :} À Douala ou Yaoundé, les grues de chantier sont des systèmes de poulies multiples (palans) : 4, 6 ou 8 brins porteurs pour soulever des tonnes avec un moteur électrique d'effort modéré.
\end{itemize}
\end{savaistu}

\begin{aretenir}[Formules à connaître pour le BEPC]
\begin{center}
\begin{tabularx}{\linewidth}{|l|X|X|}\hline
\textbf{Machine} & \textbf{Condition d'équilibre (idéal)} & \textbf{Règle pratique} \\ \hline
Levier & $F \times d_F = R \times d_R$ & $d_F, d_R$ mesurés depuis le point d'appui \\ \hline
Plan incliné & $F = \dfrac{P \times h}{L}$ & $h$ = hauteur verticale, $L$ = longueur pente \\ \hline
Poulie fixe & $F = P$ & Change la direction seulement \\ \hline
Poulie mobile ($n$ brins) & $F = \dfrac{P}{n}$ & $n$ = nombre de brins soutenant la charge \\ \hline
Treuil & $F \times L = P \times r$ $\Leftrightarrow$ $F = \dfrac{P \times r}{L}$ & $L$ = longueur manivelle, $r$ = rayon tambour \\ \hline
\end{tabularx}
\end{center}
\textbf{Règle d'or :} $F \times \Delta l_F = R \times \Delta l_R$ (Travail conservé). Gain en force = Perte en distance.
\end{aretenir}

\begin{attention}[Les 5 erreurs qui coûtent des points au BEPC]
\begin{enumerate}
\item \textbf{Mesurer les bras de levier depuis le mauvais point.} Les distances $d_F$ et $d_R$ se mesurent toujours \emph{depuis le point d'appui}, jamais depuis le milieu de la barre ni entre l'effort et la résistance.
\item \textbf{Oublier de convertir les unités.} Dans $F = \dfrac{P \times r}{L}$ ou $F \times d_F = R \times d_R$, toutes les longueurs doivent être dans la \emph{même unité} (cm avec cm, m avec m) avant le calcul ; un mélange cm/m fausse tout le résultat.
\item \textbf{Croire que la poulie fixe divise l'effort par deux.} Une poulie \emph{fixe} ne fait que changer la direction de la force ($F = P$) ; seule la poulie \emph{mobile} divise l'effort par le nombre de brins porteurs.
\item \textbf{Confondre hauteur et longueur du plan incliné.} Dans $F = \dfrac{P \times h}{L}$, $h$ est la hauteur \emph{verticale} et $L$ la longueur \emph{le long de la pente} ; les intervertir donne un effort plus grand que le poids, ce qui est absurde : vérifie toujours que $F < P$.
\item \textbf{Oublier la contrepartie du gain en force.} Toute machine simple qui réduit l'effort augmente le déplacement (corde à tirer plus longue, chemin plus long) : le préciser dans la conclusion montre que tu as compris le principe et rapporte le point de synthèse.
\end{enumerate}
\end{attention}

\begin{exercice}[Entraînement BEPC - Machines simples]
\begin{enumerate}
\item \textbf{Le casse-noix.} Un casse-noix est un levier à deux bras. La noix (résistance $R$) est placée à $d_R = \SI{2}{cm}$ du pivot. Les mains (effort $F$) agissent à $d_F = \SI{12}{cm}$. Si on appuie avec $F = \SI{50}{N}$, quelle est la force exercée sur la noix ?
\item \textbf{Rampe PMR.} Une rampe d'accès pour personne à mobilité réduite (PMR) doit respecter une pente maximale de $5\%$ ($h/L = 0,05$). Pour franchir un trottoir de $h = \SI{15}{cm}$, quelle doit être la longueur minimale $L$ de la rampe ? Si un fauteuil roulant pèse $P = \SI{800}{N}$, quel effort $F$ faut-il pour le monter (sans frottement) ?
\item \textbf{Palan de levage.} Un palan à 4 brins porteurs soulève une charge de $P = \SI{2000}{N}$. Calculer l'effort $F$. Si on veut lever la charge de $\SI{2}{m}$, quelle longueur de corde faut-il tirer ?
\item \textbf{Treuil de voiture.} Un treuil de 4x4 a un tambour de rayon $r = \SI{5}{cm}$ et une manivelle (ou rayon de la commande) de $L = \SI{30}{cm}$. Quel effort permet de tracter la voiture ($P = \SI{10000}{N}$) ?
\end{enumerate}
\end{exercice}

\begin{corrige}[Entraînement BEPC - Machines simples]
\begin{enumerate}
\item \textbf{Casse-noix.} $F \times d_F = R \times d_R \Rightarrow R = \dfrac{F \times d_F}{d_R} = \dfrac{50 \times 12}{2} = \SI{300}{N}$. La force sur la noix est 6 fois l'effort des mains.
\item \textbf{Rampe PMR.} Pente $5\% \Rightarrow h/L = 0,05 \Rightarrow L = h / 0,05 = 0,15 / 0,05 = \SI{3}{m}$. Effort $F = P \times h / L = 800 \times 0,15 / 3 = \SI{40}{N}$. Très faible effort grâce à la grande longueur.
\item \textbf{Palan.} $n=4 \Rightarrow F = P/n = 2000/4 = \SI{500}{N}$. Longueur corde $L = n \times \Delta h = 4 \times 2 = \SI{8}{m}$.
\item \textbf{Treuil 4x4.} $F = P \times r / L = 10000 \times 5 / 30 = 50000 / 30 \approx \SI{1667}{N}$. (Note : en réalité, un treuil de 4x4 utilise un réducteur de vitesse / multiplicateur de couple interne, pas seulement la manivelle).
\end{enumerate}
\end{corrige}

\newpage

\coursec{Exercices}

\courssub{Je vérifie mes connaissances}

\begin{exercice}[QCM : la bonne machine simple]
Pour chaque objet ou situation, choisis la machine simple qu'il met en jeu : \textbf{a)} levier ; \textbf{b)} plan incliné ; \textbf{c)} poulie ; \textbf{d)} treuil.
\begin{enumerate}
\item Une brouette pour transporter du sable.
\item Une rampe en planches pour monter une moto sur un camion.
\item Un casse-noix.
\item La manivelle d'un moulin à manioc.
\item Une corde passant sur une roue pour monter un seau d'eau.
\item Un pied-de-biche pour soulever une pierre.
\item Un escalier.
\item Le système qui remonte le store d'une boutique.
\end{enumerate}
\end{exercice}

\begin{exercice}[Vrai ou faux ? Justifie]
Réponds par \emph{vrai} ou \emph{faux}, puis justifie ta réponse en une phrase.
\begin{enumerate}
\item Une machine simple permet de soulever une charge avec un effort plus petit que le poids de la charge.
\item Avec une poulie fixe, l'effort nécessaire est deux fois plus petit que le poids de la charge.
\item Plus le plan incliné est long (pour une même hauteur), plus l'effort à fournir est grand.
\item Dans un levier, plus le bras de levier de l'effort est grand, plus l'effort est petit.
\end{enumerate}
\end{exercice}

\begin{exercice}[Définitions à compléter]
Recopie et complète les phrases suivantes avec les mots : \emph{point d'appui}, \emph{charge}, \emph{effort}, \emph{treuil}, \emph{machine simple}.
\begin{enumerate}
\item Un dispositif qui permet de déplacer ou de soulever une charge en réduisant l'effort à fournir est une \ldots
\item Dans un levier, la barre tourne autour d'un point fixe appelé \ldots
\item La force exercée par l'utilisateur sur la machine s'appelle l'\ldots ; le poids de l'objet à déplacer s'appelle la \ldots
\item Le \ldots est constitué d'un tambour sur lequel s'enroule une corde, actionné par une manivelle.
\end{enumerate}
\end{exercice}

\begin{exercice}[Calculs directs]
\begin{enumerate}
\item Un ouvrier exerce un effort de \SI{150}{N} sur une poulie mobile pour soulever une charge. Quel est le poids maximal de la charge (on néglige les frottements) ?
\item Sur un plan incliné, un élève mesure un effort de \SI{60}{N} pour tirer une caisse de \SI{180}{N}. De combien de newtons l'effort est-il réduit par rapport à une montée verticale ?
\item La manivelle d'un treuil a une longueur de \SI{40}{cm} et le tambour un rayon de \SI{10}{cm}. Quel est le rapport entre l'effort et la charge ?
\end{enumerate}
\end{exercice}

\courssub{Je m'entraîne}

\begin{exercice}[Le levier du chantier]
Un ouvrier utilise une barre de fer de \SI{2,5}{m} pour soulever un bloc de béton de poids \SI{400}{N}. Le point d'appui est placé à \SI{0,5}{m} de la charge.
\begin{enumerate}
\item Fais un schéma légendé du levier (charge, effort, point d'appui, bras de levier).
\item Calcule la longueur du bras de levier de l'effort.
\item En appliquant la condition d'équilibre du levier $F \times d_F = P \times d_P$, calcule l'effort $F$ nécessaire.
\item Conclus sur l'avantage du levier dans cette situation.
\end{enumerate}
\end{exercice}

\begin{exercice}[Le puits du quartier]
Un seau plein d'eau a un poids de \SI{120}{N}. On le remonte d'un puits de \SI{4}{m} de profondeur.
\begin{enumerate}
\item \textbf{Avec une poulie fixe :} quel effort faut-il exercer ? Quelle longueur de corde faut-il tirer pour que le seau monte de \SI{4}{m} ?
\item \textbf{Avec une poulie mobile :} quel effort faut-il exercer (frottements négligés) ? Quelle longueur de corde faut-il tirer pour la même montée ?
\item Compare les deux dispositifs : quel est l'avantage et quel est l'inconvénient de la poulie mobile ?
\end{enumerate}
\end{exercice}

\begin{exercice}[Exploitation de mesures au plan incliné]
Un groupe d'élèves tire un chariot de poids $P = \SI{100}{N}$ sur un plan incliné et mesure l'effort $F$ au dynamomètre pour différents angles $\alpha$ du plan :
\begin{center}
\begin{tabular}{|l|c|c|c|c|c|}
\hline
\rowcolor{popblueL} Angle $\alpha$ (degrés) & 10 & 20 & 30 & 45 & 60 \\
\hline Effort $F$ (N) & 17 & 34 & 50 & 71 & 87 \\
\hline
\end{tabular}
\end{center}
\begin{enumerate}
\item Trace la courbe $F = f(\alpha)$ (échelles : \SI{1}{cm} pour 10 degrés ; \SI{1}{cm} pour \SI{10}{N}).
\item Comment évolue l'effort $F$ quand l'angle augmente ?
\item Pour quel angle l'effort serait-il égal au poids du chariot ? Justifie.
\item Déduis-en pourquoi les routes de montagne montent en lacets.
\end{enumerate}
\end{exercice}

\begin{exercice}[Sécurité sur un chantier]
Sur un chantier à Yaoundé, on observe la scène suivante : un ouvrier soulève des sacs de ciment avec une poulie dont la corde est effilochée ; un autre ouvrier se tient juste sous la charge suspendue ; la poulie est fixée à une poutre vermoulue.
\begin{enumerate}
\item Relève trois erreurs de sécurité commises sur ce chantier.
\item Pour chaque erreur, propose la consigne de sécurité correcte.
\item Rappelle pourquoi il est dangereux de se placer sous une charge suspendue, même avec une machine simple en bon état.
\end{enumerate}
\end{exercice}

\courssub{Je me prépare au Bac}

\begin{exercice}[Type BEPC : le levier et la brouette]
\textbf{Partie A.} Un maçon veut soulever une pierre de poids $P = \SI{600}{N}$ à l'aide d'une barre de longueur totale \SI{3}{m}. Il place le point d'appui à \SI{0,6}{m} de la pierre.
\begin{enumerate}
\item Définis les termes : machine simple, point d'appui, bras de levier.
\item Schématise le dispositif en indiquant la charge, l'effort et les deux bras de levier.
\item Énonce la condition d'équilibre du levier, puis calcule l'effort $F$ que le maçon doit exercer.
\item Le maçon peut-il soulever la pierre s'il est capable de fournir un effort maximal de \SI{200}{N} ? Justifie.
\end{enumerate}
\textbf{Partie B.} Le même maçon transporte du sable dans une brouette. La charge de \SI{500}{N} est placée à \SI{0,4}{m} de la roue, et ses mains sont à \SI{1,2}{m} de la roue.
\begin{enumerate}
\item Où se trouve le point d'appui dans une brouette ?
\item Calcule l'effort exercé par chacune des deux mains du maçon (on suppose l'effort également réparti).
\item Cite deux autres objets de la vie courante qui fonctionnent comme un levier.
\end{enumerate}
\end{exercice}

\begin{exercice}[Type BEPC : le treuil du forage]
Pour remonter un seau de poids $P = \SI{240}{N}$ d'un forage de \SI{6}{m} de profondeur, un village utilise un treuil dont le tambour a un rayon $r = \SI{15}{cm}$ et la manivelle une longueur $L = \SI{60}{cm}$. On néglige les frottements.
\begin{enumerate}
\item Décris la constitution d'un treuil et fais-en un schéma légendé.
\item La condition d'équilibre du treuil s'écrit $F \times L = P \times r$. Calcule l'effort $F$ nécessaire pour remonter le seau.
\item Calcule le nombre de tours de manivelle nécessaires pour remonter le seau de \SI{6}{m} (on rappelle que la circonférence du tambour vaut $2\pi r$ ; prendre $\pi \approx 3,14$).
\item Un enfant ne peut fournir qu'un effort de \SI{40}{N}. Peut-il remonter le seau ? Justifie par un calcul.
\item Cite un avantage et un inconvénient du treuil par rapport à la poulie fixe.
\end{enumerate}
\end{exercice}

\begin{exercice}[Type BEPC : problème de synthèse — le déménagement]
Pour monter un réfrigérateur de poids $P = \SI{800}{N}$ dans un camion dont le plancher est à $h = \SI{1}{m}$ du sol, deux déménageurs hésitent entre trois méthodes.
\begin{enumerate}
\item \textbf{Méthode 1 :} soulever le réfrigérateur verticalement. Quel effort total doivent-ils fournir ?
\item \textbf{Méthode 2 :} utiliser une planche de \SI{2,5}{m} comme plan incliné. On admet que l'effort est donné par $F = P \times \dfrac{h}{L}$ où $L$ est la longueur de la planche. Calcule l'effort à fournir.
\item \textbf{Méthode 3 :} utiliser une poulie mobile accrochée au-dessus du camion. Calcule l'effort nécessaire (frottements négligés) et la longueur de corde à tirer pour monter de \SI{1}{m}.
\item Classe les trois méthodes de la plus avantageuse à la moins avantageuse en termes d'effort.
\item Un proverbe dit : « ce qu'on gagne en force, on le perd en distance ». Vérifie cette affirmation en comparant les méthodes 1 et 3 (effort et longueur déplacée).
\item Propose deux consignes de sécurité que les déménageurs doivent respecter pendant cette opération.
\end{enumerate}
\end{exercice}

\newpage

\coursec{Corrigés des exercices}

\begin{corrige}[Exercice 1 — QCM : la bonne machine simple]
\begin{enumerate}
\item \textbf{a) Levier.} La brouette tourne autour de l'axe de la roue (point d'appui) ; la charge est entre la roue et les poignées.
\item \textbf{b) Plan incliné.} La rampe permet de monter la moto avec un effort réduit, mais sur une distance plus grande.
\item \textbf{a) Levier.} Le casse-noix est un levier dont le point d'appui est la charnière ; la noix est la charge.
\item \textbf{d) Treuil.} La manivelle actionne un axe (tambour) sur lequel s'enroule la courroie ou la corde.
\item \textbf{c) Poulie.} La corde passe sur une roue : c'est une poulie fixe qui change la direction de l'effort.
\item \textbf{a) Levier.} Le pied-de-biche pivote autour d'un point d'appui placé près de la pierre.
\item \textbf{b) Plan incliné.} Un escalier est un plan incliné en escalier : on monte moins brutalement qu'à la verticale.
\item \textbf{d) Treuil} (ou poulie si le store est remonté par une corde sur une roue). Le système à manivelle qui enroule la sangle est un treuil.
\end{enumerate}
\end{corrige}

\begin{corrige}[Exercice 2 — Vrai ou faux ? Justifie]
\begin{enumerate}
\item \textbf{Vrai.} C'est précisément le rôle d'une machine simple : réduire l'effort à fournir (en l'échangeant contre une plus grande distance parcourue).
\item \textbf{Faux.} C'est la \emph{poulie mobile} qui divise l'effort par deux ; la poulie fixe ne fait que changer la direction de l'effort ($F = P$).
\item \textbf{Faux.} Plus le plan incliné est long pour une même hauteur, plus la pente est douce et plus l'effort est \emph{petit} ($F = P \times \dfrac{h}{L}$).
\item \textbf{Vrai.} D'après la condition d'équilibre $F \times d_F = P \times d_P$, à charge et à $d_P$ fixés, plus $d_F$ est grand, plus $F$ est petit.
\end{enumerate}
\end{corrige}

\begin{corrige}[Exercice 3 — Définitions à compléter]
\begin{enumerate}
\item Un dispositif qui permet de déplacer ou de soulever une charge en réduisant l'effort à fournir est une \cle{machine simple}.
\item Dans un levier, la barre tourne autour d'un point fixe appelé \cle{point d'appui}.
\item La force exercée par l'utilisateur sur la machine s'appelle l'\cle{effort} ; le poids de l'objet à déplacer s'appelle la \cle{charge}.
\item Le \cle{treuil} est constitué d'un tambour sur lequel s'enroule une corde, actionné par une manivelle.
\end{enumerate}
\end{corrige}

\begin{corrige}[Exercice 4 — Calculs directs]
\begin{enumerate}
\item Avec une poulie mobile, l'effort est divisé par deux : $F = \dfrac{P}{2}$, donc $P = 2F = 2 \times 150 = \SI{300}{N}$.
\item Réduction d'effort : $180 - 60 = \SI{120}{N}$. L'effort est réduit de \SI{120}{N} par rapport à une montée verticale.
\item Pour un treuil, $F \times L = P \times r$, donc $\dfrac{F}{P} = \dfrac{r}{L} = \dfrac{10}{40} = \dfrac{1}{4}$. L'effort vaut le quart de la charge.
\end{enumerate}
\end{corrige}

\begin{corrige}[Exercice 5 — Le levier du chantier]
\begin{enumerate}
\item Schéma légendé :
\begin{popfigure}
\begin{tikzpicture}[scale=1.1]
\fill[poppurple] (1.5,0) -- (2.1,0) -- (1.8,0.45) -- cycle;
\draw[line width=2.5pt,popdark] (0,0.55) -- (5,0.85);
\fill[poporange] (0.15,0.56) rectangle (0.85,1.16);
\draw[->,line width=1.5pt,popblue] (4.6,2.1) -- (4.6,1.0);
\node[popblue] at (4.6,2.35) {$F$ (effort)};
\node[poporange!80!black] at (0.5,1.45) {charge $P$};
\node[poppurple] at (1.8,-0.25) {point d'appui};
\draw[<->,popgreen!60!black] (0.5,0.2) -- (1.8,0.2) node[midway,below] {\small $d_P = 0{,}5$ m};
\draw[<->,popgreen!60!black] (1.8,-0.55) -- (4.6,-0.55) node[midway,below] {\small $d_F = 2$ m};
\end{tikzpicture}
\legende{Levier du chantier : la barre pivote autour du point d'appui ; la charge est proche de l'appui, l'effort est appliqué loin de l'appui.}
\end{popfigure}
\item Bras de levier de l'effort : $d_F = 2{,}5 - 0{,}5 = \SI{2}{m}$.
\item Condition d'équilibre : $F \times d_F = P \times d_P$, donc
\[ F = \dfrac{P \times d_P}{d_F} = \dfrac{400 \times 0{,}5}{2} = \SI{100}{N}. \]
\item L'ouvrier n'exerce que \SI{100}{N} pour soulever \SI{400}{N} : l'effort est divisé par quatre. Le levier est d'autant plus avantageux que le point d'appui est proche de la charge.
\end{enumerate}
\end{corrige}

\begin{corrige}[Exercice 6 — Le puits du quartier]
\begin{enumerate}
\item \textbf{Poulie fixe :} elle ne fait que changer la direction de l'effort, donc $F = P = \SI{120}{N}$. Pour que le seau monte de \SI{4}{m}, il faut tirer \SI{4}{m} de corde.
\item \textbf{Poulie mobile :} la charge est supportée par deux brins de corde, donc $F = \dfrac{P}{2} = \dfrac{120}{2} = \SI{60}{N}$. Mais chaque brin doit raccourcir de \SI{4}{m} : il faut tirer $2 \times 4 = \SI{8}{m}$ de corde.
\item \textbf{Comparaison :} l'avantage de la poulie mobile est de diviser l'effort par deux (\SI{60}{N} au lieu de \SI{120}{N}) ; son inconvénient est de devoir tirer deux fois plus de corde (\SI{8}{m} au lieu de \SI{4}{m}) : ce qu'on gagne en force, on le perd en distance.
\end{enumerate}
\end{corrige}

\begin{corrige}[Exercice 7 — Exploitation de mesures au plan incliné]
\begin{enumerate}
\item Courbe $F = f(\alpha)$ :
\begin{popfigure}
\begin{tikzpicture}
\begin{axis}[width=0.85\linewidth,height=6.5cm,xlabel={Angle $\alpha$ (degrés)},ylabel={Effort $F$ (N)},xmin=0,xmax=65,ymin=0,ymax=100,xtick={0,10,20,30,40,50,60},ytick={0,10,20,30,40,50,60,70,80,90,100},grid=both,grid style={popline!40}]
\addplot[popcurve,mark=*,mark size=1.5pt] coordinates {(10,17) (20,34) (30,50) (45,71) (60,87)};
\end{axis}
\end{tikzpicture}
\legende{Effort nécessaire pour tirer le chariot en fonction de l'angle du plan incliné.}
\end{popfigure}
\item Quand l'angle $\alpha$ augmente, l'effort $F$ augmente : la courbe est croissante. La croissance est de plus en plus rapide pour les grands angles.
\item L'effort serait égal au poids ($F = P = \SI{100}{N}$) pour un angle de $90$ degrés : le plan serait alors vertical et tirer le chariot reviendrait à le soulever directement.
\item Les routes de montagne montent en lacets pour \emph{allonger le trajet} et donc \emph{réduire la pente} : à hauteur égale, un plan plus long exige un effort plus petit, ce qui permet aux véhicules de monter sans forcer excessivement.
\end{enumerate}
\end{corrige}

\begin{corrige}[Exercice 8 — Sécurité sur un chantier]
\begin{enumerate}
\item Les trois erreurs de sécurité :
\begin{itemize}
\item utiliser une corde \emph{effilochée} (usée) ;
\item se tenir \emph{sous la charge suspendue} ;
\item fixer la poulie sur une \emph{poutre vermoulue} (fragilisée).
\end{itemize}
\item Consignes correctes :
\begin{itemize}
\item vérifier l'état de la corde avant chaque utilisation et la remplacer dès qu'elle est usée ;
\item interdire à toute personne de se placer sous une charge suspendue (zone interdite balisée) ;
\item fixer la poulie sur un support solide, sain et vérifié, capable de supporter la charge.
\end{itemize}
\item Même avec du matériel en bon état, une rupture imprévue (corde, crochet, nœud) reste possible : si la charge tombe, elle écrase toute personne située en dessous. Le poids de la charge transforme une simple chute en accident grave, voire mortel.
\end{enumerate}
\end{corrige}

\begin{corrige}[Exercice 9 — Type BEPC : le levier et la brouette]
\textbf{Partie A.}
\begin{enumerate}
\item \textbf{Définitions.} Une \cle{machine simple} est un dispositif qui permet de déplacer ou de soulever une charge en réduisant l'effort à fournir. Le \cle{point d'appui} est le point fixe autour duquel la barre du levier pivote. Le \cle{bras de levier} est la distance entre le point d'appui et le point d'application d'une force (charge ou effort).
\item Schéma :
\begin{popfigure}
\begin{tikzpicture}[scale=1.05]
\fill[poppurple] (1.2,0) -- (1.8,0) -- (1.5,0.4) -- cycle;
\draw[line width=2.5pt,popdark] (0,0.5) -- (6,0.95);
\fill[poporange] (0.1,0.5) rectangle (0.8,1.1);
\node[poporange!80!black] at (0.45,1.4) {$P = 600$ N};
\draw[->,line width=1.5pt,popblue] (5.6,2.2) -- (5.6,1.05);
\node[popblue] at (5.6,2.45) {$F$};
\node[poppurple] at (1.5,-0.25) {point d'appui};
\draw[<->,popgreen!60!black] (0.45,0.15) -- (1.5,0.15) node[midway,below] {\small $d_P = 0{,}6$ m};
\draw[<->,popgreen!60!black] (1.5,-0.55) -- (5.6,-0.55) node[midway,below] {\small $d_F = 2{,}4$ m};
\end{tikzpicture}
\legende{Levier : charge près de l'appui, effort à l'autre extrémité de la barre.}
\end{popfigure}
\item \textbf{Condition d'équilibre du levier :} le produit de l'effort par son bras de levier est égal au produit de la charge par son bras de levier : $F \times d_F = P \times d_P$.\\
Ici $d_F = 3 - 0{,}6 = \SI{2,4}{m}$, donc :
\[ F = \dfrac{P \times d_P}{d_F} = \dfrac{600 \times 0{,}6}{2{,}4} = \dfrac{360}{2{,}4} = \SI{150}{N}. \]
\item Oui : l'effort nécessaire (\SI{150}{N}) est inférieur à l'effort maximal du maçon (\SI{200}{N}), donc il peut soulever la pierre.
\end{enumerate}
\textbf{Partie B.}
\begin{enumerate}
\item Dans une brouette, le point d'appui est l'\emph{axe de la roue}, autour duquel la caisse bascule.
\item Condition d'équilibre par rapport à la roue : $F_{\text{total}} \times 1{,}2 = P \times 0{,}4$, donc
\[ F_{\text{total}} = \dfrac{500 \times 0{,}4}{1{,}2} = \dfrac{200}{1{,}2} \approx \SI{166,7}{N}. \]
L'effort étant également réparti sur les deux mains : $F_{\text{main}} = \dfrac{166{,}7}{2} \approx \SI{83,3}{N}$ par main.
\item Autres leviers de la vie courante : les ciseaux, la pince à linge, le casse-noix, le pied-de-biche, la balançoire à bascule, l'ouvre-bouteille.
\end{enumerate}
\textbf{Barème indicatif (sur 10) :} définitions 1,5 pt ; schéma légendé 1,5 pt ; condition d'équilibre énoncée 1 pt ; calcul de $F$ 1,5 pt ; justification A.4 1 pt ; point d'appui de la brouette 0,5 pt ; calcul Partie B 1,5 pt ; exemples 1 pt ; unités et rédaction 0,5 pt.\\
\textbf{Points de vigilance :} citer la condition d'équilibre \emph{avant} de calculer ; convertir toutes les longueurs en mètres ; ne pas oublier l'unité (N) dans le résultat ; pour A.4, comparer explicitement \SI{150}{N} et \SI{200}{N}.
\end{corrige}

\begin{corrige}[Exercice 10 — Type BEPC : le treuil du forage]
\begin{enumerate}
\item Un treuil est constitué d'un \emph{tambour} (cylindre) sur lequel s'enroule une corde, et d'une \emph{manivelle} fixée à l'axe du tambour, qui permet de le faire tourner.
\begin{popfigure}
\begin{tikzpicture}[scale=0.95]
\draw[fill=popblueL] (0,0) circle (0.75);
\draw[fill=white] (0,0) circle (0.12);
\draw[line width=2pt,popdark] (0,0) -- (1.8,0) -- (1.8,0.6);
\draw[fill=poporange] (1.65,0.6) rectangle (1.95,0.95);
\draw[->,line width=1.3pt,popblue] (2.3,1.3) -- (1.95,0.85);
\node[popblue] at (2.6,1.5) {$F$};
\draw[popgold,line width=1.5pt] (-0.75,0) -- (-0.75,-2.2);
\draw[fill=popgreenL] (-1.15,-2.9) rectangle (-0.35,-2.2);
\node at (-0.75,-3.2) {\small seau ($P$)};
\draw[<->,popgreen!60!black] (0,1.15) -- (1.8,1.15) node[midway,above] {\small $L = 60$ cm};
\draw[<->,popgreen!60!black] (0.15,-0.15) -- (0.68,-0.62) node[midway,right] {\small $r$};
\end{tikzpicture}
\legende{Treuil : la manivelle de longueur $L$ fait tourner le tambour de rayon $r$ ; la corde s'enroule et remonte le seau.}
\end{popfigure}
\item Condition d'équilibre : $F \times L = P \times r$, donc
\[ F = \dfrac{P \times r}{L} = \dfrac{240 \times 15}{60} = \SI{60}{N}. \]
\item Circonférence du tambour : $C = 2\pi r = 2 \times 3{,}14 \times 15 = \SI{94,2}{cm} = \SI{0,942}{m}$. À chaque tour de manivelle, la corde s'enroule d'une circonférence. Nombre de tours :
\[ n = \dfrac{6}{0{,}942} \approx 6{,}4 \text{ tours}. \]
Il faut donc environ 7 tours de manivelle (6 tours complets ne suffisent pas).
\item Oui : l'effort nécessaire est $F = \SI{60}{N}$... mais l'enfant ne peut fournir que \SI{40}{N}. Or $40 < 60$, donc \textbf{non}, l'enfant ne peut pas remonter le seau seul. (Il lui faudrait une manivelle plus longue : $L = \dfrac{P \times r}{F} = \dfrac{240 \times 15}{40} = \SI{90}{cm}$.)
\item \textbf{Avantage} du treuil sur la poulie fixe : il réduit fortement l'effort (ici divisé par 4) alors que la poulie fixe ne le réduit pas. \textbf{Inconvénient} : il faut tourner la manivelle de nombreux tours (grande distance parcourue par la main) pour une montée courte, et le mécanisme est plus lourd et plus encombrant.
\end{enumerate}
\textbf{Barème indicatif (sur 10) :} description et schéma 2 pts ; calcul de $F$ 2 pts ; circonférence et nombre de tours 2 pts ; justification 4 2 pts ; avantage et inconvénient 1,5 pt ; unités et rédaction 0,5 pt.\\
\textbf{Points de vigilance :} dans $F \times L = P \times r$, $r$ et $L$ doivent être dans la \emph{même} unité (ici cm) ; pour le nombre de tours, convertir la circonférence en mètres avant de diviser ; arrondir le nombre de tours à l'entier supérieur ; pour la question 4, la comparaison chiffrée est exigée.
\end{corrige}

\begin{corrige}[Exercice 11 — Type BEPC : problème de synthèse — le déménagement]
\begin{enumerate}
\item \textbf{Méthode 1 :} pour soulever verticalement le réfrigérateur, il faut compenser entièrement son poids : $F_1 = P = \SI{800}{N}$.
\item \textbf{Méthode 2 :} plan incliné de longueur $L = \SI{2,5}{m}$ :
\[ F_2 = P \times \dfrac{h}{L} = 800 \times \dfrac{1}{2{,}5} = \SI{320}{N}. \]
\item \textbf{Méthode 3 :} poulie mobile : l'effort est divisé par deux :
\[ F_3 = \dfrac{P}{2} = \dfrac{800}{2} = \SI{400}{N}. \]
Pour monter de \SI{1}{m}, chacun des deux brins de corde doit raccourcir de \SI{1}{m} : il faut tirer $2 \times 1 = \SI{2}{m}$ de corde.
\item Classement de la plus avantageuse (plus petit effort) à la moins avantageuse :
\[ F_2 = \SI{320}{N} \;<\; F_3 = \SI{400}{N} \;<\; F_1 = \SI{800}{N}, \]
soit : plan incliné, puis poulie mobile, puis soulèvement direct.
\item Vérification du proverbe avec les méthodes 1 et 3 :
\begin{itemize}
\item Méthode 1 : effort \SI{800}{N}, déplacement de la charge sur \SI{1}{m} ;
\item Méthode 3 : effort \SI{400}{N} (deux fois plus petit), mais il faut tirer \SI{2}{m} de corde (deux fois plus long).
\end{itemize}
L'effort a été divisé par deux exactement au prix d'une distance doublée : « ce qu'on gagne en force, on le perd en distance » est bien vérifié.
\item Consignes de sécurité (deux parmi) :
\begin{itemize}
\item vérifier l'état de la corde, de la poulie et de son point de fixation avant l'opération ;
\item ne jamais se placer sous le réfrigérateur suspendu ;
\item porter des gants et des chaussures de protection, se tenir le dos droit pour pousser ;
\item fixer solidement la planche pour qu'elle ne glisse pas pendant la montée.
\end{itemize}
\end{enumerate}
\textbf{Barème indicatif (sur 10) :} méthode 1 : 1 pt ; méthode 2 : 2 pts ; méthode 3 (effort + longueur) : 2 pts ; classement : 1 pt ; vérification du proverbe : 2 pts ; consignes de sécurité : 1,5 pt ; unités et rédaction : 0,5 pt.\\
\textbf{Points de vigilance :} appliquer la formule $F = P \times \dfrac{h}{L}$ avec $h$ et $L$ dans la même unité ; ne pas oublier que la poulie mobile double la longueur de corde à tirer ; le classement doit être justifié par les valeurs calculées ; la vérification du proverbe exige une comparaison chiffrée (facteur 2), pas une simple affirmation.
\end{corrige}

\newpage

\coursec{Fiche bilan}

\begin{popfigure}
\begin{tikzpicture}[mindmap, grow cyclic, every node/.style={concept, font=\small},
  root concept/.append style={concept color=popblue, font=\bfseries\small, minimum size=2.6cm},
  level 1/.append style={level distance=3.4cm, sibling angle=60, font=\scriptsize},
  level 1 concept/.append style={minimum size=2.2cm}]
\node[root concept]{Machines simples}
  child[concept color=poporange]{ node {Le levier} }
  child[concept color=popgreen]{ node {Plan incliné} }
  child[concept color=poppurple]{ node {La poulie} }
  child[concept color=poppink]{ node {Le treuil} }
  child[concept color=popgold]{ node {Avantages et inconvénients} }
  child[concept color=popblue!70]{ node {Sécurité} };
\end{tikzpicture}
\legende{Carte mentale de la leçon : les machines simples (levier, plan incliné, poulie, treuil) permettent de soulever ou déplacer des charges avec moins d'effort, en respectant les règles de sécurité.}
\end{popfigure}

\begin{bilan}
\textbf{Définitions clés}
\begin{itemize}
  \item Une \cle{machine simple} est un dispositif qui permet d'accomplir un travail (soulever, déplacer une charge) en \textbf{réduisant l'effort} à fournir ou en \textbf{changeant la direction} de la force.
  \item Le \cle{levier} est une barre rigide qui tourne autour d'un point fixe appelé \cle{pivot} (ou point d'appui). Exemples : pied-de-biche, brouette, casse-noix, balançoire.
  \item Le \cle{plan incliné} est une surface plane inclinée par rapport à l'horizontale : il permet de monter une charge avec un effort plus faible, mais sur une distance plus longue. Exemples : rampe de chargement, route en lacets, escalier.
  \item La \cle{poulie} est une roue à gorge sur laquelle passe une corde. Une poulie \textbf{fixe} change le sens de l'effort ; une poulie \textbf{mobile} divise l'effort par deux.
  \item Le \cle{treuil} est formé d'un tambour (cylindre) sur lequel s'enroule une corde, mû par une manivelle. Exemples : treuil de puits, treuil de dépannage.
\end{itemize}

\textbf{Formules et propriétés à connaître par cœur}
\begin{center}
\begin{tabularx}{\linewidth}{|l|X|X|}
\hline
\rowcolor{popblueL} \textbf{Machine} & \textbf{Condition d'équilibre / relation} & \textbf{Gain} \\
\hline
Levier & $F \times d_{F} = R \times d_{R}$ (moments) & $F < R$ si $d_{F} > d_{R}$ \\
\hline
Plan incliné & $F < P$ ; l'effort diminue quand la pente diminue & Effort réduit, distance augmentée \\
\hline
Poulie fixe & $F = P$ & Change le sens de l'effort \\
\hline
Poulie mobile & $F = \dfrac{P}{2}$ & Effort divisé par 2 \\
\hline
Treuil & $F \times R = P \times r$ & $F < P$ si $R > r$ (manivelle longue) \\
\hline
\end{tabularx}
\end{center}
où $F$ est l'effort (force motrice), $R$ ou $P$ la résistance (poids de la charge), $d_{F}$ et $d_{R}$ les distances au pivot, $R$ le rayon de la manivelle et $r$ le rayon du tambour.

\textbf{Méthodes express}
\begin{itemize}
  \item \textbf{Identifier} la machine simple dans une situation : chercher le pivot (levier), la pente (plan incliné), la corde et la roue (poulie), le tambour et la manivelle (treuil).
  \item \textbf{Calculer un effort} avec le levier : écrire l'égalité des moments $F \times d_F = R \times d_R$, puis isoler $F$.
  \item \textbf{Bilan général} : aucune machine simple ne supprime le travail ; on gagne en \textbf{force}, on perd en \textbf{distance} (et inversement).
  \item \textbf{Sécurité} : ne jamais dépasser la charge maximale, vérifier les cordes et les fixations, garder une position stable.
\end{itemize}
\end{bilan}

\begin{aretenir}[Checklist avant le BEPC]
\begin{itemize}
  \item $\square$ Je sais définir une machine simple et citer son rôle (réduire l'effort ou changer sa direction).
  \item $\square$ Je sais définir le levier et repérer le pivot, l'effort et la résistance sur un schéma.
  \item $\square$ Je connais la condition d'équilibre du levier : $F \times d_F = R \times d_R$.
  \item $\square$ Je sais expliquer pourquoi le plan incliné réduit l'effort (mais allonge la distance).
  \item $\square$ Je distingue poulie fixe ($F = P$, changement de sens) et poulie mobile ($F = \dfrac{P}{2}$).
  \item $\square$ Je sais décrire le treuil (tambour + manivelle + corde) et citer un exemple (puits).
  \item $\square$ Je sais citer au moins deux exemples de la vie courante pour chaque machine simple.
  \item $\square$ Je sais comparer avantages (effort réduit) et inconvénients (distance plus grande, encombrement).
  \item $\square$ Je sais résoudre un calcul simple d'effort avec le levier ou la poulie mobile.
  \item $\square$ Je connais les consignes de sécurité : charge maximale, état des cordes, position stable.
  \item $\square$ Je sais que le travail ne disparaît jamais : ce qu'on gagne en force, on le perd en distance.
\end{itemize}
\end{aretenir}

\findecours

\end{document}
